% ThinkAlgo — Số học và Tư duy số — Lớp 7
% Edition: student (no solutions), VI.  Generated by build_grade7.py — edit body/ and rebuild.
\documentclass[11pt]{article}
\newif\ifsol \solfalse
\usepackage{iftex}
\ifPDFTeX
  \usepackage[utf8]{inputenc}\usepackage[T5]{fontenc}
\else
  \usepackage{fontspec}
\fi
\usepackage[a4paper,top=2.95cm,headheight=1.5cm,headsep=0.45cm,bottom=3.1cm,left=1.9cm,right=1.9cm]{geometry}
\usepackage{amsmath,amssymb}
\usepackage{tikz}\usetikzlibrary{calc,shadows,shapes.geometric}
\usepackage[most]{tcolorbox}
\usepackage{enumitem}
\usepackage{array}
\usepackage{tabularx}
\usepackage{colortbl}
\usepackage{comment}
\usepackage{needspace}
\usepackage{fancyhdr}
\usepackage{fontawesome5}
\usepackage{graphicx}
\usepackage{microtype}
% PDF outline (bookmarks) so each chapter shows in the reader's outline — does not change the page layout
\usepackage[hidelinks,unicode,bookmarksnumbered=false,bookmarksopen=true,bookmarksopenlevel=0]{hyperref}
\usepackage{bookmark}
\usetikzlibrary{decorations.pathreplacing,arrows.meta,angles,quotes,patterns,positioning,through,intersections}

% ---------------- palette: navy family + subtle light tint; orange ONLY in the logo ----------------
% ----- EXACT palette from the user's original main (36).tex (vivid, not pale) -----
\definecolor{tmBlue}{RGB}{66,133,244}   % problem / header / section
\definecolor{tmGreen}{RGB}{52,168,83}   % solution
\definecolor{tmOrange}{RGB}{255,152,0}  % example
\colorlet{Navy}{tmBlue}                 % header rule / section heading / badge / divider
\definecolor{NavyD}{HTML}{222222}       % solution body text (dark, readable)
\colorlet{Gold}{Navy}
\colorlet{ProbBlue}{tmBlue}             % problem boxes
\colorlet{SolGreen}{tmGreen}            % solution boxes
\definecolor{AnsRed}{HTML}{E02424}      % vivid red for answers / totals
\definecolor{CardGreen}{HTML}{34A853}   % digit-card green
\definecolor{StarGold}{HTML}{F2B134}    % difficulty star (gold)
\definecolor{KidOrange}{HTML}{DE8A2C}   % logo accent only
\definecolor{Ink}{HTML}{1F2733}       % body text
\definecolor{Slate}{HTML}{5B6A82}     % muted labels
\definecolor{Line}{HTML}{C6D0DF}      % light hairlines / student boxes
\definecolor{Tint}{HTML}{EDF1F7}      % pale panel fill (worked solution / answer)
% aliases so older names keep working
\colorlet{TextDark}{Ink}\colorlet{TextGray}{Slate}\colorlet{Soft}{Slate}
\colorlet{CardLine}{Line}\colorlet{Teal}{Navy}\colorlet{ProbB}{Navy}\colorlet{SolG}{Navy}
% bright, kid-friendly colours for bar-model (sơ đồ đoạn thẳng) blocks
\definecolor{tmPurple}{RGB}{124,77,226}
\definecolor{tmTeal}{RGB}{0,150,136}
\definecolor{tmPink}{RGB}{233,30,99}
\AtBeginDocument{\color{Ink}}
\raggedbottom  % natural flow, even spacing (no stretched gaps)
\setlength{\parindent}{0pt}\setlength{\parskip}{0.18em}
\renewcommand{\arraystretch}{1.35}
\providecommand{\AnsWord}{Đáp số}

% ---------------- brand logo (original mark) ----------------
\newcommand{\logomark}[1]{\begin{tikzpicture}[baseline=-0.55ex,scale=#1]
  \draw[KidOrange,line width=1.5pt] (-0.27,0) .. controls (-0.27,0.23) and (0,0.23) .. (0,0)
        .. controls (0,-0.23) and (-0.27,-0.23) .. (-0.27,0);
  \draw[Navy,line width=1.5pt] (0.27,0) .. controls (0.27,0.23) and (0,0.23) .. (0,0)
        .. controls (0,-0.23) and (0.27,-0.23) .. (0.27,0);
\end{tikzpicture}}
% real brand logo image (gold infinity + "think algo" wordmark)
\newcommand{\logohead}{\IfFileExists{logo.png}{\includegraphics[height=1.35cm]{logo.png}}%
  {\raisebox{0.35cm}{\logomark{1.5}\hspace{0.5em}{\Large\bfseries\color{KidOrange}think\,\color{Navy}algo}}}}

% ---------------- running header (logo + Name/Date + topic) on EVERY page, single navy rule ----------------
\pagestyle{fancy}\fancyhf{}
\renewcommand{\headrulewidth}{1.4pt}\renewcommand{\footrulewidth}{0pt}
\renewcommand{\headrule}{{\color{Navy}\hrule height 1.3pt}}
\fancyhead[L]{\raisebox{-0.95cm}{\logohead}}
\fancyhead[R]{\setlength{\tabcolsep}{0pt}%
  \begin{tabular}{@{}r@{}}
    {\small\color{Ink}{\color{tmBlue}\faUserEdit}~\NameWord:~\rule{3.0cm}{0.5pt}\quad{\color{tmBlue}\faCalendarDay}~\DateWord:~\rule{1.9cm}{0.5pt}}\\[5pt]
    {\large\bfseries\color{Navy}\CurTopic}
  \end{tabular}}

% ---------------- footer (rounded page border + Topic/Grade + thinkalgo.com + page + tagline) ----------------
% NOTE: a remember-picture overlay must be placed DIRECTLY in \fancyfoot, NOT inside \def/\newcommand.
\fancyfoot[C]{\begin{tikzpicture}[remember picture, overlay]
  \draw[Navy!40, line width=1.4pt, rounded corners=12pt]
     ($(current page.north west)+(0.8cm,-0.8cm)$) rectangle ($(current page.south east)+(-0.8cm,0.45cm)$);
  \coordinate (FL) at ($(current page.south west)+(1.2cm, 2.18cm)$);
  \coordinate (FR) at ($(current page.south east)+(-1.2cm, 2.18cm)$);
  \draw[Navy!55, line width=1pt] (FL) -- (FR);
  \node[fill=white, draw=Navy, line width=1.2pt, rounded corners=8pt, font=\bfseries\small,
        text=Navy, inner xsep=11pt, inner ysep=4pt] at ($(current page.south)+(0,2.18cm)$) {{\color{tmBlue}\faGlobe}\ \ thinkalgo.com};
  \node[text=Slate, font=\footnotesize, anchor=north west, align=left, inner sep=0pt, text width=7.6cm]
        at ($(FL)+(0.1cm,-0.30cm)$) {{\color{tmBlue}\faTag}\ \TopicWord: \textbf{\color{Ink}\CurTopic}\\[1.5pt]{\color{tmBlue}\faGraduationCap}\ \DocGrade};
  \node[text=Slate, font=\footnotesize] at ($(current page.south)+(0,1.20cm)$) {\thepage};
  \node[text=Slate, font=\scriptsize, anchor=north east, align=right, inner sep=0pt, text width=7.4cm]
        at ($(FR)+(-0.1cm,-0.34cm)$) {{\normalsize\bfseries\color{Navy}ThinkAlgo}\ — From Math Thinking to Code};
\end{tikzpicture}}

% ---------------- page-1 lead-in (topic is already in the header; keep a slim spacer) ----------------
\newcommand{\titlehead}{\vspace*{-0.2cm}}
\providecommand{\CurTopic}{}

% chapter banner + sets the running-header topic. #1 = number, #2 = title
\newcommand{\chapterhead}[2]{\renewcommand{\CurTopic}{#2}%
  \setcounter{seccnt}{0}% each chapter is independent: section badges restart at 1
  \phantomsection\pdfbookmark[0]{#1. #2}{chap#1}\label{chap:#1}% add this chapter to the PDF outline + TOC ref
  \par\vspace{0.1cm}\noindent\begin{tikzpicture}[baseline=-0.5ex]
    \node[circle,fill=tmBlue,text=white,font=\Large\bfseries,minimum size=0.95cm,inner sep=0pt](cn){#1};
    \node[anchor=west,font=\LARGE\bfseries,text=tmBlue] at ([xshift=12pt]cn.east){#2};
  \end{tikzpicture}\par\vspace{4pt}\noindent\textcolor{tmBlue}{\rule{\linewidth}{2pt}}\par\vspace{0.5cm}}

% ---------------- section heading (outline navy number badge, no fill) ----------------
\newcounter{seccnt}
\newcommand{\shead}[1]{\par\vspace{0.7em}\refstepcounter{seccnt}%
  \noindent\begin{tikzpicture}[baseline=-0.6ex]
    \node[circle,draw=Gold,line width=1.2pt,text=Gold,font=\bfseries,minimum size=1.05em,inner sep=1pt]{\theseccnt};
  \end{tikzpicture}\hspace{0.4em}{\large\bfseries\color{Navy} #1}\par\vspace{0.25em}}

% ============ boxes EXACTLY per the original main (36).tex: solid-blue PROBLEM, solid-orange EXAMPLE,
%              outline-green SOLUTION, all with drop shadow + rounded corners ============
% ---------------- knowledge box: white, tmBlue border + left accent + drop shadow ----------------
\newtcolorbox{knowbox}{enhanced,breakable,colback=white,colframe=tmBlue!70,boxrule=0.9pt,arc=5pt,
  left=12pt,right=11pt,top=6pt,bottom=6pt,drop shadow,
  borderline west={3.5pt}{0pt}{tmBlue}}

% ---------------- EXAMPLE (Ví dụ): SOLID ORANGE tab, white text ----------------
\newtcolorbox{exbox}[1]{enhanced,breakable,colback=white,colframe=tmOrange,boxrule=1pt,arc=5pt,
  left=12pt,right=11pt,top=14pt,bottom=7pt,drop shadow,
  attach boxed title to top left={xshift=10pt,yshift=-\tcboxedtitleheight/2},
  boxed title style={colback=tmOrange,colframe=tmOrange,arc=4pt,boxrule=0pt,drop shadow},
  title={#1},fonttitle=\bfseries\color{white}}
\newcommand{\guide}{\par\vspace{2pt}{\bfseries\color{tmBlue}\GuideWord}\par\nopagebreak}

% ---------------- PROBLEM (Bài): SOLID BLUE tab, white text ----------------
\newtcolorbox{probbox}[2]{enhanced,breakable,colback=white,colframe=tmBlue,boxrule=1pt,arc=5pt,
  left=12pt,right=10pt,top=14pt,bottom=7pt,drop shadow,
  attach boxed title to top left={xshift=10pt,yshift=-\tcboxedtitleheight/2},
  boxed title style={colback=tmBlue,colframe=tmBlue,arc=4pt,boxrule=0pt,drop shadow},
  fonttitle=\bfseries\color{white},title={\ProbWord\ #1\ifx\relax#2\relax\else#2\fi}}

% ---------------- SOLUTION (Lời giải): OUTLINE GREEN tab (white fill, green border + text) ----------------
% Teacher: a green tcolorbox. Student: the body is dropped entirely (comment pkg — robust with tables/\\).
\ifsol
  \newtcolorbox{solbox}{enhanced,breakable,colback=white,colframe=tmGreen,boxrule=1pt,arc=5pt,
    left=12pt,right=10pt,top=14pt,bottom=7pt,drop shadow,
    attach boxed title to top left={xshift=10pt,yshift=-\tcboxedtitleheight/2},
    boxed title style={colback=white,colframe=tmGreen,boxrule=1pt,arc=4pt,drop shadow},
    fonttitle=\bfseries\color{tmGreen!88!black},title={\SolTitle},before upper={\color{NavyD}}}
\else
  \excludecomment{solbox}
\fi

% ---------------- answer highlight: label green, value RED ----------------
\newcommand{\ans}[1]{\par\smallskip{\bfseries\color{tmGreen!88!black}\AnsWord:}\ {\bfseries\color{AnsRed}#1}\par}

% ---------------- digit cards (exact original style: fill=c!10, border=c!60!black, drop shadow) ----------------
% \digitcards{0/gray,1/cyan,2/green,3/cyan} : value/colour pairs
\newcommand{\digitcards}[1]{\begin{center}\begin{tikzpicture}[scale=0.7,transform shape,font=\sffamily]
  \def\dx{0}\foreach \v/\c in {#1}{
    \node[draw=\c!60!black,line width=1.5pt,fill=\c!10,rounded corners=4pt,
          minimum width=1cm,minimum height=1.4cm,drop shadow] at (\dx,0){\Huge\bfseries \v};
    \pgfmathsetmacro\dx{\dx+1.5}\xdef\dx{\dx}}
\end{tikzpicture}\end{center}}

% ---------------- student working space (white, light border) ----------------
\newcommand{\worksp}[1]{\ifsol\else\begin{tcolorbox}[colback=white,colframe=Line,boxrule=0.8pt,
  arc=5pt,height=\dimexpr#1+0.05cm\relax,left=8pt,right=8pt,top=4pt,bottom=4pt]{\footnotesize\color{Slate}\itshape\WorkWord}\end{tcolorbox}\fi}

% ---------------- \probsol{solution} : Teacher shows the green solution box; Student shows a BLANK
%   "Bài làm" box whose height matches that solution's height (writing room = work needed), so both
%   versions have identical per-problem height and paginate the same. ----------------
\newsavebox{\solmeasure}\newlength{\solht}
\newcommand{\probsol}[1]{%
  \begin{lrbox}{\solmeasure}\begin{minipage}{\dimexpr\linewidth-24pt\relax}#1\end{minipage}\end{lrbox}%
  \ifsol
    \begin{solbox}#1\end{solbox}%
  \else
    \setlength{\solht}{\dimexpr\ht\solmeasure+\dp\solmeasure+0.85cm\relax}%
    \ifdim\solht<1.35cm\setlength{\solht}{1.35cm}\fi
    \ifdim\solht>\StudentMaxWork\setlength{\solht}{\StudentMaxWork}\fi
    \begin{tcolorbox}[colback=white,colframe=Line,boxrule=0.8pt,arc=5pt,
      left=8pt,right=8pt,top=4pt,bottom=4pt,height=\solht]{\footnotesize\color{Slate}\itshape\WorkWord}\end{tcolorbox}%
  \fi}

% ---------------- practice divider: blue left accent bar, no fill ----------------
\newcommand{\dividerbar}[1]{\par\vspace{0.6em}\noindent
  \begin{tcolorbox}[enhanced,frame hidden,interior hidden,borderline west={5pt}{0pt}{ProbBlue},
    colback=white,left=12pt,right=12pt,top=5pt,bottom=5pt,arc=4pt]
  {\large\bfseries\color{ProbBlue} #1}\end{tcolorbox}\vspace{0.4em}}

% place-value choice-count diagram (Cambridge style): #1 = "label/count/note/color" list, #2 = result
% notes must NOT contain commas or slashes (they separate \foreach fields)
\newcommand{\countboxes}[2]{\begin{center}\begin{tikzpicture}[scale=0.62,transform shape,font=\sffamily]
  \def\xx{0}%
  \foreach \lab/\cnt/\note/\col [count=\i] in {#1}{%
    \ifnum\i>1\relax \node[text=gray,font=\Huge] at (\xx-1.8,0){$\times$};\fi
    \node[draw=\col!65,fill=\col!8,rounded corners=3pt,line width=1pt,minimum width=2.7cm,minimum height=1.7cm] at (\xx,0){};
    \node[text=\col!85,font=\bfseries] at (\xx,0.58){\lab};
    \node[text=\col!88,font=\Huge\bfseries] at (\xx,-0.1){\cnt};
    \node[text=gray,font=\scriptsize] at (\xx,-1.05){\note};
    \pgfmathsetmacro\xx{\xx+3.6}\xdef\xx{\xx}}%
  \node[text=gray,font=\Huge] at (\xx-1.8,0){$=$};
  \node[text=red,font=\Huge\bfseries] at (\xx,0){#2};
\end{tikzpicture}\end{center}}

% ======== bar model (sơ đồ đoạn thẳng): lightly-rounded blocks, darker border, tiny gaps, bright colours ========
% one block: \blk{x1}{x2}{yc}{height}{colour}{label}  (a small gap is left on each side)
\newcommand{\blk}[6]{%
  \draw[rounded corners=1.3pt,fill=#5!16,draw=#5!80!black,line width=1pt]
    (#1+0.045,#3-#4/2) rectangle (#2-0.045,#3+#4/2);
  \node[font=\small\bfseries,text=#5!52!black] at ({(#1+#2)/2},#3){#6};}
% measurement with a double-headed ARROW (not a brace) + label above the span: \brcup{x1}{x2}{y}{label}
\newcommand{\brcup}[4]{%
  \draw[<->,>=stealth,tmBlue!72!black,line width=0.7pt] (#1+0.05,#3)--(#2-0.05,#3);
  \node[above=-1.5pt,font=\footnotesize,text=tmBlue!58!black] at ({(#1+#2)/2},#3){#4};}
% measurement with a double-headed ARROW + label below the span: \brcdn{x1}{x2}{y}{label}
\newcommand{\brcdn}[4]{%
  \draw[<->,>=stealth,tmBlue!72!black,line width=0.7pt] (#1+0.05,#3)--(#2-0.05,#3);
  \node[below=-1.5pt,font=\footnotesize,text=tmBlue!58!black] at ({(#1+#2)/2},#3){#4};}
% row label at the left of a bar: \rowlab{x}{y}{text}
\newcommand{\rowlab}[3]{\node[anchor=east,font=\small\bfseries,text=Ink] at (#1,#2){#3};}
% vertical right-pointing CURLY BRACE spanning two stacked bars: \vbrace{x}{ytop}{ybot}{label}
\newcommand{\vbrace}[4]{%
  \pgfmathsetmacro\ym{(#2+#3)/2}%
  \draw[tmBlue!72!black,line width=0.9pt,line cap=round]
    (#1,#2) .. controls (#1+0.18,#2) and (#1+0.18,#2-0.18) .. (#1+0.18,#2-0.22)
    -- (#1+0.18,{\ym+0.22})
    .. controls (#1+0.18,\ym) and (#1+0.30,{\ym+0.08}) .. (#1+0.42,\ym)
    .. controls (#1+0.30,{\ym-0.08}) and (#1+0.18,\ym) .. (#1+0.18,{\ym-0.22})
    -- (#1+0.18,#3+0.22)
    .. controls (#1+0.18,#3+0.18) and (#1+0.18,#3) .. (#1,#3);
  \node[right=4pt,font=\footnotesize,text=tmBlue!58!black] at (#1+0.42,\ym){#4};}
% blank 3x3 grid for a problem statement (student fills it in)
\newcommand{\blankgrid}{\begin{tikzpicture}[font=\sffamily]
  \foreach \r in {0,1,2}{\foreach \c in {0,1,2}{
    \draw[rounded corners=1.3pt,fill=Tint,draw=Slate!55,line width=0.9pt]
      (\c+0.05,\r+0.05) rectangle (\c+0.95,\r+0.95);}}
\end{tikzpicture}}
% wrapper: LEFT-aligned (not centred), a little air above/below
\newenvironment{barfig}{\par\vspace{3pt}\noindent\kern0.35cm\begin{tikzpicture}[font=\sffamily,>=stealth]}%
  {\end{tikzpicture}\par\vspace{3pt}}
% wrapper: CENTRED (for illustrative diagrams in theory / examples)
\newenvironment{barfigC}{\par\vspace{3pt}\begin{center}\begin{tikzpicture}[font=\sffamily,>=stealth]}%
  {\end{tikzpicture}\end{center}\vspace{2pt}}

% ---------------- difficulty stars (nicely drawn) : \stars{1}=hard, \stars{2}=very hard ----------------
\newcommand{\diffstar}{\tikz[baseline=-0.55ex]{\node[star,star points=5,star point ratio=0.46,
  draw=StarGold!72!black,fill=StarGold,line width=0.3pt,inner sep=1.75pt]{};}}
\newcommand{\stars}[1]{\ifcase#1\or\kern2pt\diffstar\or\kern2pt\diffstar\kern1.5pt\diffstar\fi}

\newcommand{\nd}[1]{\overline{#1}}

% number line: #1 last index, #2 ticks "i/label", #3 open (excluded) indices, #4 filled (solution) indices
\newcommand{\numlinex}[4]{\begin{center}\begin{tikzpicture}[x=1.55cm,>=stealth]
  \draw[->,thick,Navy](-0.45,0)--(#1+0.45,0);
  \foreach \i/\lab in {#2}{\draw[Navy,line width=0.8pt](\i,0.08)--(\i,-0.08);
     \node[below,font=\small,text=Ink]at(\i,-0.1){\lab};}
  \foreach \i in {#3}{\draw[Navy,fill=white,line width=1pt](\i,0)circle(2.3pt);}
  \foreach \i in {#4}{\fill[KidOrange](\i,0)circle(2.7pt);}
\end{tikzpicture}\end{center}}


% ======================================================================
% Grade 7 additions (shared by all four editions)
% ======================================================================
% student edition: blank "work" box height = min(solution height, this cap)
\newlength{\StudentMaxWork}\setlength{\StudentMaxWork}{7.5cm}

% ThinkAlgo tip box (label set per language by \TipWord)
\newtcolorbox{tipbox}{enhanced,breakable,colback=StarGold!8,colframe=StarGold!75!black,boxrule=1pt,arc=4pt,
  left=11pt,right=10pt,top=6pt,bottom=6pt,
  before upper={\textbf{\color{StarGold!65!black}\faLightbulb\ \TipWord}\ }}

% "Recall" box: re-states a piece of theory inside a solution / example
\newtcolorbox{recallbox}{enhanced,breakable,colback=tmBlue!5,colframe=tmBlue!45,boxrule=0.6pt,arc=3pt,
  left=8pt,right=8pt,top=4pt,bottom=4pt,
  before upper={\textbf{\color{tmBlue!80!black}\faBookOpen\ \RecallWord}\ }}

% "Why is it true? / Proof" box
\newtcolorbox{proofbox}{enhanced,breakable,colback=tmPurple!4,colframe=tmPurple!45,boxrule=0.6pt,arc=3pt,
  left=8pt,right=8pt,top=4pt,bottom=4pt,
  before upper={\textbf{\color{tmPurple!85!black}\faSearch\ \ProofWord}\ },
  after upper={\hfill$\square$}}

% "Common mistake" box
\newtcolorbox{warnbox}{enhanced,breakable,colback=AnsRed!4,colframe=AnsRed!45,boxrule=0.6pt,arc=3pt,
  left=8pt,right=8pt,top=4pt,bottom=4pt,
  before upper={\textbf{\color{AnsRed!85!black}\faExclamationTriangle\ \WarnWord}\ }}

% "Note for parents" box (short hint on how to explain the idea at home)
\newtcolorbox{parentbox}{enhanced,breakable,colback=tmTeal!5,colframe=tmTeal!50,boxrule=0.6pt,arc=3pt,
  left=8pt,right=8pt,top=4pt,bottom=4pt,
  before upper={\textbf{\color{tmTeal!80!black}\faUsers\ \ParentWord}\ }}

% numbered step inside a solution:  \step{1}{Short title}  -> "Bước 1. Short title"
\newcommand{\step}[2]{\par\smallskip\noindent{\bfseries\color{tmBlue!85!black}\StepWord~#1.}\ {\bfseries #2}\par\nopagebreak}
% verification line:  \verify{...}
\newcommand{\verify}[1]{\par\smallskip\noindent{\bfseries\color{tmGreen!80!black}\faCheckCircle\ \CheckWord}\ #1\par}

% multiple-choice options
% \choices{..}{..}{..}{..}{..}      : five short options on one line
% \choicesII{..}{..}{..}{..}{..}    : two/three per line (medium length)
% \choicesL{..}{..}{..}{..}{..}     : one option per line (long text)
\newcommand{\choices}[5]{\par\smallskip\noindent
  \begin{tabularx}{\linewidth}{@{}XXXXX@{}}(A)\ #1&(B)\ #2&(C)\ #3&(D)\ #4&(E)\ #5\end{tabularx}\par}
\newcommand{\choicesII}[5]{\par\smallskip\noindent
  \begin{tabularx}{\linewidth}{@{}XXX@{}}(A)\ #1&(B)\ #2&(C)\ #3\\(D)\ #4&(E)\ #5&\end{tabularx}\par}
\newcommand{\choicesL}[5]{\par\smallskip\noindent
  \begin{tabular}{@{}l@{}}(A)\ #1\\(B)\ #2\\(C)\ #3\\(D)\ #4\\(E)\ #5\end{tabular}\par}

% centred figure inside a problem / solution
\newenvironment{figC}{\par\vspace{2pt}\begin{center}\begin{tikzpicture}[font=\sffamily\small,>=stealth,line join=round]}%
  {\end{tikzpicture}\end{center}\vspace{1pt}}

\def\DocGrade{Lớp 7}
\def\ProbWord{Bài}
\def\SolWord{Lời giải.}
\def\GuideWord{Hướng dẫn giải.}
\def\PageWord{Trang}
\def\NameWord{Họ và tên}
\def\DateWord{Ngày}
\def\SolTitle{Lời giải chi tiết}
\def\WorkWord{Bài làm:}
\def\TopicWord{Chủ đề}
\def\AnsWord{Đáp số}
\def\TipWord{Mẹo ThinkAlgo.}
\def\RecallWord{Nhắc lại kiến thức.}
\def\ProofWord{Vì sao lại đúng?}
\def\WarnWord{Lỗi thường gặp.}
\def\ParentWord{Gợi ý cho phụ huynh.}
\def\StepWord{Bước}
\def\CheckWord{Thử lại:}

\def\DocTopic{Số học và Tư duy số}
\begin{document}
\titlehead
\chapterhead{4}{Số học và Tư duy số}

% ---------------- part a ----------------
% =====================================================================
% Chủ đề 4 — Số học và Tư duy số — phần a: lý thuyết, ví dụ, bài 1–15
% =====================================================================
\begin{knowbox}
\textbf{Ý tưởng lớn.} ``Tư duy số'' (cảm nhận về số) là khả năng \emph{hiểu} con số chứ không chỉ \emph{tính} con số:
biết số nào lớn hơn, một kết quả có hợp lý không, nên làm tròn thế nào để ước lượng nhanh, và nhìn ra cấu trúc ẩn
bên trong một số (các chữ số, các ước, các bội). Trong chủ đề này, em sẽ ôn lại và nối kết các mảng kiến thức:
\begin{itemize}[leftmargin=1.4em,itemsep=1pt,topsep=2pt]
  \item \textbf{số thập phân và số nguyên âm}: so sánh, sắp xếp trên trục số;
  \item \textbf{thứ tự thực hiện phép tính} và \textbf{ước lượng} kết quả;
  \item \textbf{phân số, tỉ số, tỉ số phần trăm}: ``một phần của một số'', phân số bằng nhau, giảm giá;
  \item \textbf{tốc độ} -- quãng đường -- thời gian và đổi đơn vị (km/h, m/s);
  \item \textbf{chia hết, ước chung, chữ số tận cùng, cấu tạo số} (giải các phép tính đố chữ);
  \item \textbf{trung bình cộng, trung vị} và các \textbf{dãy số có quy luật} (Fibonacci, cấp số cộng).
\end{itemize}
Bí quyết chung: \emph{trước khi tính, hãy ước lượng; sau khi tính, hãy thử lại}. Một học sinh có tư duy số tốt
thường loại được ngay vài phương án vô lý trong câu trắc nghiệm mà chưa cần tính chính xác.
\end{knowbox}
\vspace{0.3cm}

% ---------------------------------------------------------------------
\shead{Số thập phân: đọc, so sánh, sắp xếp, cộng trừ}
\begin{knowbox}
\textbf{1. Giá trị theo hàng.} Trong số thập phân, mỗi chữ số đứng ở một \emph{hàng}. Sau dấu phẩy lần lượt là
hàng phần mười, hàng phần trăm, hàng phần nghìn, \ldots
\begin{center}
\renewcommand{\arraystretch}{1.25}
\begin{tabular}{|c|c|c|c|c|c|}
\hline
\rowcolor{tmBlue!10} Chục & Đơn vị & , & Phần mười & Phần trăm & Phần nghìn\\ \hline
 & $2$ & , & $2$ & $3$ & \\ \hline
 & $2$ & , & $3$ & $0$ & \\ \hline
\end{tabular}
\end{center}
Ví dụ: $2{,}23 = 2 + \dfrac{2}{10} + \dfrac{3}{100}$, còn $2{,}3 = 2 + \dfrac{3}{10} = 2{,}30$.

\textbf{2. Thêm hoặc bớt chữ số $0$ ở tận cùng phần thập phân không làm thay đổi giá trị:}
$2{,}3 = 2{,}30 = 2{,}300$. Lý do: $\dfrac{3}{10} = \dfrac{30}{100} = \dfrac{300}{1000}$ (nhân cả tử và mẫu với $10$).

\textbf{3. Quy tắc so sánh hai số thập phân dương.}
\begin{itemize}[leftmargin=1.4em,itemsep=1pt,topsep=2pt]
  \item So sánh \emph{phần nguyên} trước: số nào có phần nguyên lớn hơn thì lớn hơn.
  \item Nếu phần nguyên bằng nhau, \emph{viết thêm chữ số $0$} để hai phần thập phân có cùng số chữ số, rồi so sánh
  từng hàng từ trái sang phải: hàng phần mười, rồi hàng phần trăm, \ldots
\end{itemize}
Ví dụ: so sánh $0{,}304$ và $0{,}34$. Viết $0{,}34 = 0{,}340$. Hàng phần mười: $3 = 3$; hàng phần trăm: $0 < 4$.
Vậy $0{,}304 < 0{,}340 = 0{,}34$.

\textbf{4. Cộng, trừ số thập phân:} đặt tính sao cho \emph{các dấu phẩy thẳng cột} (hàng thẳng hàng), rồi cộng/trừ
như số tự nhiên.

\textbf{5. Làm tròn số thập phân.} Muốn làm tròn đến hàng phần trăm (hai chữ số sau dấu phẩy), ta nhìn chữ số ngay
bên phải hàng đó (hàng phần nghìn): nếu chữ số này nhỏ hơn $5$ thì giữ nguyên, nếu lớn hơn hoặc bằng $5$ thì cộng thêm
$1$ vào hàng phần trăm; sau đó bỏ các chữ số phía sau. Ví dụ: $4{,}6875 \approx 4{,}69$ (vì chữ số hàng phần nghìn là
$7 \ge 5$), còn $3{,}1416 \approx 3{,}14$ (vì $1 < 5$). Kí hiệu ``$\approx$'' đọc là ``xấp xỉ''.
\end{knowbox}
\begin{proofbox}
Vì sao không được so sánh ``theo độ dài'' (số nào nhiều chữ số hơn thì lớn hơn)? Hãy viết theo phần nghìn:
$0{,}304 = \dfrac{304}{1000}$ và $0{,}34 = \dfrac{340}{1000}$. Hai phân số cùng mẫu, tử $304 < 340$ nên
$0{,}304 < 0{,}34$, dù $0{,}304$ ``trông dài hơn''. Việc thêm chữ số $0$ chính là đưa về cùng mẫu số.
\end{proofbox}
\begin{warnbox}
Nhiều bạn nghĩ $2{,}23 > 2{,}3$ vì ``$23 > 3$''. Sai! Phải so sánh từng hàng: $2{,}23$ có $2$ phần mười,
còn $2{,}3$ có $3$ phần mười, nên $2{,}23 < 2{,}3$.
\end{warnbox}
\begin{tipbox}
Muốn tìm số nhỏ nhất (lớn nhất) trong một tập hợp số thập phân, hãy viết tất cả các số thành \emph{một cột}, dấu phẩy
thẳng nhau, thêm chữ số $0$ cho đủ dài. Khi đó em chỉ việc so sánh như các số tự nhiên.
\end{tipbox}

% ---------------------------------------------------------------------
\shead{Số nguyên âm và trục số}
\begin{knowbox}
\textbf{1. Trục số.} Các số nguyên được xếp trên một đường thẳng: số $0$ ở giữa, các số dương ở bên phải, các số âm
ở bên trái.
\numlinex{8}{0/$-4$,1/$-3$,2/$-2$,3/$-1$,4/$0$,5/$1$,6/$2$,7/$3$,8/$4$}{}{4}
\textbf{2. Quy tắc so sánh.} Trên trục số (nằm ngang), \emph{số bên trái nhỏ hơn số bên phải}. Do đó:
\begin{itemize}[leftmargin=1.4em,itemsep=1pt,topsep=2pt]
  \item mọi số âm đều nhỏ hơn $0$, và $0$ nhỏ hơn mọi số dương;
  \item mọi số âm nhỏ hơn mọi số dương;
  \item với hai số âm, số nào \emph{xa số $0$ hơn} thì \emph{nhỏ hơn}: $-8 < -2$ vì $8 > 2$.
\end{itemize}
\textbf{3. Hình ảnh đời thường.} Nhiệt độ $-8^\circ$C lạnh hơn $-2^\circ$C; nợ $8$ nghìn đồng thì ``nghèo hơn'' nợ
$2$ nghìn đồng.

\textbf{4. Mặt phẳng tọa độ.} Một điểm được xác định bởi cặp số $(x;y)$: $x$ là hoành độ (đi sang phải/trái),
$y$ là tung độ (đi lên/xuống). Di chuyển ngang chỉ làm thay đổi $x$, di chuyển dọc chỉ làm thay đổi $y$.
\end{knowbox}
\begin{proofbox}
Vì sao $-8 < -2$? Ta có $-2 = -8 + 6$: muốn đi từ $-8$ đến $-2$ trên trục số ta phải đi \emph{sang phải} $6$ đơn vị.
Đi sang phải nghĩa là tăng lên, nên $-2$ lớn hơn $-8$.
\end{proofbox}
\begin{warnbox}
Đừng so sánh số âm ``theo cảm giác của số dương''. $-8$ \emph{nhỏ hơn} $-2$, dù $8$ lớn hơn $2$. Hãy luôn vẽ
nhanh một trục số khi phân vân.
\end{warnbox}
\begin{tipbox}
\textbf{Tính chẵn -- lẻ} là một công cụ ``tư duy số'' rất mạnh: chẵn $+$ chẵn $=$ chẵn, lẻ $+$ lẻ $=$ chẵn, chẵn $+$ lẻ $=$ lẻ;
tích của các số lẻ luôn là số lẻ. Chẳng hạn, tổng của $5$ số lẻ là số lẻ, tổng của $4$ số lẻ là số chẵn. Khi một điểm
di chuyển trên lưới tọa độ, xét tính chẵn -- lẻ của tọa độ thường cho biết ngay một vị trí có thể đến được hay không.
\end{tipbox}

% ---------------------------------------------------------------------
\shead{Thứ tự thực hiện phép tính và ước lượng}
\begin{knowbox}
\textbf{1. Thứ tự thực hiện phép tính.}
\begin{enumerate}[leftmargin=1.6em,itemsep=1pt,topsep=2pt]
  \item Trong \emph{ngoặc} trước: $(\;)$, rồi $[\;]$, rồi $\{\;\}$.
  \item Tiếp theo là \emph{lũy thừa}.
  \item Rồi đến \emph{nhân và chia}, làm lần lượt từ trái sang phải.
  \item Cuối cùng là \emph{cộng và trừ}, làm lần lượt từ trái sang phải.
\end{enumerate}
Ví dụ: $6\times 2 - 3 = 12 - 3 = 9$ (nhân trước, trừ sau), còn $6\times(2-3) = 6\times(-1) = -6$.

\textbf{2. Làm tròn để ước lượng.} Để biết nhanh một kết quả ``khoảng bao nhiêu'', ta làm tròn các thừa số đến số
tròn chục, tròn trăm gần nhất rồi tính nhẩm. Ví dụ: $99 \times 9 \approx 100 \times 9 = 900$.

\textbf{3. Mẹo tính nhanh:} $99 \times 9 = (100-1)\times 9 = 900 - 9 = 891$ (tính chất phân phối của phép nhân
đối với phép trừ).
\end{knowbox}
\begin{proofbox}
Vì sao $(100-1)\times 9 = 100\times 9 - 1\times 9$? Hãy tưởng tượng $9$ hàng, mỗi hàng $100$ chiếc kẹo,
nhưng mỗi hàng bị lấy đi $1$ chiếc. Số kẹo còn lại là $9$ hàng $\times$ $99$ chiếc, cũng bằng
($9\times 100$ chiếc) trừ đi ($9$ chiếc bị lấy). Đó là \emph{tính chất phân phối}: $a\times(b-c) = a\times b - a\times c$.
\end{proofbox}
\begin{warnbox}
Tính ``từ trái sang phải bất kể phép tính gì'' là sai: $2 + 3\times 4 \ne 5\times 4$. Đúng là
$2 + 3\times 4 = 2 + 12 = 14$.
\end{warnbox}
\begin{tipbox}
Trong câu trắc nghiệm hỏi ``số nào gần nhất với\ldots'', các phương án thường cách nhau gấp $10$ lần
($10$; $100$; $1000$; \ldots). Khi đó chỉ cần ước lượng thô là đủ chọn đúng.
\end{tipbox}

% ---------------------------------------------------------------------
\shead{Phân số: một phần của một số, phân số bằng nhau}
\begin{knowbox}
\textbf{1. Phân số là một phần của cái toàn thể.} $\dfrac{a}{b}$ nghĩa là chia toàn thể thành $b$ phần bằng nhau
và lấy $a$ phần.

\textbf{2. Tìm giá trị phân số của một số:} muốn tìm $\dfrac{m}{n}$ của $a$, ta tính $a \times \dfrac{m}{n}$
(chia $a$ thành $n$ phần rồi lấy $m$ phần). Ví dụ: $\dfrac{1}{3}$ của $396$ là $396 : 3 = 132$.

\textbf{3. Phần còn lại.} Nếu đã làm được $\dfrac{m}{n}$ công việc thì phần còn lại là $1 - \dfrac{m}{n} = \dfrac{n-m}{n}$.

\textbf{4. Phân số bằng nhau (tính chất cơ bản).} Nhân (hoặc chia) cả tử và mẫu với cùng một số khác $0$ thì được
phân số bằng nó:
\[ \frac{3}{7} = \frac{3\times 9}{7\times 9} = \frac{27}{63}, \qquad \frac{12}{36} = \frac{12:12}{36:12} = \frac{1}{3}. \]
\textbf{5. Rút gọn phân số:} chia cả tử và mẫu cho ước chung lớn nhất của chúng để được phân số tối giản.

\textbf{6. Bài toán ``tổng và tỉ''.} Nếu hai số có tỉ số $\dfrac{5}{8}$ thì có thể coi số thứ nhất gồm $5$ phần,
số thứ hai gồm $8$ phần bằng nhau. Biết tổng (hoặc hiệu) thì tìm được giá trị một phần.

\textbf{7. Phân số của một phân số.} ``Một nửa của một phần tư'' là $\dfrac12\times\dfrac14 = \dfrac18$: nhân tử với tử,
mẫu với mẫu. Điều này hay gặp khi một hình được chia đôi nhiều lần liên tiếp: mỗi lần chia đôi, phần diện tích nhân thêm
với $\dfrac12$.

\textbf{8. So sánh phân số.} Quy đồng mẫu số rồi so sánh tử số, hoặc đổi ra số thập phân. Ví dụ:
$\dfrac23 = \dfrac{8}{12} > \dfrac{7}{12}$; \ $\dfrac35 = 0{,}6 < \dfrac34 = 0{,}75$.
\end{knowbox}
\begin{barfigC}
  \rowlab{-0.2}{0}{Số thứ nhất}
  \foreach \i in {0,...,4}{\blk{\i*0.9}{\i*0.9+0.9}{0}{0.55}{tmBlue}{}}
  \rowlab{-0.2}{-0.8}{Số thứ hai}
  \foreach \i in {0,...,7}{\blk{\i*0.9}{\i*0.9+0.9}{-0.8}{0.55}{tmOrange}{}}
  \vbrace{7.35}{0.3}{-1.1}{tổng $= 13$ phần}
\end{barfigC}
\begin{proofbox}
Vì sao nhân cả tử và mẫu với cùng một số thì phân số không đổi? $\dfrac{3}{7}$ là ``chia cái bánh thành $7$ phần,
lấy $3$ phần''. Nếu cắt \emph{mỗi} phần thành $9$ miếng nhỏ hơn, cái bánh có $7\times 9 = 63$ miếng, còn phần ta lấy
có $3\times 9 = 27$ miếng. Lượng bánh không đổi, nên $\dfrac{3}{7} = \dfrac{27}{63}$.
\end{proofbox}
\begin{warnbox}
Đọc kỹ câu hỏi: đề hỏi phần \emph{đã làm} hay phần \emph{còn lại}? Nhiều bạn tính đúng $\dfrac13$ của $396$ là $132$
rồi khoanh luôn, trong khi đề hỏi số trang \emph{còn lại} là $396 - 132 = 264$.
\end{warnbox}
\begin{parentbox}
Với phân số, hình vẽ luôn thắng lời giảng. Hãy vẽ một đoạn thẳng (thanh sô-cô-la) chia đều thành các phần và tô
màu phần đã dùng; con sẽ tự thấy phần còn lại. Đây cũng là ``sơ đồ đoạn thẳng'' quen thuộc ở tiểu học.
\end{parentbox}

% ---------------------------------------------------------------------
\shead{Tỉ số phần trăm, đơn giá và giảm giá}
\begin{knowbox}
\textbf{1. Phần trăm nghĩa là ``trên một trăm''.} $80\% = \dfrac{80}{100} = 0{,}8$.

\textbf{2. Tìm $p\%$ của một số $a$:} tính $a \times \dfrac{p}{100}$. Ví dụ: $80\%$ của $50$ là
$50 \times \dfrac{80}{100} = 40$.

\textbf{3. Tìm tỉ số phần trăm của $a$ so với $b$:} tính $\dfrac{a}{b}\times 100\%$. Ví dụ: đúng $32$ câu trên $40$ câu
là $\dfrac{32}{40}\times 100\% = 80\%$.

\textbf{4. Đơn giá.} Đơn giá là giá của \emph{một} đơn vị (một ki-lô-gam, một quyển, \ldots):
\[ \text{Thành tiền} = \text{đơn giá} \times \text{số lượng}, \qquad \text{đơn giá} = \text{thành tiền} : \text{số lượng}. \]
\textbf{5. Giảm giá $p\%$.} Giá mới $= $ giá cũ $\times \left(1 - \dfrac{p}{100}\right)$. Giảm $10\%$ tức là còn
$90\%$: nhân với $0{,}9$.

\textbf{6. Giảm giá nhiều lần liên tiếp.} Mỗi lần giảm $10\%$ \emph{của giá hiện tại} thì ta nhân thêm một lần với
$0{,}9$. Ví dụ, giá $100$ nghìn đồng sau hai lần giảm: $100 \to 90 \to 81$ (nghìn đồng), \emph{không phải} $80$.
\end{knowbox}
\begin{proofbox}
Vì sao hai lần giảm $10\%$ không bằng một lần giảm $20\%$? Lần thứ hai giảm $10\%$ của $90$ (là $9$), chứ không phải
$10\%$ của $100$ (là $10$). Vì giá đã nhỏ đi nên số tiền được giảm lần sau cũng nhỏ đi: tổng cộng chỉ giảm
$10 + 9 = 19$ nghìn đồng, tức $19\%$.
\end{proofbox}
\begin{tipbox}
Tìm $10\%$ thật nhanh: chỉ cần chia cho $10$ (dời dấu phẩy sang trái một chữ số). Từ $10\%$ suy ra
$5\%$ (chia đôi), $20\%$ (gấp đôi), $80\% = 100\% - 20\%$, \ldots
\end{tipbox}

% ---------------------------------------------------------------------
\shead{Tốc độ, quãng đường, thời gian và đổi đơn vị}
\begin{knowbox}
\textbf{1. Ba công thức gắn với nhau.} Với chuyển động đều:
\[ s = v \times t, \qquad v = s : t, \qquad t = s : v, \]
trong đó $s$ là quãng đường, $v$ là vận tốc (tốc độ), $t$ là thời gian.

\textbf{2. Ai nhanh hơn?} Người nhanh hơn là người đi được quãng đường \emph{dài hơn} trong \emph{cùng một khoảng thời
gian}. Cách chắc chắn nhất: tính vận tốc của từng người theo cùng một đơn vị (ví dụ m/s) rồi so sánh.

\textbf{3. Chạy nhiều vòng.} Chạy $n$ vòng quanh đường dài $L$ mét thì quãng đường là $n \times L$ mét.

\textbf{4. Đổi km/h sang m/s.} $1$ km $= 1000$ m, $1$ giờ $= 3600$ giây, nên
\[ 1 \text{ km/h} = \frac{1000 \text{ m}}{3600 \text{ s}} = \frac{5}{18} \text{ m/s}, \qquad
   18 \text{ km/h} = 18\times\frac{5}{18} = 5 \text{ m/s}, \qquad 1 \text{ m/s} = 3{,}6 \text{ km/h}. \]
\end{knowbox}
\begin{figC}
  \draw[tmBlue,line width=1.2pt,rounded corners=8pt] (0,0) rectangle (3.6,1.6);
  \node[draw=tmOrange,fill=tmOrange!15,rounded corners=3pt,minimum width=1.6cm] at (1.8,0.8) {$s = v\times t$};
  \node[text=Ink,align=left,anchor=west] at (4.2,0.8) {Che đại lượng cần tìm:\\ $v = s : t$ \quad và \quad $t = s : v$};
\end{figC}
\begin{warnbox}
Hai đại lượng phải \emph{cùng hệ đơn vị}: nếu quãng đường tính bằng mét và thời gian bằng giây thì vận tốc ra m/s;
đừng chia mét cho giờ. Trước khi so sánh hai vận tốc, hãy đưa về cùng đơn vị.
\end{warnbox}

% ---------------------------------------------------------------------
\shead{Chia hết, ước chung, chữ số tận cùng và cấu tạo số}
\begin{knowbox}
\textbf{1. Bội và ước.} $a$ là \emph{bội} của $b$ (hay $b$ là \emph{ước} của $a$) nếu $a = b\times k$ với $k$ là số
tự nhiên. Ví dụ: $45 = 9\times 5$ nên $45$ là bội của $9$.

\textbf{2. Dấu hiệu chia hết.}
\begin{itemize}[leftmargin=1.4em,itemsep=1pt,topsep=2pt]
  \item cho $2$: chữ số tận cùng là $0, 2, 4, 6, 8$; \quad cho $5$: tận cùng là $0$ hoặc $5$; \quad cho $10$: tận cùng là $0$;
  \item cho $3$: tổng các chữ số chia hết cho $3$; \quad cho $9$: \emph{tổng các chữ số chia hết cho $9$}.
\end{itemize}
\textbf{3. Ước chung.} Số chia hết cả $a$ lẫn $b$ gọi là ước chung của $a$ và $b$. Mọi ước chung đều là ước của
\emph{ước chung lớn nhất} ƯCLN$(a,b)$. Ví dụ: ƯCLN$(40,72)=8$, nên các ước chung của $40$ và $72$ là
$1, 2, 4, 8$.

\textbf{4. Chữ số tận cùng của một tích} chỉ phụ thuộc vào chữ số tận cùng của các thừa số:
$\ldots 3 \times \ldots 7$ tận cùng bằng chữ số tận cùng của $3\times 7 = 21$, tức là $1$.

\textbf{5. Cấu tạo số.} Một số viết bằng các chữ số được phân tích theo hàng:
\[ \nd{abc} = 100a + 10b + c, \qquad \nd{PQQ} = 100P + 10Q + Q = 100P + 11Q. \]
(Dấu gạch ngang trên đầu cho biết đây là số có các chữ số $a, b, c$, không phải tích $a\times b\times c$.)
\end{knowbox}
\begin{proofbox}
Vì sao dấu hiệu chia hết cho $9$ đúng? Xét số $\nd{abc} = 100a + 10b + c$. Ta tách:
\[ 100a + 10b + c = (99a + 9b) + (a + b + c) = 9\times(11a + b) + (a+b+c). \]
Phần $9\times(11a+b)$ luôn chia hết cho $9$. Vậy $\nd{abc}$ chia hết cho $9$ khi và chỉ khi phần còn lại
$a+b+c$ (tổng các chữ số) chia hết cho $9$. Với số có nhiều chữ số hơn, lập luận y hệt vì
$10 = 9+1$, $100 = 99+1$, $1000 = 999+1$, \ldots
\end{proofbox}
\begin{tipbox}
Gặp bài ``đố chữ'' (mỗi chữ cái là một chữ số) như $\nd{PQQ}+\nd{PPQ}+\nd{QQQ}=876$, hãy làm \emph{từ hàng đơn vị
sang trái}, nhớ để ý các số nhớ. Hàng đơn vị thường có ít khả năng nhất, nên tìm được chữ số đầu tiên rất nhanh.
Cũng có thể dùng cấu tạo số để đưa về một phương trình. Khi phải đếm hoặc tìm hết các khả năng (các bộ chữ số,
các kết quả khi tung đồng xu, \ldots), hãy \emph{liệt kê có hệ thống} (theo thứ tự tăng dần, hoặc bằng bảng/sơ đồ
cây) để không sót, không trùng.
\end{tipbox}
\begin{knowbox}
\textbf{Liệt kê có hệ thống và khả năng xảy ra.}
\begin{itemize}[leftmargin=1.4em,itemsep=1pt,topsep=2pt]
  \item Muốn đếm các số có ba chữ số mà \emph{tích} các chữ số bằng một số cho trước, ta tìm các bộ ba chữ số (viết theo
  thứ tự tăng dần để không trùng), rồi đếm số cách sắp xếp mỗi bộ: ba chữ số khác nhau có $3\times 2\times 1 = 6$ cách xếp;
  hai chữ số giống nhau có $3$ cách xếp. Ví dụ: tích bằng $12$ với bộ $\{1;3;4\}$ cho $6$ số $134, 143, 314, 341, 413, 431$;
  bộ $\{2;2;3\}$ cho $3$ số $223, 232, 322$.
  \item Khi mọi kết quả đều \emph{có khả năng như nhau}, xác suất của một sự kiện bằng
  \[ \frac{\text{số kết quả thuận lợi}}{\text{tổng số kết quả có thể}}. \]
  Ví dụ: tung $2$ đồng xu có $4$ kết quả đồng khả năng: SS, SN, NS, NN (S: sấp, N: ngửa). Xác suất cả hai đồng xu cùng
  một mặt là $\dfrac{2}{4} = \dfrac12$. Dùng sơ đồ cây để liệt kê:
\end{itemize}
\begin{figC}
  \node[draw=Navy,circle,inner sep=1.5pt] (r) at (0,0) {\scriptsize bắt đầu};
  \node[draw=tmBlue,fill=tmBlue!10,circle,inner sep=2pt] (a) at (2,0.8) {S};
  \node[draw=tmOrange,fill=tmOrange!10,circle,inner sep=2pt] (b) at (2,-0.8) {N};
  \foreach \n/\y/\l/\c in {aa/1.2/S/tmBlue,ab/0.4/N/tmOrange}{\node[draw=\c,fill=\c!10,circle,inner sep=2pt] (\n) at (4,\y) {\l};}
  \foreach \n/\y/\l/\c in {ba/-0.4/S/tmBlue,bb/-1.2/N/tmOrange}{\node[draw=\c,fill=\c!10,circle,inner sep=2pt] (\n) at (4,\y) {\l};}
  \draw[Slate] (r)--(a) (r)--(b) (a)--(aa) (a)--(ab) (b)--(ba) (b)--(bb);
  \node[right,font=\scriptsize] at (4.3,1.2) {SS \ (cùng mặt)};
  \node[right,font=\scriptsize] at (4.3,0.4) {SN};
  \node[right,font=\scriptsize] at (4.3,-0.4) {NS};
  \node[right,font=\scriptsize] at (4.3,-1.2) {NN \ (cùng mặt)};
\end{figC}
\end{knowbox}
\begin{warnbox}
``Tổng các chữ số chia hết cho $9$'' là dấu hiệu của $9$, \emph{không} dùng cho $4$, $6$ hay $8$. Và một số chia hết
cho $3$ chưa chắc chia hết cho $9$: $12$ chia hết cho $3$ nhưng không chia hết cho $9$.
\end{warnbox}

% ---------------------------------------------------------------------
\shead{Dãy số có quy luật, trung bình cộng và trung vị}
\begin{knowbox}
\textbf{1. Dãy Fibonacci.} $1, 1, 2, 3, 5, 8, 13, 21, 34, 55, \ldots$ Mỗi số (kể từ số thứ ba) bằng \emph{tổng hai
số đứng ngay trước nó}: $2 = 1+1$, $3 = 1+2$, $5 = 2+3$, \ldots

\textbf{2. Cấp số cộng (dãy cách đều).} Mỗi số (kể từ số thứ hai) bằng số liền trước cộng với một số không đổi $d$
(công sai). Ví dụ: $1, 4, 7, 10$ (với $d = 3$). Số thứ $n$ bằng $a_1 + (n-1)\times d$.

\textbf{3. Trung bình cộng.} Trung bình cộng của $n$ số $=$ (tổng các số) $: n$. Suy ra
\[ \text{Tổng} = \text{trung bình cộng} \times \text{số các số}. \]
\textbf{4. Phương pháp ``độ lệch''.} Tổng các độ lệch so với trung bình cộng luôn bằng $0$: phần ``thừa'' của các số
lớn hơn trung bình đúng bằng phần ``thiếu'' của các số nhỏ hơn trung bình.

\textbf{5. Trung vị.} Sắp xếp các số theo thứ tự tăng dần; số đứng chính giữa là trung vị (nếu có một số lẻ các số).
Ví dụ: trung vị của $42, 43, 46, 49, 50$ là $46$. Nếu có một số chẵn các số thì trung vị là trung bình cộng của
hai số đứng giữa: trung vị của $6, 7, 9, 10$ là $(7+9):2 = 8$. Lưu ý: khi thêm hoặc thay đổi một số, vị trí ``ở giữa''
có thể thay đổi, nên phải sắp xếp lại.

\textbf{6. Các số nguyên liên tiếp.} Với một số lẻ các số nguyên liên tiếp, trung bình cộng chính là số đứng giữa:
trung bình của $19, 20, 21, 22, 23$ là $21$.
\end{knowbox}
\begin{proofbox}
Vì sao tổng các độ lệch bằng $0$? Giả sử $4$ số có trung bình $80$, tức tổng là $4\times 80 = 320$. Viết mỗi số thành
$80 + $ (độ lệch). Cộng lại: $320 + $ (tổng các độ lệch) $= 320$, nên tổng các độ lệch phải bằng $0$.
\end{proofbox}
\begin{figC}
  \draw[Slate!60] (-0.3,0)--(5.3,0);
  \node[left,font=\scriptsize,text=Slate] at (-0.3,0) {TB $=40$};
  \foreach \x/\v/\c in {0/38/tmOrange,1.5/43/tmGreen,3/41/tmGreen,4.5/38/tmOrange}{
    \pgfmathsetmacro\h{(\v-40)*0.3}
    \fill[\c!70] (\x,0) rectangle (\x+0.7,\h);
    \node[font=\scriptsize] at (\x+0.35,{\h+0.25*sign(\h+0.001)}) {$\v$};}
\end{figC}
\begin{center}\small Bốn số $38, 43, 41, 38$ có trung bình cộng $40$. Độ lệch so với $40$: $-2,\ +3,\ +1,\ -2$; tổng $= 0$.\end{center}
\begin{parentbox}
Có thể minh họa trung bình cộng bằng các chồng xu: ``san bằng'' các chồng cao xuống chồng thấp cho đến khi các chồng
cao bằng nhau. Chiều cao chung đó là trung bình cộng. Số xu chuyển đi từ chồng cao đúng bằng số xu thêm vào chồng thấp
-- đó chính là ý ``tổng độ lệch bằng $0$''.
\end{parentbox}

% =====================================================================
% VÍ DỤ MẪU
% =====================================================================
\par\vspace{0.7em}\noindent{\large\bfseries\color{tmOrange!90!black}\faPencil*\ \ Ví dụ mẫu}\par\vspace{0.3em}

\begin{exbox}{Ví dụ 1}
Cho tập hợp $\{0{,}5;\ 0{,}45;\ 0{,}405;\ 0{,}054;\ 0{,}54\}$.
\begin{enumerate}[label=\alph*),leftmargin=1.6em,itemsep=0pt,topsep=2pt]
  \item Sắp xếp các số theo thứ tự tăng dần.
  \item Tính tổng của số lớn nhất và số nhỏ nhất.
\end{enumerate}
\guide
\step{1}{Viết các số thành một cột, thêm chữ số $0$ cho đủ ba chữ số sau dấu phẩy}
Thêm chữ số $0$ vào tận cùng phần thập phân không làm thay đổi giá trị, nên ta có:
\[ 0{,}5 = 0{,}500;\quad 0{,}45 = 0{,}450;\quad 0{,}405;\quad 0{,}054;\quad 0{,}54 = 0{,}540. \]
\step{2}{So sánh như các số tự nhiên}
Bỏ ``$0{,}$'' đi, ta so sánh các số $500;\ 450;\ 405;\ 54;\ 540$:
\[ 54 < 405 < 450 < 500 < 540. \]
Vậy thứ tự tăng dần là: $0{,}054 < 0{,}405 < 0{,}45 < 0{,}5 < 0{,}54$.
\step{3}{Cộng số lớn nhất với số nhỏ nhất (đặt dấu phẩy thẳng cột)}
\[ \begin{array}{r} 0{,}540\\ +\ 0{,}054\\ \hline 0{,}594 \end{array} \]
\verify{$0{,}594 - 0{,}54 = 0{,}054$. Đúng.}
\ans{a) $0{,}054 < 0{,}405 < 0{,}45 < 0{,}5 < 0{,}54$;\quad b) $0{,}594$}
\end{exbox}

\begin{exbox}{Ví dụ 2}
Chị Hoa đạp xe trên quãng đường dài $45$ km. Buổi sáng chị đi được $\dfrac{2}{5}$ quãng đường. Hỏi buổi chiều chị còn
phải đi bao nhiêu ki-lô-mét nữa, và quãng đường đó bằng mấy phần của cả quãng đường?
\guide
\step{1}{Vẽ sơ đồ đoạn thẳng}
Chia quãng đường thành $5$ phần bằng nhau; buổi sáng đi $2$ phần, buổi chiều còn $3$ phần.
\begin{barfigC}
  \foreach \i in {0,1}{\blk{\i*1.6}{\i*1.6+1.6}{0}{0.6}{tmBlue}{$9$ km}}
  \foreach \i in {2,3,4}{\blk{\i*1.6}{\i*1.6+1.6}{0}{0.6}{tmOrange}{$9$ km}}
  \brcup{0}{3.2}{0.5}{buổi sáng}
  \brcup{3.2}{8}{0.5}{buổi chiều}
  \brcdn{0}{8}{-0.5}{$45$ km}
\end{barfigC}
\step{2}{Tìm giá trị một phần}
Một phần dài $45 : 5 = 9$ (km).
\step{3}{Tính quãng đường buổi sáng và buổi chiều}
Buổi sáng: $2\times 9 = 18$ (km). Buổi chiều: $45 - 18 = 27$ (km), cũng chính là $3\times 9 = 27$ (km).
\step{4}{Phần quãng đường còn lại}
$1 - \dfrac{2}{5} = \dfrac{5}{5} - \dfrac{2}{5} = \dfrac{3}{5}$. Kiểm tra: $\dfrac{27}{45} = \dfrac{27:9}{45:9} = \dfrac{3}{5}$.
\ans{$27$ km, bằng $\dfrac{3}{5}$ quãng đường}
\end{exbox}

\begin{exbox}{Ví dụ 3}
a) Tính $20 - 4\times 3 + 8 : 2$ \quad và \quad $(20-4)\times(3+8) : 2$.\\
b) Không tính chính xác, hãy cho biết $101 \times 49$ gần nhất với số nào: $50$; $500$; $5000$; $50\,000$? Sau đó
tính chính xác bằng cách nhanh.
\guide
\step{1}{Biểu thức không có ngoặc: nhân, chia trước; cộng, trừ sau}
\[ 20 - \underbrace{4\times 3}_{12} + \underbrace{8 : 2}_{4} = 20 - 12 + 4 = 8 + 4 = 12. \]
(Phép trừ và phép cộng làm lần lượt từ trái sang phải: $20 - 12 = 8$, rồi $8 + 4 = 12$.)
\step{2}{Biểu thức có ngoặc: làm trong ngoặc trước}
\[ (20 - 4)\times(3+8) : 2 = 16 \times 11 : 2 = 176 : 2 = 88. \]
\step{3}{Ước lượng}
$101 \approx 100$ và $49 \approx 50$, nên $101 \times 49 \approx 100\times 50 = 5000$. Chọn $5000$.
\step{4}{Tính chính xác bằng tính chất phân phối}
$101 \times 49 = (100 + 1)\times 49 = 4900 + 49 = 4949$, quả thật rất gần $5000$.
\ans{a) $12$ và $88$;\quad b) gần $5000$ nhất; $101\times 49 = 4949$}
\end{exbox}

\begin{exbox}{Ví dụ 4}
a) Gạo giá $18$ nghìn đồng một ki-lô-gam. Mẹ mua $15$ kg và được giảm $10\%$ trên tổng hóa đơn. Mẹ phải trả bao nhiêu tiền?\\
b) Một chiếc áo giá $400$ nghìn đồng. Cứ sau mỗi tuần, cửa hàng lại giảm $10\%$ \emph{giá hiện tại}. Sau ít nhất bao
nhiêu tuần thì giá áo thấp hơn $300$ nghìn đồng?
\guide
\step{1}{Thành tiền khi chưa giảm}
Thành tiền $=$ đơn giá $\times$ số lượng $= 18 \times 15 = 270$ (nghìn đồng).
\step{2}{Trừ đi $10\%$}
$10\%$ của $270$ là $270 : 10 = 27$ (nghìn đồng). Mẹ trả $270 - 27 = 243$ (nghìn đồng) $= 243\,000$ đồng.
(Hoặc: còn $90\%$, tức $270 \times 0{,}9 = 243$.)
\step{3}{Giảm giá liên tiếp: mỗi tuần nhân với $0{,}9$}
\begin{center}
\begin{tabular}{|c|c|c|c|c|}
\hline
\rowcolor{tmBlue!10} Thời điểm & Ban đầu & Sau 1 tuần & Sau 2 tuần & Sau 3 tuần\\ \hline
Giá (nghìn đồng) & $400$ & $360$ & $324$ & $291{,}6$\\ \hline
Được giảm & -- & $40$ & $36$ & $32{,}4$\\ \hline
\end{tabular}
\end{center}
Sau $2$ tuần giá là $324 > 300$; sau $3$ tuần giá là $291{,}6 < 300$.
\begin{warnbox}
Không được nghĩ ``mỗi tuần giảm $40$ nghìn'' (vì $10\%$ của $400$). Mỗi tuần giảm $10\%$ của \emph{giá lúc đó}, nên số
tiền được giảm ngày càng ít: $40$, rồi $36$, rồi $32{,}4$, \ldots
\end{warnbox}
\ans{a) $243\,000$ đồng;\quad b) $3$ tuần}
\end{exbox}

\begin{exbox}{Ví dụ 5}
a) Bình đạp xe $3$ vòng quanh một sân vận động có chu vi $400$ m hết $4$ phút. An chạy xe đạp với vận tốc $15$ km/h.
Ai đi nhanh hơn?\\
b) Một ô tô chạy với vận tốc $72$ km/h. Trong $3$ giây, ô tô đi được bao nhiêu mét?
\guide
\step{1}{Vận tốc của Bình (m/s)}
Quãng đường: $3 \times 400 = 1200$ (m). Thời gian: $4$ phút $= 4\times 60 = 240$ (giây).
\[ v_{\text{Bình}} = 1200 : 240 = 5 \text{ (m/s)}. \]
\step{2}{Đổi vận tốc của An sang m/s}
\begin{recallbox}$1$ km/h $= \dfrac{1000 \text{ m}}{3600 \text{ s}} = \dfrac{5}{18}$ m/s.\end{recallbox}
$v_{\text{An}} = 15 \times \dfrac{5}{18} = \dfrac{75}{18} \approx 4{,}17$ (m/s). Vì $5 > 4{,}17$ nên \textbf{Bình nhanh hơn}.
(Cách khác: $5$ m/s $= 5\times 3{,}6 = 18$ km/h $> 15$ km/h.)
\step{3}{Ô tô}
$72$ km/h $= 72 \times \dfrac{5}{18} = 20$ m/s. Trong $3$ giây đi được $20 \times 3 = 60$ (m).
\ans{a) Bình nhanh hơn ($5$ m/s so với khoảng $4{,}17$ m/s);\quad b) $60$ m}
\end{exbox}

\begin{exbox}{Ví dụ 6}
a) Trong các số $72;\ 84;\ 117;\ 135$, số nào chia hết cho $9$?\\
b) Tìm chữ số $x$ để số $\nd{3x5}$ chia hết cho $9$.\\
c) Có bao nhiêu số tự nhiên là ước của cả $24$ và $36$?
\guide
\step{1}{Dùng dấu hiệu chia hết cho $9$: tính tổng các chữ số}
$72$: $7+2 = 9$ (chia hết cho $9$). \quad $84$: $8 + 4 = 12$ (không chia hết cho $9$).\\
$117$: $1+1+7 = 9$ (chia hết). \quad $135$: $1+3+5 = 9$ (chia hết).\\
Vậy $72$, $117$, $135$ chia hết cho $9$; $84$ thì không ($84 = 9\times 9 + 3$).
\step{2}{Tìm chữ số $x$}
Tổng các chữ số của $\nd{3x5}$ là $3 + x + 5 = 8 + x$. Vì $x$ là chữ số nên $8 \le 8 + x \le 17$. Trong khoảng này chỉ có
$9$ chia hết cho $9$, nên $8 + x = 9$, tức $x = 1$. Số cần tìm là $315 = 9 \times 35$.
\step{3}{Ước chung của $24$ và $36$}
Liệt kê ước của $24$: $1, 2, 3, 4, 6, 8, 12, 24$. Ước của $36$: $1, 2, 3, 4, 6, 9, 12, 18, 36$.\\
Các số có mặt ở cả hai danh sách: $1, 2, 3, 4, 6, 12$ -- đó cũng chính là các ước của ƯCLN$(24,36) = 12$.
\ans{a) $72;\ 117;\ 135$;\quad b) $x = 1$;\quad c) $6$ số}
\end{exbox}

\begin{exbox}{Ví dụ 7}
Trong phép tính $\nd{AB} + \nd{AB} + \nd{AB} = \nd{1AB}$, mỗi chữ cái đại diện cho một chữ số. Tìm $A$ và $B$.
\guide
\textbf{Cách 1 -- Dùng cấu tạo số.} Ta có $\nd{1AB} = 100 + \nd{AB}$. Vậy
\[ 3\times\nd{AB} = 100 + \nd{AB} \;\Rightarrow\; 2\times\nd{AB} = 100 \;\Rightarrow\; \nd{AB} = 50. \]
Vậy $A = 5$, $B = 0$.

\textbf{Cách 2 -- Xét từng hàng.} Hàng đơn vị: $B + B + B = 3B$ có chữ số tận cùng là $B$, tức $3B - B = 2B$ tận cùng
bằng $0$, nên $B = 0$ hoặc $B = 5$.
\begin{itemize}[leftmargin=1.4em,itemsep=1pt,topsep=1pt]
  \item Nếu $B = 5$: $3\times 5 = 15$, viết $5$ nhớ $1$. Hàng chục: $3A + 1 = \nd{1A} = 10 + A$, suy ra $2A = 9$ -- không có
  chữ số nào thỏa mãn. Loại.
  \item Nếu $B = 0$: không có nhớ. Hàng chục: $3A = 10 + A$, suy ra $2A = 10$, $A = 5$.
\end{itemize}
\verify{$50 + 50 + 50 = 150 = \nd{1AB}$ với $A = 5$, $B = 0$. Đúng.}
\ans{$A = 5$, $B = 0$}
\end{exbox}

\begin{exbox}{Ví dụ 8}
a) Bốn bài kiểm tra của Nam được $6$; $9$; $10$ và $x$ điểm. Điểm trung bình của bốn bài là $8$. Tìm $x$.\\
b) Trung bình cộng của năm số tự nhiên liên tiếp là $30$. Tìm số nhỏ nhất.
\guide
\step{1}{Cách tính tổng}
\begin{recallbox}Tổng $=$ trung bình cộng $\times$ số các số.\end{recallbox}
Tổng điểm bốn bài là $8 \times 4 = 32$. Vậy $x = 32 - 6 - 9 - 10 = 7$.
\step{2}{Cách dùng độ lệch (nhanh hơn)}
So với $8$: bài $6$ thiếu $2$, bài $9$ thừa $1$, bài $10$ thừa $2$. Tổng độ lệch: $-2 + 1 + 2 = +1$. Để tổng các độ lệch
bằng $0$, bài thứ tư phải thiếu $1$: $x = 8 - 1 = 7$.
\step{3}{Năm số liên tiếp}
Với năm số liên tiếp, trung bình cộng là số đứng giữa. Số giữa là $30$, nên năm số là $28, 29, 30, 31, 32$. Số nhỏ nhất là $28$.
\verify{$28 + 29 + 30 + 31 + 32 = 150 = 5\times 30$. Đúng.}
\ans{a) $x = 7$;\quad b) $28$}
\end{exbox}

% =====================================================================
\clearpage
\dividerbar{Luyện tập cơ bản}

% ---------------------------------------------------------------- Bài 1
\begin{probbox}{1}{\relax}
Số nhỏ nhất trong tập hợp $\{3{,}2;\ 2{,}3;\ 3;\ 2{,}23;\ 3{,}22\}$ là
\choices{$3{,}2$}{$2{,}3$}{$3$}{$2{,}23$}{$3{,}22$}
\end{probbox}
\probsol{
\textbf{Phân tích đề.} Có $5$ số thập phân. Ta cần tìm số \emph{nhỏ nhất}.
\begin{recallbox}So sánh số thập phân: so sánh phần nguyên trước; nếu phần nguyên bằng nhau thì so sánh lần lượt hàng
phần mười, hàng phần trăm, \ldots\ Có thể thêm chữ số $0$ vào cuối phần thập phân mà giá trị không đổi.\end{recallbox}
\step{1}{So sánh phần nguyên}
Phần nguyên của các số: $3{,}2 \to 3$; \ $2{,}3 \to 2$; \ $3 \to 3$; \ $2{,}23 \to 2$; \ $3{,}22 \to 3$.\\
Số nhỏ nhất phải có phần nguyên nhỏ nhất, tức là bằng $2$. Chỉ còn hai ứng viên: $2{,}3$ và $2{,}23$.
\step{2}{Viết cùng số chữ số thập phân và so sánh}
$2{,}3 = 2{,}30$ \ và \ $2{,}23$. Hàng phần mười: $3$ so với $2$; vì $2 < 3$ nên $2{,}23 < 2{,}30$.
\step{3}{Kết luận}
Sắp xếp toàn bộ cho chắc chắn: $2{,}23 < 2{,}3 < 3 < 3{,}2 < 3{,}22$. Số nhỏ nhất là $2{,}23$.
\begin{warnbox}Đừng chọn $2{,}3$ chỉ vì ``$3$ nhỏ hơn $23$''. Chữ số $3$ trong $2{,}3$ là $3$ phần mười, còn trong
$2{,}23$ chỉ có $2$ phần mười.\end{warnbox}
\verify{$2{,}23 = \dfrac{223}{100}$ và $2{,}3 = \dfrac{230}{100}$; $223 < 230$.}
Số $2{,}23$ là phương án \textbf{(D)}.
\ans{(D) $2{,}23$}
}

% ---------------------------------------------------------------- Bài 2
\begin{probbox}{2}{\relax}
Bạn Minh tham gia một chuyến chèo thuyền kayak dài $36$ km. Trước bữa trưa, Minh đã chèo được $12$ km đầu tiên. Hỏi
sau bữa trưa, Minh còn phải chèo bao nhiêu phần của cả chuyến đi?
\choices{$\dfrac{1}{2}$}{$\dfrac{5}{6}$}{$\dfrac{3}{4}$}{$\dfrac{2}{3}$}{$\dfrac{3}{5}$}
\end{probbox}
\probsol{
\textbf{Phân tích đề.} Cả chuyến dài $36$ km, đã chèo $12$ km. Đề hỏi phần \emph{còn lại} (chưa chèo), viết dưới dạng phân số.
\step{1}{Tính quãng đường còn lại}
$36 - 12 = 24$ (km).
\step{2}{Lập phân số ``phần còn lại trên toàn bộ''}
\begin{recallbox}Phân số chỉ ``bao nhiêu phần của toàn thể'' $=\dfrac{\text{phần đang xét}}{\text{toàn thể}}$. Rút gọn bằng cách
chia cả tử và mẫu cho ước chung lớn nhất.\end{recallbox}
\[ \frac{24}{36} = \frac{24 : 12}{36 : 12} = \frac{2}{3}. \]
\step{3}{Minh họa bằng sơ đồ}
Chia $36$ km thành $3$ đoạn, mỗi đoạn $12$ km: Minh đã chèo $1$ đoạn, còn $2$ đoạn.
\begin{barfigC}
  \blk{0}{2.4}{0}{0.6}{tmBlue}{$12$ km}
  \blk{2.4}{4.8}{0}{0.6}{tmOrange}{$12$ km}
  \blk{4.8}{7.2}{0}{0.6}{tmOrange}{$12$ km}
  \brcup{0}{2.4}{0.5}{trước bữa trưa}
  \brcup{2.4}{7.2}{0.5}{sau bữa trưa: $\frac{2}{3}$}
\end{barfigC}
\begin{warnbox}Phân số $\dfrac{12}{36} = \dfrac{1}{3}$ là phần \emph{đã chèo}. Đề hỏi phần \emph{còn lại}, nên đáp án là
$1 - \dfrac13 = \dfrac23$.\end{warnbox}
\verify{$\dfrac{2}{3}\times 36 = 24$ km; $24 + 12 = 36$ km. Đúng.}
\begin{parentbox}Hãy hỏi con hai câu: ``Đã làm được bao nhiêu phần?'' và ``Còn lại bao nhiêu phần?''. Tập thói quen
trả lời cả hai câu giúp con tránh được bẫy chọn nhầm $\dfrac13$.\end{parentbox}
Phân số $\dfrac{2}{3}$ là phương án \textbf{(D)}.
\ans{(D) $\dfrac{2}{3}$}
}

% ---------------------------------------------------------------- Bài 3
\begin{probbox}{3}{\relax}
Khẳng định nào sau đây là \textbf{đúng}?
\choicesII{$0$ nhỏ hơn $-5$}{$7$ nhỏ hơn $-1$}{$10$ nhỏ hơn $\dfrac{1}{4}$}{$-1$ nhỏ hơn $-3$}{$-8$ nhỏ hơn $-2$}
\end{probbox}
\probsol{
\textbf{Phân tích đề.} Mỗi phương án là một phép so sánh. Ta kiểm tra từng phương án, tìm phương án đúng.
\begin{recallbox}Trên trục số nằm ngang, số bên trái nhỏ hơn số bên phải. Mọi số âm nhỏ hơn $0$ và nhỏ hơn mọi số dương.
Trong hai số âm, số nào xa $0$ hơn thì nhỏ hơn.\end{recallbox}
\step{1}{Đặt các số lên trục số}
\begin{figC}
  \draw[->,thick,Navy] (-9.4*0.6,0)--(10.6*0.6,0);
  \foreach \x in {-9,...,10}{\draw[Navy] (\x*0.6,0.08)--(\x*0.6,-0.08);}
  \foreach \x in {-8,-5,-3,-2,-1,0,7,10}{\node[below,font=\scriptsize] at (\x*0.6,-0.1) {$\x$};}
  \foreach \x in {-8,-2}{\fill[tmOrange] (\x*0.6,0) circle (3pt);}
  \draw[->,tmGreen!70!black,thick] (-8*0.6,0.35) to[bend left=30] node[above,font=\scriptsize]{sang phải $6$ đơn vị} (-2*0.6,0.35);
\end{figC}
\step{2}{Kiểm tra từng phương án}
(A) $0$ nhỏ hơn $-5$? Sai: $-5$ nằm bên trái $0$, nên $-5 < 0$.\\
(B) $7$ nhỏ hơn $-1$? Sai: số dương luôn lớn hơn số âm.\\
(C) $10$ nhỏ hơn $\dfrac14$? Sai: $\dfrac14 = 0{,}25 < 10$.\\
(D) $-1$ nhỏ hơn $-3$? Sai: $-3$ xa $0$ hơn nên $-3 < -1$.\\
(E) $-8$ nhỏ hơn $-2$? \textbf{Đúng}: $-8$ nằm bên trái $-2$ (xa $0$ hơn).
\begin{warnbox}Nhiều bạn chọn (D) vì nghĩ ``$1 < 3$ nên $-1 < -3$''. Với số âm thì ngược lại: $-1 > -3$
(nhiệt độ $-1^\circ$C ấm hơn $-3^\circ$C).\end{warnbox}
\verify{$-2 - (-8) = 6 > 0$, nghĩa là $-2$ lớn hơn $-8$.}
Khẳng định đúng duy nhất là \textbf{(E)}.
\ans{(E) $-8$ nhỏ hơn $-2$}
}

% ---------------------------------------------------------------- Bài 4
\begin{probbox}{4}{\relax}
Giá trị của biểu thức $6 \times 2 - 3$ là
\choices{$9$}{$-6$}{$12$}{$15$}{$10$}
\end{probbox}
\probsol{
\textbf{Phân tích đề.} Biểu thức có hai phép tính: nhân và trừ, không có ngoặc.
\begin{recallbox}Trong biểu thức không có ngoặc: nhân, chia trước; cộng, trừ sau.\end{recallbox}
\step{1}{Thực hiện phép nhân trước}
$6\times 2 = 12$.
\step{2}{Rồi thực hiện phép trừ}
$12 - 3 = 9$.
\begin{warnbox}Nếu làm phép trừ trước, ta sẽ được $6\times(2-3) = 6\times(-1) = -6$ -- đó là phương án (B), một ``bẫy''
dành cho bạn quên thứ tự thực hiện phép tính. Biểu thức đề bài \emph{không có ngoặc}, nên phải nhân trước.\end{warnbox}
\verify{Đọc biểu thức như một câu chuyện: ``$6$ hộp, mỗi hộp $2$ cái bánh, ăn mất $3$ cái'' thì còn $12 - 3 = 9$ cái.}
Giá trị $9$ là phương án \textbf{(A)}.
\ans{(A) $9$}
}

% ---------------------------------------------------------------- Bài 5
\begin{probbox}{5}{\relax}
Bài kiểm tra có tổng điểm là $50$. Bạn Mai đạt $80\%$ số điểm. Mai được bao nhiêu điểm?
\choices{$40$}{$62{,}5$}{$10$}{$45$}{$35$}
\end{probbox}
\probsol{
\textbf{Phân tích đề.} Tổng điểm tối đa là $50$; Mai đạt $80\%$ của $50$. Cần tìm số điểm của Mai.
\begin{recallbox}Muốn tìm $p\%$ của số $a$ ta tính $a \times \dfrac{p}{100}$.\end{recallbox}
\step{1}{Viết phép tính}
$80\%$ của $50$ là $50 \times \dfrac{80}{100}$.
\step{2}{Tính}
$50 \times \dfrac{80}{100} = \dfrac{50\times 80}{100} = \dfrac{4000}{100} = 40$.
\step{3}{Cách nhẩm nhanh}
$10\%$ của $50$ là $5$; vậy $80\% = 8\times 10\%$ là $8\times 5 = 40$.
\begin{warnbox}Phương án (B) $62{,}5$ là kết quả của phép chia $50 : 0{,}8$ -- sai, vì $62{,}5$ còn lớn hơn cả điểm tối đa $50$!
Phương án (C) $10$ là số điểm Mai \emph{bị mất} ($20\%$ của $50$).\end{warnbox}
\verify{$\dfrac{40}{50} = \dfrac{80}{100} = 80\%$. Đúng.}
Số điểm $40$ là phương án \textbf{(A)}.
\ans{(A) $40$}
}

% ---------------------------------------------------------------- Bài 6
\begin{probbox}{6}{\relax}
Thịt bò xay được bán với giá $120\,000$ đồng một ki-lô-gam. Mua $12$ kg thịt bò xay thì hết bao nhiêu tiền?
\choicesII{$120\,000$ đồng}{$12\,000$ đồng}{$1\,440\,000$ đồng}{$132\,000$ đồng}{$10\,000$ đồng}
\end{probbox}
\probsol{
\textbf{Phân tích đề.} Đơn giá (giá của $1$ kg) là $120\,000$ đồng; số lượng là $12$ kg. Cần tìm thành tiền.
\begin{recallbox}Thành tiền $=$ đơn giá $\times$ số lượng.\end{recallbox}
\step{1}{Viết phép tính}
Thành tiền $= 120\,000 \times 12$ (đồng).
\step{2}{Tính (tách $12 = 10 + 2$)}
$120\,000\times 10 = 1\,200\,000$; \quad $120\,000 \times 2 = 240\,000$.\\
$120\,000 \times 12 = 1\,200\,000 + 240\,000 = 1\,440\,000$ (đồng).
\step{3}{Ước lượng để kiểm tra}
$12$ kg, mỗi kg hơn $100$ nghìn đồng, nên phải trả hơn $1\,200\,000$ đồng. Chỉ có phương án (C) thỏa mãn.
\begin{warnbox}Các phương án sai đến từ việc dùng nhầm phép tính: (A) là giá $1$ kg; (D) là $120\,000 + 12\,000$ (cộng thay vì
nhân); (E) là $120\,000 : 12$ (chia thay vì nhân).\end{warnbox}
\verify{$1\,440\,000 : 12 = 120\,000$ đồng mỗi kg. Đúng.}
Số tiền $1\,440\,000$ đồng là phương án \textbf{(C)}.
\ans{(C) $1\,440\,000$ đồng}
}

% ---------------------------------------------------------------- Bài 7
\begin{probbox}{7}{\relax}
Giá trị của biểu thức $(8 \times 4) + 3$ là
\choices{$96$}{$15$}{$56$}{$35$}{$28$}
\end{probbox}
\probsol{
\textbf{Phân tích đề.} Biểu thức có ngoặc chứa phép nhân, sau đó cộng thêm $3$.
\begin{recallbox}Biểu thức có ngoặc: làm phép tính trong ngoặc trước.\end{recallbox}
\step{1}{Tính trong ngoặc}
$8\times 4 = 32$.
\step{2}{Cộng thêm $3$}
$32 + 3 = 35$.
\begin{warnbox}Đọc nhầm thành $8 \times (4 + 3) = 8\times 7 = 56$ sẽ dẫn đến phương án (C). Ngoặc trong đề bao quanh
$8\times 4$, không phải $4 + 3$. Phương án (A) $96 = 8\times 4\times 3$ là do nhầm dấu cộng thành dấu nhân.\end{warnbox}
\verify{Không có ngoặc thì kết quả vẫn là $8\times 4 + 3 = 35$, vì nhân luôn được làm trước cộng; ngoặc ở đây chỉ để nhấn mạnh.}
Giá trị $35$ là phương án \textbf{(D)}.
\ans{(D) $35$}
}

% ---------------------------------------------------------------- Bài 8
\begin{probbox}{8}{\relax}
Số nguyên nào sau đây gần với $99 \times 9$ nhất?
\choices{$10\,000$}{$100$}{$100\,000$}{$1000$}{$10$}
\end{probbox}
\probsol{
\textbf{Phân tích đề.} Các phương án là $10$; $100$; $1000$; $10\,000$; $100\,000$ -- cách nhau gấp $10$ lần. Chỉ cần ước
lượng $99\times 9$ là đủ, nhưng ta cũng tính chính xác để chắc chắn.
\step{1}{Ước lượng}
$99 \approx 100$ nên $99 \times 9 \approx 100\times 9 = 900$.
\step{2}{Tính chính xác bằng tính chất phân phối}
\begin{recallbox}Tính chất phân phối: $(a - b)\times c = a\times c - b\times c$.\end{recallbox}
$99\times 9 = (100 - 1)\times 9 = 900 - 9 = 891$.
\step{3}{So sánh khoảng cách đến các phương án gần nhất}
Khoảng cách từ $891$ đến $1000$ là $1000 - 891 = 109$.\\
Khoảng cách từ $891$ đến $100$ là $891 - 100 = 791$.\\
Các phương án còn lại ($10$; $10\,000$; $100\,000$) còn xa hơn nhiều. Vậy $1000$ là gần nhất.
\begin{figC}
  \draw[->,thick,Navy] (-0.3,0)--(10.6,0);
  \foreach \x/\l in {0/0,1/100,10/1000}{\draw[Navy] (\x,0.1)--(\x,-0.1); \node[below,font=\scriptsize] at (\x,-0.1) {$\l$};}
  \fill[tmOrange] (8.91,0) circle (3pt); \node[above,font=\scriptsize] at (8.91,0.1) {$891$};
  \draw[<->,tmGreen!70!black] (8.91,-0.55)--(10,-0.55) node[midway,below,font=\scriptsize]{$109$};
  \draw[<->,AnsRed!70] (1,-0.55)--(8.91,-0.55) node[midway,below,font=\scriptsize]{$791$};
\end{figC}
\verify{$891 : 9 = 99$. Đúng.}
Số $1000$ là phương án \textbf{(D)}.
\ans{(D) $1000$}
}

% ---------------------------------------------------------------- Bài 9
\begin{probbox}{9}{\relax}
Bạn Huy đang đọc một cuốn sách dày $396$ trang. Huy mới chỉ đọc xong $\dfrac{1}{3}$ đầu tiên của cuốn sách. Hỏi Huy còn
phải đọc bao nhiêu trang nữa để hết cuốn sách?
\choices{$264$}{$124$}{$250$}{$199$}{$244$}
\end{probbox}
\probsol{
\textbf{Phân tích đề.} Cả cuốn sách: $396$ trang. Đã đọc: $\dfrac13$ cuốn. Hỏi: số trang \emph{còn lại}.
\begin{recallbox}Muốn tìm $\dfrac{m}{n}$ của một số, ta chia số đó cho $n$ rồi nhân với $m$.\end{recallbox}
\step{1}{Tính số trang đã đọc}
$\dfrac13$ của $396$ là $396 : 3 = 132$ (trang).
\step{2}{Tính số trang còn lại}
$396 - 132 = 264$ (trang).
\step{3}{Cách khác: tính trực tiếp phần còn lại}
Phần còn lại là $1 - \dfrac13 = \dfrac23$ cuốn sách, tức $396 : 3 \times 2 = 132\times 2 = 264$ (trang).
\begin{barfigC}
  \blk{0}{2.6}{0}{0.6}{tmBlue}{$132$ trang}
  \blk{2.6}{5.2}{0}{0.6}{tmOrange}{$132$ trang}
  \blk{5.2}{7.8}{0}{0.6}{tmOrange}{$132$ trang}
  \brcup{0}{2.6}{0.5}{đã đọc}
  \brcup{2.6}{7.8}{0.5}{còn phải đọc}
  \brcdn{0}{7.8}{-0.5}{$396$ trang}
\end{barfigC}
\begin{warnbox}Dừng lại ở $132$ là trả lời nhầm câu hỏi (đó là số trang \emph{đã đọc}; và $132$ cũng không có trong các
phương án). Luôn đọc lại câu hỏi trước khi chọn đáp án.\end{warnbox}
\verify{$264 + 132 = 396$, và $132\times 3 = 396$. Đúng.}
Số $264$ là phương án \textbf{(A)}.
\ans{(A) $264$}
}

% ---------------------------------------------------------------- Bài 10
\begin{probbox}{10}{\relax}
Số nào sau đây là bội của $9$?
\choices{$50$}{$40$}{$35$}{$45$}{$55$}
\end{probbox}
\probsol{
\textbf{Phân tích đề.} Cần tìm số chia hết cho $9$ (là bội của $9$) trong năm số đã cho.
\begin{recallbox}Một số chia hết cho $9$ khi và chỉ khi tổng các chữ số của nó chia hết cho $9$.\end{recallbox}
\step{1}{Tính tổng các chữ số của từng số}
\begin{center}
\begin{tabular}{|c|c|c|c|c|c|}
\hline
\rowcolor{tmBlue!10} Số & $50$ & $40$ & $35$ & $45$ & $55$\\ \hline
Tổng các chữ số & $5$ & $4$ & $8$ & $9$ & $10$\\ \hline
Chia hết cho $9$? & không & không & không & \textbf{có} & không\\ \hline
\end{tabular}
\end{center}
\step{2}{Kiểm tra bằng bảng nhân $9$}
Các bội của $9$ gần đó: $36,\ 45,\ 54$. Chỉ có $45 = 9\times 5$ xuất hiện trong các phương án.
\begin{proofbox}Vì sao dấu hiệu đúng với số có hai chữ số $\nd{ab}$? Ta có $\nd{ab} = 10a + b = 9a + (a + b)$. Phần $9a$
chia hết cho $9$, nên $\nd{ab}$ chia hết cho $9$ đúng khi $a + b$ chia hết cho $9$. Với $45$: $45 = 9\times 4 + 9$.\end{proofbox}
\verify{$45 : 9 = 5$, không dư.}
Số $45$ là phương án \textbf{(D)}.
\ans{(D) $45$}
}

% ---------------------------------------------------------------- Bài 11
\begin{probbox}{11}{\relax}
Cho tập hợp $\{0{,}34;\ 0{,}304;\ 0{,}034;\ 0{,}43\}$. Tổng của số nhỏ nhất và số lớn nhất trong tập hợp là
\choices{$0{,}77$}{$0{,}734$}{$0{,}077$}{$0{,}464$}{$0{,}338$}
\end{probbox}
\probsol{
\textbf{Phân tích đề.} Bước 1: tìm số nhỏ nhất và số lớn nhất. Bước 2: cộng hai số đó.
\begin{recallbox}Để so sánh số thập phân, thêm chữ số $0$ cho các phần thập phân dài bằng nhau rồi so sánh từng hàng
từ trái sang phải.\end{recallbox}
\step{1}{Viết các số với ba chữ số thập phân}
\begin{center}
\begin{tabular}{|c|c|c|c|}
\hline
\rowcolor{tmBlue!10} Số ban đầu & Viết lại & Phần mười & Phần trăm \\ \hline
$0{,}34$ & $0{,}340$ & $3$ & $4$ \\ \hline
$0{,}304$ & $0{,}304$ & $3$ & $0$ \\ \hline
$0{,}034$ & $0{,}034$ & $0$ & $3$ \\ \hline
$0{,}43$ & $0{,}430$ & $4$ & $3$ \\ \hline
\end{tabular}
\end{center}
\step{2}{Tìm số nhỏ nhất và lớn nhất}
So sánh như số tự nhiên: $34 < 304 < 340 < 430$, nên $0{,}034 < 0{,}304 < 0{,}34 < 0{,}43$.\\
Số nhỏ nhất: $0{,}034$ (hàng phần mười là $0$). Số lớn nhất: $0{,}43$ (hàng phần mười là $4$).
\step{3}{Cộng, đặt dấu phẩy thẳng cột}
\[ \begin{array}{r} 0{,}430\\ +\ 0{,}034\\ \hline 0{,}464 \end{array} \]
\begin{warnbox}Đặt tính lệch hàng (cộng $43$ với $34$ ``theo hàng cuối'') sẽ cho $0{,}77$ hoặc $0{,}077$ -- phương án (A)
và (C). Phải để dấu phẩy thẳng cột: $4$ phần mười cộng $0$ phần mười, $3$ phần trăm cộng $3$ phần trăm, \ldots\end{warnbox}
\verify{$0{,}464 - 0{,}43 = 0{,}034$. Đúng.}
Tổng $0{,}464$ là phương án \textbf{(D)}.
\ans{(D) $0{,}464$}
}

% ---------------------------------------------------------------- Bài 12
\begin{probbox}{12}{\relax}
Lan và Hà chạy trên một đường chạy vòng tròn dài $150$ m. Lan chạy $3$ vòng hết $120$ giây, còn Hà chạy $5$ vòng hết
$160$ giây. Bạn nào chạy nhanh hơn, và vận tốc của bạn đó xấp xỉ bao nhiêu?
\choicesII{Lan, $3{,}75$ m/s}{Lan, $2{,}4$ m/s}{Hà, $3{,}3$ m/s}{Hà, $4{,}69$ m/s}{Hà, $3{,}75$ m/s}
\end{probbox}
\probsol{
\textbf{Phân tích đề.} Một vòng dài $150$ m. Lan: $3$ vòng trong $120$ s. Hà: $5$ vòng trong $160$ s. Cần tính vận tốc
(m/s) của mỗi bạn, so sánh, rồi chọn phương án.
\begin{recallbox}Vận tốc $=$ quãng đường $:$ thời gian. Quãng đường chạy $n$ vòng $= n\times$ độ dài một vòng.\end{recallbox}
\step{1}{Quãng đường mỗi bạn chạy}
Lan: $3\times 150 = 450$ (m). \quad Hà: $5\times 150 = 750$ (m).
\step{2}{Vận tốc của Lan}
$v_{\text{Lan}} = 450 : 120 = 3{,}75$ (m/s).
\step{3}{Vận tốc của Hà}
$v_{\text{Hà}} = 750 : 160 = 4{,}6875 \approx 4{,}69$ (m/s).
\step{4}{So sánh}
$4{,}69 > 3{,}75$ nên \textbf{Hà chạy nhanh hơn}, với vận tốc khoảng $4{,}69$ m/s.
\begin{center}
\begin{tabular}{|c|c|c|c|}
\hline
\rowcolor{tmBlue!10} & Quãng đường & Thời gian & Vận tốc\\ \hline
Lan & $450$ m & $120$ s & $3{,}75$ m/s\\ \hline
Hà & $750$ m & $160$ s & $\approx 4{,}69$ m/s\\ \hline
\end{tabular}
\end{center}
\begin{warnbox}Phương án (E) ``Hà, $3{,}75$ m/s'' ghép đúng tên người nhanh hơn với vận tốc của người chậm hơn. Phương án (B)
$2{,}4$ m/s là kết quả của phép tính ngược $120 : 50$ (chia thời gian cho quãng đường). Hãy luôn viết rõ đơn vị:
m $:$ s $=$ m/s.\end{warnbox}
\verify{$4{,}6875 \times 160 = 750$ m $= 5$ vòng; $3{,}75\times 120 = 450$ m $= 3$ vòng. Đúng.}
Kết luận ``Hà, $4{,}69$ m/s'' là phương án \textbf{(D)}.
\ans{(D) Hà, $4{,}69$ m/s}
}

% ---------------------------------------------------------------- Bài 13
\begin{probbox}{13}{\relax}
Trong phép cộng dưới đây, $P$ và $Q$ là hai chữ số. Giá trị của $P + Q$ là
\[ \begin{array}{r@{\;}c@{\;}c@{\;}c} & P & Q & Q\\ & P & P & Q\\ + & Q & Q & Q\\ \hline & 8 & 7 & 6 \end{array} \]
\choices{$3$}{$5$}{$7$}{$6$}{$4$}
\end{probbox}
\probsol{
\textbf{Phân tích đề.} Ba số $\nd{PQQ}$, $\nd{PPQ}$, $\nd{QQQ}$ có tổng là $876$. Cần tìm hai chữ số $P$, $Q$ rồi tính $P + Q$.
Ta xét từng hàng, bắt đầu từ hàng đơn vị.
\step{1}{Hàng đơn vị: $Q + Q + Q = 3Q$ có chữ số tận cùng là $6$}
$Q$ là chữ số nên $3Q$ có thể là $0, 3, 6, \ldots, 27$. Các số tận cùng bằng $6$ trong đó: chỉ có $6$ (vì $16$ và $26$
không chia hết cho $3$). Vậy $3Q = 6$, $Q = 2$, và \emph{không có nhớ} sang hàng chục.
\step{2}{Hàng chục: $Q + P + Q = P + 4$ có chữ số tận cùng là $7$}
$P + 4$ nằm trong khoảng $4$ đến $13$; tận cùng bằng $7$ thì $P + 4 = 7$, tức $P = 3$. Không có nhớ sang hàng trăm.
\step{3}{Hàng trăm: kiểm tra $P + P + Q = 8$}
$3 + 3 + 2 = 8$. Khớp với chữ số hàng trăm của $876$.
\begin{recallbox}Cấu tạo số: $\nd{PQQ} = 100P + 11Q$, \ $\nd{PPQ} = 110P + Q$, \ $\nd{QQQ} = 111Q$.\end{recallbox}
\step{4}{Cách khác bằng cấu tạo số}
Cộng lại: $(100P + 11Q) + (110P + Q) + 111Q = 210P + 123Q = 876$. Thử $Q = 2$: $210P = 876 - 246 = 630$, nên $P = 3$.
(Các giá trị khác của $Q$ không cho $P$ là chữ số.)
\verify{$322 + 332 + 222 = 876$. Đúng.}
\begin{parentbox}Bài đố chữ giống một trò chơi trinh thám: mỗi hàng là một ``manh mối''. Hãy cho con viết riêng từng
hàng thành một câu (``ba lần $Q$ tận cùng bằng $6$'') rồi thử các chữ số từ $0$ đến $9$. Đừng quên kiểm tra số nhớ.\end{parentbox}
Vậy $P + Q = 3 + 2 = 5$, là phương án \textbf{(B)}.
\ans{(B) $5$}
}

% ---------------------------------------------------------------- Bài 14
\begin{probbox}{14}{\relax}
Trung bình cộng của bốn số nguyên $78$; $83$; $82$ và $x$ là $80$. Khẳng định nào sau đây đúng?
\choicesL{$x$ lớn hơn số trung bình cộng $2$ đơn vị}{$x$ nhỏ hơn số trung bình cộng $1$ đơn vị}{$x$ nhỏ hơn số trung bình cộng $2$ đơn vị}{$x$ nhỏ hơn số trung bình cộng $3$ đơn vị}{$x$ bằng số trung bình cộng}
\end{probbox}
\probsol{
\textbf{Phân tích đề.} Biết trung bình cộng của bốn số là $80$ và ba trong bốn số. Tìm $x$ rồi so sánh với $80$.
\begin{recallbox}Tổng các số $=$ trung bình cộng $\times$ số các số.\end{recallbox}
\step{1}{Tổng của bốn số}
$80\times 4 = 320$.
\step{2}{Tìm $x$}
$78 + 83 + 82 = 243$, nên $x = 320 - 243 = 77$.
\step{3}{So sánh với trung bình cộng}
$80 - 77 = 3$: $x$ nhỏ hơn trung bình cộng $3$ đơn vị.
\step{4}{Cách nhanh -- phương pháp độ lệch}
So với $80$: $78$ thiếu $2$; $83$ thừa $3$; $82$ thừa $2$. Tổng độ lệch của ba số: $-2 + 3 + 2 = +3$.
Vì tổng tất cả các độ lệch phải bằng $0$, nên $x$ phải ``thiếu'' $3$: $x = 80 - 3 = 77$.
\begin{proofbox}Mỗi số viết được thành $80 + $(độ lệch). Tổng bốn số là $4\times 80 + $(tổng độ lệch). Mà tổng bốn số đúng bằng
$4\times 80$, nên tổng độ lệch bằng $0$.\end{proofbox}
\verify{$(78 + 83 + 82 + 77) : 4 = 320 : 4 = 80$. Đúng.}
\begin{parentbox}Cách ``độ lệch'' rất hợp để tính nhẩm điểm trung bình trong đời sống. Có thể đố con: ``Ba bài
được $7$, $9$, $8$; muốn trung bình $8$ thì bài thứ tư cần mấy điểm?'' (Đáp: $8$, vì $-1 + 1 + 0 = 0$.)\end{parentbox}
Khẳng định ``$x$ nhỏ hơn số trung bình cộng $3$ đơn vị'' là phương án \textbf{(D)}.
\ans{(D)}
}

% ---------------------------------------------------------------- Bài 15
\begin{probbox}{15}{\relax}
Trong dãy Fibonacci $1, 1, 2, 3, 5, \ldots$, mỗi số kể từ số $2$ trở đi bằng tổng của hai số đứng ngay trước nó. Chẳng hạn,
số tiếp theo của dãy là $3 + 5 = 8$. Số nào sau đây có mặt trong dãy?
\choices{$20$}{$21$}{$22$}{$23$}{$24$}
\end{probbox}
\probsol{
\textbf{Phân tích đề.} Quy luật: số sau $=$ tổng hai số liền trước. Ta viết tiếp dãy cho đến khi vượt quá $24$ (số lớn
nhất trong các phương án), rồi xem phương án nào xuất hiện.
\step{1}{Viết tiếp dãy}
\begin{center}
\begin{tabular}{|c|c|c|c|c|c|c|c|c|c|}
\hline
\rowcolor{tmBlue!10} Thứ tự & 1 & 2 & 3 & 4 & 5 & 6 & 7 & 8 & 9\\ \hline
Số & $1$ & $1$ & $2$ & $3$ & $5$ & $8$ & $13$ & $\mathbf{21}$ & $34$\\ \hline
\end{tabular}
\end{center}
Cụ thể: $5 + 8 = 13$; \ $8 + 13 = 21$; \ $13 + 21 = 34$.
\step{2}{Đối chiếu với các phương án}
Số $34$ đã lớn hơn $24$, nên các số tiếp theo đều lớn hơn $24$. Giữa $13$ và $34$ chỉ có đúng một số của dãy là $21$.
Vậy $20$, $22$, $23$, $24$ không thuộc dãy; chỉ có $21$ thuộc dãy.
\begin{tipbox}Dãy Fibonacci luôn tăng (từ số thứ hai trở đi) vì ta cộng thêm số dương. Nhờ vậy, khi đã vượt qua một số,
ta biết chắc số đó không thể xuất hiện về sau nữa.\end{tipbox}
\verify{$8 + 13 = 21$. Đúng.}
Số $21$ là phương án \textbf{(B)}.
\ans{(B) $21$}
}

% ---------------- part b ----------------
% ============================================================
% Topic 4 -- Số học và Tư duy số -- part b: problems 16--30 (VI)
% ============================================================

\begin{probbox}{16}{\relax}
Trung bình cộng của năm số nguyên liên tiếp bằng $21$. Số nhỏ nhất trong năm số đó là
\choices{$17$}{$21$}{$1$}{$18$}{$19$}
\end{probbox}
\probsol{
\textbf{Phân tích đề.} Ta có \emph{năm số nguyên liên tiếp} (mỗi số hơn số đứng trước đúng $1$ đơn vị). Biết trung bình cộng của chúng là $21$. Cần tìm số \emph{nhỏ nhất}.

\begin{recallbox}Trung bình cộng của một nhóm số $=$ (tổng các số) $:$ (số các số). Do đó: tổng các số $=$ trung bình cộng $\times$ số các số.\end{recallbox}

\step{1}{Tìm tổng của năm số}
Tổng của năm số là $21\times 5=105$.

\step{2}{Gọi tên năm số theo số ở giữa}
Gọi số đứng giữa (số thứ ba) là $n$. Khi đó năm số liên tiếp là
\[n-2,\quad n-1,\quad n,\quad n+1,\quad n+2.\]
Cộng lại: $(n-2)+(n-1)+n+(n+1)+(n+2)=5n$ (các số $-2$ và $+2$, $-1$ và $+1$ triệt tiêu nhau).

\step{3}{Tìm số ở giữa rồi suy ra số nhỏ nhất}
Ta có $5n=105$, nên $n=105:5=21$. Năm số là $19,\ 20,\ 21,\ 22,\ 23$. Số nhỏ nhất là $n-2=19$.

\begin{figC}
\draw[->,thick,Navy] (-0.5,0)--(6.6,0);
\foreach \x/\lab in {0/18,1/19,2/20,3/21,4/22,5/23,6/24}{\draw[Navy,line width=0.8pt] (\x,0.08)--(\x,-0.08); \node[below,text=Ink] at (\x,-0.1) {$\lab$};}
\foreach \x in {1,2,3,4,5}{\fill[KidOrange] (\x,0) circle (2.7pt);}
\draw[<->,tmBlue] (1,0.45)--(3,0.45) node[midway,above,font=\footnotesize]{lùi $2$};
\draw[<->,tmBlue] (3,0.45)--(5,0.45) node[midway,above,font=\footnotesize]{tiến $2$};
\node[tmGreen!70!black,font=\footnotesize\bfseries] at (3,1.1) {số ở giữa $=$ trung bình cộng};
\end{figC}

\begin{proofbox}Với một số \emph{lẻ} các số nguyên liên tiếp, các số đứng đối xứng hai bên số ở giữa ``bù trừ'' cho nhau (một số kém $1$, một số hơn $1$; một số kém $2$, một số hơn $2$\dots). Vì thế tổng bằng (số ở giữa) $\times$ (số các số), tức là trung bình cộng đúng bằng số ở giữa.\end{proofbox}

\begin{warnbox}Đừng chọn ngay $21$: đó là số \emph{ở giữa}, không phải số nhỏ nhất. Cũng đừng lấy $21-4=17$: từ số giữa đến số nhỏ nhất chỉ lùi $2$ bước, không phải $4$.\end{warnbox}

\verify{$19+20+21+22+23=105$ và $105:5=21$. Đúng.}
Số nhỏ nhất là $19$, ứng với phương án \textbf{(E)}.
\ans{(E) $19$}
}

\begin{probbox}{17}{\relax}
Có bao nhiêu số nguyên dương (kể cả số $1$) là ước của cả hai số $40$ và $72$?
\choices{$9$}{$12$}{$4$}{$2$}{$5$}
\end{probbox}
\probsol{
\textbf{Phân tích đề.} Một số ``là ước của cả $40$ và $72$'' nghĩa là $40$ chia hết cho số đó \emph{và} $72$ cũng chia hết cho số đó. Đó chính là các \emph{ước chung} của $40$ và $72$. Ta cần đếm chúng.

\begin{recallbox}Số $d$ là ước của $a$ nếu $a$ chia hết cho $d$. Ước chung của hai số là số vừa là ước của số này vừa là ước của số kia. Ước chung lớn nhất (ƯCLN) là số lớn nhất trong các ước chung.\end{recallbox}

\step{1}{Liệt kê các ước của $40$}
Ta tìm theo từng cặp nhân bằng $40$: $1\times 40$, $2\times 20$, $4\times 10$, $5\times 8$. Vậy
\[\text{Ư}(40)=\{1;\ 2;\ 4;\ 5;\ 8;\ 10;\ 20;\ 40\}.\]

\step{2}{Liệt kê các ước của $72$}
Các cặp nhân bằng $72$: $1\times 72$, $2\times 36$, $3\times 24$, $4\times 18$, $6\times 12$, $8\times 9$. Vậy
\[\text{Ư}(72)=\{1;\ 2;\ 3;\ 4;\ 6;\ 8;\ 9;\ 12;\ 18;\ 24;\ 36;\ 72\}.\]

\step{3}{Chọn các số có mặt trong cả hai danh sách}
So sánh hai danh sách, các số chung là $1,\ 2,\ 4,\ 8$. Vậy có $4$ ước chung.

\begin{tipbox}Có cách nhanh hơn: các ước chung của hai số chính là các ước của ƯCLN của chúng. Ta có $40=2^3\cdot 5$ và $72=2^3\cdot 3^2$, nên $\text{ƯCLN}(40,72)=2^3=8$. Các ước của $8$ là $1,2,4,8$: đúng $4$ số.\end{tipbox}

\begin{proofbox}Nếu $d$ là ước của cả $40$ và $72$ thì $d$ cũng là ước của hiệu $72-40=32$, rồi của $40-32=8$. Ngược lại, mọi ước của $8$ đều là ước của $40=8\cdot 5$ và của $72=8\cdot 9$. Vậy ``ước chung của $40$ và $72$'' và ``ước của $8$'' là cùng một tập hợp.\end{proofbox}

\begin{warnbox}Phương án $12$ là số ước của riêng $72$; phương án $9$ hay $5$ thường do đếm sót hoặc đếm nhầm. Câu hỏi hỏi ước \emph{chung}.\end{warnbox}

\verify{$40:1,\ 40:2=20,\ 40:4=10,\ 40:8=5$; $72:1,\ 72:2=36,\ 72:4=18,\ 72:8=9$ -- tất cả đều chia hết. Còn $3,6,9$ không là ước của $40$; $5,10$ không là ước của $72$.}
Có $4$ số, ứng với phương án \textbf{(C)}.
\ans{(C) $4$}
}

\begin{probbox}{18}{\relax}
Số cần điền vào ô trống để có đẳng thức đúng $\dfrac{3}{7}=\dfrac{\square}{63}$ là
\choices{$27$}{$9$}{$59$}{$63$}{$3$}
\end{probbox}
\probsol{
\textbf{Phân tích đề.} Hai phân số bằng nhau; mẫu số đã đổi từ $7$ thành $63$. Ta cần tìm tử số mới.

\begin{recallbox}\textbf{Tính chất cơ bản của phân số:} nếu nhân cả tử và mẫu của một phân số với cùng một số khác $0$ thì được một phân số bằng nó: $\dfrac{a}{b}=\dfrac{a\cdot m}{b\cdot m}\ (m\neq 0)$.\end{recallbox}

\step{1}{Mẫu số đã được nhân với bao nhiêu?}
$63:7=9$, nên mẫu $7$ đã được nhân với $9$ để thành $63$.

\step{2}{Nhân tử số với cùng số đó}
Để phân số không đổi, tử số cũng phải nhân với $9$: $3\times 9=27$. Vậy
\[\frac{3}{7}=\frac{3\cdot 9}{7\cdot 9}=\frac{27}{63}.\]

\begin{warnbox}Phương án $9$ là \emph{số nhân}, không phải tử số mới. Phương án $59$ đến từ lỗi ``cộng thêm'': $63-7=56$ rồi $3+56=59$ -- cộng cùng một số vào tử và mẫu \emph{không} giữ nguyên phân số.\end{warnbox}

\verify{Nhân chéo: $3\times 63=189$ và $7\times 27=189$. Hai tích bằng nhau nên hai phân số bằng nhau.}
Số cần điền là $27$, ứng với phương án \textbf{(A)}.
\ans{(A) $27$}
}

\begin{probbox}{19}{\relax}
Một phân số bằng phân số $\dfrac{5}{8}$. Tổng của tử số và mẫu số của nó bằng $91$. Hiệu của mẫu số và tử số của phân số đó là bao nhiêu?
\choices{$21$}{$3$}{$33$}{$13$}{$19$}
\end{probbox}
\probsol{
\textbf{Phân tích đề.} Phân số cần tìm bằng $\dfrac58$, nên tử số và mẫu số ``chia phần'' theo tỉ lệ $5:8$. Biết tử $+$ mẫu $=91$. Cần tìm mẫu $-$ tử.

\begin{recallbox}Mọi phân số bằng $\dfrac{5}{8}$ đều có dạng $\dfrac{5k}{8k}$ với $k$ là số nguyên dương nào đó (nhân cả tử và mẫu với $k$).\end{recallbox}

\step{1}{Vẽ sơ đồ đoạn thẳng}
Coi tử số gồm $5$ phần bằng nhau, mẫu số gồm $8$ phần như thế.
\begin{barfig}
\rowlab{-0.15}{0.5}{Tử số}
\foreach \i in {0,...,4}{\blk{\i*0.8}{\i*0.8+0.8}{0.5}{0.5}{tmBlue}{}}
\rowlab{-0.15}{-0.3}{Mẫu số}
\foreach \i in {0,...,7}{\blk{\i*0.8}{\i*0.8+0.8}{-0.3}{0.5}{tmOrange}{}}
\brcup{0}{4}{0.9}{$5$ phần}
\brcdn{0}{6.4}{-0.7}{$8$ phần}
\vbrace{6.6}{0.75}{-0.55}{tổng $=91$}
\end{barfig}

\step{2}{Tìm giá trị một phần}
Tổng số phần là $5+8=13$ phần, ứng với $91$. Một phần là $91:13=7$.

\step{3}{Tìm tử số, mẫu số và hiệu}
Tử số $=5\times 7=35$; mẫu số $=8\times 7=56$. Phân số là $\dfrac{35}{56}$.\\
Hiệu mẫu $-$ tử $=56-35=21$ (cũng chính là $3$ phần: $3\times 7=21$).

\begin{warnbox}Phương án $3$ là hiệu $8-5$ của phân số $\dfrac58$ ban đầu, chưa nhân với giá trị một phần. Phương án $13$ là số phần, không phải hiệu.\end{warnbox}

\verify{$35+56=91$; rút gọn $\dfrac{35}{56}=\dfrac{35:7}{56:7}=\dfrac58$. Đúng.}
Hiệu là $21$, ứng với phương án \textbf{(A)}.
\ans{(A) $21$}
}

\begin{probbox}{20}{\relax}
Chị Lan lái ô tô với vận tốc $80$ km/h trên một con đường có các cột điện dựng cách đều nhau. Cứ $4$ giây xe lại đi qua một cột điện. Khoảng cách giữa hai cột điện liền nhau gần nhất với giá trị nào?
\choices{$50$ m}{$60$ m}{$70$ m}{$80$ m}{$90$ m}
\end{probbox}
\probsol{
\textbf{Phân tích đề.} Trong $4$ giây, xe đi được đúng khoảng cách giữa hai cột liền nhau. Vậy ta cần tính quãng đường xe đi trong $4$ giây khi vận tốc là $80$ km/h. Chú ý: vận tốc tính theo \emph{km và giờ}, còn thời gian tính theo \emph{giây}, đáp án tính theo \emph{mét} -- phải đổi đơn vị.

\begin{recallbox}Quãng đường $=$ vận tốc $\times$ thời gian. $1$ km $=1000$ m; $1$ giờ $=60$ phút $=3600$ giây.\end{recallbox}

\step{1}{Đổi vận tốc ra mét mỗi giây}
$80$ km/h nghĩa là trong $3600$ giây xe đi $80\,000$ m. Trong $1$ giây xe đi
\[80\,000:3600=\frac{80\,000}{3600}=\frac{200}{9}\approx 22{,}2\ \text{m}.\]

\step{2}{Tính quãng đường đi trong $4$ giây}
\[\frac{200}{9}\times 4=\frac{800}{9}\approx 88{,}9\ \text{m}.\]

\step{3}{So sánh với các phương án}
$88{,}9$ m cách $90$ m khoảng $1{,}1$ m, cách $80$ m khoảng $8{,}9$ m. Vậy gần nhất là $90$ m.

\begin{tipbox}Cách nghĩ khác: trong $1$ giờ $=3600$ giây xe đi qua $3600:4=900$ khoảng cách giữa hai cột. $900$ khoảng dài tổng cộng $80\,000$ m, nên mỗi khoảng dài $80\,000:900\approx 88{,}9$ m.\end{tipbox}

\begin{warnbox}Nếu quên đổi đơn vị và lấy $80\times 4=320$ thì sẽ không khớp phương án nào. Nếu đổi $1$ giờ $=60$ giây (sai, phải là $3600$ giây) cũng ra kết quả vô lí.\end{warnbox}

\verify{$90$ m mỗi $4$ giây là $22{,}5$ m/s, tức $22{,}5\times 3600=81\,000$ m/h $=81$ km/h, rất gần $80$ km/h. Hợp lí.}
Chọn phương án \textbf{(E)}.
\ans{(E) $90$ m}
}

\begin{probbox}{21}{\relax}
Các số $334$ và $419$ đều có tích các chữ số bằng $36$ (vì $3\cdot 3\cdot 4=36$ và $4\cdot 1\cdot 9=36$). Có tất cả bao nhiêu số nguyên dương có ba chữ số mà tích các chữ số bằng $36$?
\choices{$21$}{$15$}{$18$}{$24$}{$12$}
\end{probbox}
\probsol{
\textbf{Phân tích đề.} Một số có ba chữ số $\nd{abc}$ có tích các chữ số $a\cdot b\cdot c=36$. Không chữ số nào được bằng $0$ (nếu có chữ số $0$ thì tích bằng $0$). Cách làm gọn gàng: \emph{trước hết} tìm các bộ ba chữ số (không kể thứ tự) có tích $36$, \emph{sau đó} đếm xem mỗi bộ xếp được thành bao nhiêu số.

\step{1}{Tìm mọi bộ ba chữ số (từ $1$ đến $9$) có tích bằng $36$}
Ta sắp mỗi bộ theo thứ tự $a\le b\le c$ và xét chữ số nhỏ nhất $a$.
\begin{itemize}[leftmargin=1.4em,itemsep=1pt]
\item $a=1$: cần $b\cdot c=36$ với $1\le b\le c\le 9$: $36=4\cdot 9=6\cdot 6$ (còn $1\cdot 36$, $2\cdot 18$, $3\cdot 12$ có số lớn hơn $9$). Được $\{1,4,9\}$ và $\{1,6,6\}$.
\item $a=2$: cần $b\cdot c=18$ với $2\le b\le c\le 9$: $18=2\cdot 9=3\cdot 6$. Được $\{2,2,9\}$ và $\{2,3,6\}$.
\item $a=3$: cần $b\cdot c=12$ với $3\le b\le c$: $12=3\cdot 4$. Được $\{3,3,4\}$.
\item $a\ge 4$: khi đó $a\cdot b\cdot c\ge 4\cdot 4\cdot 4=64>36$, không có bộ nào.
\end{itemize}
Vậy có đúng $5$ bộ: $\{1,4,9\},\ \{1,6,6\},\ \{2,2,9\},\ \{2,3,6\},\ \{3,3,4\}$.

\begin{recallbox}\textbf{Đếm cách xếp:} ba chữ số \emph{khác nhau} xếp được $3\cdot 2\cdot 1=6$ số (chọn chữ số hàng trăm: $3$ cách, hàng chục: $2$ cách, hàng đơn vị: $1$ cách). Ba chữ số mà có \emph{hai chữ số giống nhau} chỉ xếp được $3$ số (chỉ cần chọn vị trí cho chữ số khác biệt).\end{recallbox}

\step{2}{Đếm số các số tạo được từ mỗi bộ}
\begin{center}
\begin{tabular}{|c|c|l|c|}
\hline
\rowcolor{Tint}\textbf{Bộ chữ số} & \textbf{Loại} & \textbf{Các số tạo được} & \textbf{Số lượng}\\
\hline
$\{1,4,9\}$ & khác nhau & $149,194,419,491,914,941$ & $6$\\
$\{2,3,6\}$ & khác nhau & $236,263,326,362,623,632$ & $6$\\
$\{1,6,6\}$ & có $2$ số giống & $166,616,661$ & $3$\\
$\{2,2,9\}$ & có $2$ số giống & $229,292,922$ & $3$\\
$\{3,3,4\}$ & có $2$ số giống & $334,343,433$ & $3$\\
\hline
\multicolumn{3}{|r|}{\textbf{Tổng}} & $\mathbf{21}$\\
\hline
\end{tabular}
\end{center}

\step{3}{Cộng lại}
$6+6+3+3+3=21$ số.

\begin{warnbox}Hai lỗi hay gặp: (1) quên bộ có chữ số $1$ (như $\{1,4,9\}$, $\{1,6,6\}$) vì chỉ nghĩ đến cách tách $36$ thành ba số ``lớn hơn $1$''; (2) coi bộ có hai chữ số giống nhau cũng xếp được $6$ số -- khi đó sẽ đếm trùng (ví dụ $334$ bị đếm hai lần).\end{warnbox}

\verify{Bảng trên liệt kê đủ $21$ số, có cả hai số $334$ và $419$ mà đề bài nêu. Kiểm tra bằng máy tính (duyệt tất cả các số từ $100$ đến $999$) cũng cho đúng $21$ số.}
Chọn phương án \textbf{(A)}.
\ans{(A) $21$}
}

\begin{probbox}{22}{\relax}
Một \emph{dãy cách đều} là dãy số mà mỗi số hạng (từ số hạng thứ hai) bằng số hạng đứng trước cộng thêm một số không đổi. Ví dụ: $2,4,6,8$ và $1,4,7,10$ là các dãy cách đều.

Trong bảng dưới đây, các số trên mỗi hàng (đọc từ trái sang phải) tạo thành một dãy cách đều, và các số trên mỗi cột (đọc từ trên xuống) cũng tạo thành một dãy cách đều. Giá trị của $x$ là
\begin{figC}
\foreach \rr in {0,1,2,3}{\foreach \cc in {0,1,2,3}{\draw[Navy,line width=0.9pt,fill=white] (\cc*1.1,-\rr*0.8) rectangle (\cc*1.1+1.1,-\rr*0.8-0.8);}}
\node at (0.55,-0.4) {$1$};
\node at (0.55,-1.2) {$4$}; \node at (3.85,-1.2) {$25$};
\node at (0.55,-2.0) {$7$}; \node at (3.85,-2.0) {$x$};
\node at (0.55,-2.8) {$10$}; \node at (2.75,-2.8) {$36$};
\end{figC}
\choices{$37$}{$28$}{$36$}{$43{,}75$}{$46$}
\end{probbox}
\probsol{
\textbf{Phân tích đề.} Bảng có $4$ hàng, $4$ cột. Mỗi hàng và mỗi cột là một dãy cách đều. Ta cần tìm số $x$ ở hàng $3$, cột $4$. Ý tưởng: dùng các hàng đã biết đủ hai số để tìm thêm số, rồi dùng cột $4$ để tìm $x$.

\begin{recallbox}Một dãy cách đều có $4$ số hạng viết được là $a,\ a+d,\ a+2d,\ a+3d$ ($d$ là khoảng cách không đổi). Từ số đầu đến số cuối có $3$ khoảng cách, nên $d=(\text{số cuối}-\text{số đầu}):3$. Nếu biết số thứ nhất và số thứ ba thì $d=(\text{số thứ ba}-\text{số thứ nhất}):2$.\end{recallbox}

\step{1}{Hàng $2$: $4,\ \_,\ \_,\ 25$}
Từ $4$ đến $25$ có $3$ khoảng: $d=(25-4):3=21:3=7$. Hàng $2$ là $4,\ 11,\ 18,\ 25$.

\step{2}{Hàng $4$: $10,\ \_,\ 36,\ \_$}
Từ số thứ nhất $10$ đến số thứ ba $36$ có $2$ khoảng: $d=(36-10):2=26:2=13$. Hàng $4$ là $10,\ 23,\ 36,\ 49$. Vậy ô cuối của hàng $4$ là $49$.

\step{3}{Cột $4$: $\_,\ 25,\ x,\ 49$}
Trong cột $4$, số $x$ đứng giữa $25$ và $49$, cách mỗi số một khoảng bằng nhau. Do đó $x$ là trung bình cộng của chúng:
\[x=\frac{25+49}{2}=\frac{74}{2}=37.\]
(Khoảng cách của cột $4$ là $37-25=12$, nên số đầu cột là $25-12=13$.)

\step{4}{Điền đầy đủ bảng để kiểm tra}
\begin{figC}
\foreach \rr in {0,1,2,3}{\foreach \cc in {0,1,2,3}{\draw[Navy,line width=0.9pt,fill=white] (\cc*1.1,-\rr*0.8) rectangle (\cc*1.1+1.1,-\rr*0.8-0.8);}}
\foreach \cc/\val in {0/1,1/5,2/9,3/13}{\node at (\cc*1.1+0.55,-0.4) {$\val$};}
\foreach \cc/\val in {0/4,1/11,2/18,3/25}{\node at (\cc*1.1+0.55,-1.2) {$\val$};}
\foreach \cc/\val in {0/7,1/17,2/27}{\node at (\cc*1.1+0.55,-2.0) {$\val$};}
\node[text=AnsRed,font=\bfseries] at (3.85,-2.0) {$37$};
\foreach \cc/\val in {0/10,1/23,2/36,3/49}{\node at (\cc*1.1+0.55,-2.8) {$\val$};}
\node[anchor=west,font=\footnotesize,text=Slate] at (4.6,-0.4) {$+4$};
\node[anchor=west,font=\footnotesize,text=Slate] at (4.6,-1.2) {$+7$};
\node[anchor=west,font=\footnotesize,text=Slate] at (4.6,-2.0) {$+10$};
\node[anchor=west,font=\footnotesize,text=Slate] at (4.6,-2.8) {$+13$};
\end{figC}
Cột $1$: $1,4,7,10$ (cách $3$); cột $2$: $5,11,17,23$ (cách $6$); cột $3$: $9,18,27,36$ (cách $9$); cột $4$: $13,25,37,49$ (cách $12$). Mọi hàng, mọi cột đều cách đều.

\begin{warnbox}Không được tìm $x$ bằng cách nhìn hàng $3$ trước: hàng $3$ mới chỉ biết một số ($7$), chưa đủ để tìm khoảng cách. Phải đi ``vòng'' qua hàng $2$, hàng $4$ rồi mới đến cột $4$.\end{warnbox}

\verify{Hàng $3$: $7,17,27,37$ cách đều $10$; cột $4$: $25,37,49$ cách đều $12$. Đúng.}
$x=37$, ứng với phương án \textbf{(A)}.
\ans{(A) $37$}
}

\begin{probbox}{23}{\relax}
Trong số có sáu chữ số $\nd{1ABCDE}$, mỗi chữ cái biểu thị một chữ số. Biết rằng $\nd{1ABCDE}\times 3=\nd{ABCDE1}$. Giá trị của $A+B+C+D+E$ là
\choices{$29$}{$26$}{$22$}{$30$}{$28$}
\end{probbox}
\probsol{
\textbf{Phân tích đề.} Khi nhân số $\nd{1ABCDE}$ với $3$, chữ số $1$ ở đầu ``chạy'' xuống cuối, còn khối $ABCDE$ dịch sang trái một hàng. Ta cần tìm các chữ số $A,B,C,D,E$ rồi tính tổng. Có hai cách giải; cả hai đều rất đáng học.

\begin{recallbox}\textbf{Cấu tạo số:} viết thêm một chữ số $c$ vào bên phải số $N$ thì được số $10\cdot N+c$; viết chữ số $1$ vào bên trái một số $N$ có năm chữ số thì được số $100\,000+N$.\end{recallbox}

\textbf{Cách 1: Đặt ẩn cho cả khối.}
\step{1}{Gọi $N=\nd{ABCDE}$ (một số có năm chữ số)}
Khi đó $\nd{1ABCDE}=100\,000+N$ và $\nd{ABCDE1}=10N+1$.
\step{2}{Lập đẳng thức và tìm $N$}
\[3\cdot(100\,000+N)=10N+1\ \Longrightarrow\ 300\,000+3N=10N+1.\]
Bớt $3N$ và bớt $1$ ở cả hai vế: $299\,999=7N$, nên $N=299\,999:7=42\,857$.
\step{3}{Đọc các chữ số}
$A=4,\ B=2,\ C=8,\ D=5,\ E=7$. Tổng: $4+2+8+5+7=26$.

\textbf{Cách 2: Nhân từng hàng từ phải sang trái (như đặt tính).}
\[\begin{array}{r@{\;}c@{\;}c@{\;}c@{\;}c@{\;}c@{\;}c}
 & 1 & A & B & C & D & E\\
\times & & & & & & 3\\ \hline
 & A & B & C & D & E & 1
\end{array}\]
\begin{itemize}[leftmargin=1.4em,itemsep=1pt]
\item Hàng đơn vị: $3\cdot E$ tận cùng là $1$ $\Rightarrow$ $E=7$ ($3\cdot 7=21$, viết $1$ nhớ $2$).
\item Hàng chục: $3\cdot D+2$ tận cùng là $E=7$ $\Rightarrow$ $3D$ tận cùng là $5$ $\Rightarrow$ $D=5$ ($15+2=17$, viết $7$ nhớ $1$).
\item Hàng trăm: $3\cdot C+1$ tận cùng là $D=5$ $\Rightarrow$ $3C$ tận cùng là $4$ $\Rightarrow$ $C=8$ ($24+1=25$, viết $5$ nhớ $2$).
\item Hàng nghìn: $3\cdot B+2$ tận cùng là $C=8$ $\Rightarrow$ $3B$ tận cùng là $6$ $\Rightarrow$ $B=2$ ($6+2=8$, nhớ $0$).
\item Hàng chục nghìn: $3\cdot A+0$ tận cùng là $B=2$ $\Rightarrow$ $A=4$ ($12$, viết $2$ nhớ $1$).
\item Hàng trăm nghìn: $3\cdot 1+1=4$, đúng bằng $A=4$. Phép nhân khớp hoàn toàn.
\end{itemize}
(Mỗi bước chỉ có một chữ số thỏa mãn, vì các tích $3\cdot 0, 3\cdot 1,\dots,3\cdot 9$ có chữ số tận cùng khác nhau.)

\begin{tipbox}Số $142\,857$ là một số rất nổi tiếng: $142\,857\times 2=285\,714$, $\times 3=428\,571$, $\times 4=571\,428$, $\times 5=714\,285$, $\times 6=857\,142$ -- luôn là các chữ số $1,4,2,8,5,7$ ``xoay vòng''. Nó liên quan đến $\dfrac17=0{,}(142857)$.\end{tipbox}

\verify{$142\,857\times 3=428\,571$. Đúng là $\nd{ABCDE1}$ với $ABCDE=42857$.}
$A+B+C+D+E=26$, ứng với phương án \textbf{(B)}.
\ans{(B) $26$}
}

\begin{probbox}{24}{\relax}
Trên mặt phẳng tọa độ, một chấm tròn xuất phát từ điểm $(20;19)$. Ở lần di chuyển thứ nhất, chấm đi $1$ đơn vị theo phương ngang hoặc phương dọc, tức là tới một trong các điểm $(21;19)$, $(19;19)$, $(20;20)$ hoặc $(20;18)$. Ở lần thứ hai, chấm đi $2$ đơn vị theo phương vuông góc với lần thứ nhất (theo chiều nào cũng được). Mỗi lần tiếp theo, quãng di chuyển dài hơn lần trước $1$ đơn vị ($3$ đơn vị, $4$ đơn vị, $5$ đơn vị,\dots) và phải vuông góc với phương của lần ngay trước đó. Chấm dừng lại sau $10$ lần di chuyển. Điểm nào sau đây \textbf{không thể} là vị trí cuối cùng của chấm?
\choices{$(27;33)$}{$(30;40)$}{$(21;21)$}{$(42;44)$}{$(37;37)$}
\end{probbox}
\probsol{
\textbf{Phân tích đề.} ``Vuông góc với lần trước'' nghĩa là các lần di chuyển \emph{xen kẽ}: ngang, dọc, ngang, dọc,\dots\ (hoặc dọc, ngang, dọc,\dots). Độ dài các lần là $1,2,3,\dots,10$. Ta theo dõi riêng sự thay đổi của hoành độ (tọa độ thứ nhất, do các bước ngang) và tung độ (tọa độ thứ hai, do các bước dọc).

\begin{recallbox}Điểm $(x;y)$: $x$ là hoành độ (đi ngang sang phải thì $x$ tăng, sang trái thì $x$ giảm), $y$ là tung độ (đi lên thì $y$ tăng, đi xuống thì $y$ giảm). Tổng của một số lẻ các số lẻ là số lẻ; tổng (hay hiệu) các số chẵn luôn là số chẵn.\end{recallbox}

\step{1}{Chia các bước thành hai nhóm}
Các lần thứ $1,3,5,7,9$ cùng một phương, có độ dài $1,3,5,7,9$ (đều lẻ). Các lần thứ $2,4,6,8,10$ cùng phương còn lại, có độ dài $2,4,6,8,10$ (đều chẵn).
\begin{itemize}[leftmargin=1.4em,itemsep=1pt]
\item Nếu lần đầu đi \emph{ngang}: hoành độ thay đổi $\pm1\pm3\pm5\pm7\pm9$, tung độ thay đổi $\pm2\pm4\pm6\pm8\pm10$.
\item Nếu lần đầu đi \emph{dọc}: ngược lại, hoành độ thay đổi $\pm2\pm4\pm6\pm8\pm10$, tung độ thay đổi $\pm1\pm3\pm5\pm7\pm9$.
\end{itemize}
(Dấu $+$ hay $-$ tùy chấm đi theo chiều nào.)

\step{2}{Tính chất chẵn -- lẻ}
$\pm1\pm3\pm5\pm7\pm9$ là tổng/hiệu của $5$ số lẻ nên luôn là số \emph{lẻ}, và nhiều nhất là $25$. Còn $\pm2\pm4\pm6\pm8\pm10$ luôn là số \emph{chẵn}, nhiều nhất là $30$. Vậy trong hai độ thay đổi (của hoành độ và tung độ) phải có đúng một số lẻ (lấy từ nhóm lẻ) và một số chẵn (lấy từ nhóm chẵn).

\begin{proofbox}Lấy tổng ``toàn dấu cộng'' $1+3+5+7+9=25$. Đổi dấu một số $k$ từ $+$ sang $-$ làm tổng giảm đi $2k$. Vậy mọi giá trị có dạng $25-2s$, trong đó $s$ là tổng các số bị đổi dấu. Tương tự, nhóm chẵn cho các giá trị $30-2s$.\end{proofbox}

\step{3}{Tính độ thay đổi cho từng phương án}
Lấy tọa độ cuối trừ $(20;19)$:
\begin{center}
\small\begin{tabularx}{\linewidth}{|c|c|c|X|}
\hline
\rowcolor{Tint}\textbf{Phương án} & \textbf{Hoành độ thay đổi} & \textbf{Tung độ thay đổi} & \textbf{Có đạt được không?}\\
\hline
(A) $(27;33)$ & $7$ & $14$ & $1+3+5+7-9=7$; $2+4+6-8+10=14$ \\
(B) $(30;40)$ & $10$ & $21$ & phải có $\pm1\pm3\pm5\pm7\pm9=21$ \dots\\
(C) $(21;21)$ & $1$ & $2$ & $1-3+5+7-9=1$; $2-4+6+8-10=2$\\
(D) $(42;44)$ & $22$ & $25$ & $2-4+6+8+10=22$; $1+3+5+7+9=25$\\
(E) $(37;37)$ & $17$ & $18$ & $-1-3+5+7+9=17$; $2+4-6+8+10=18$\\
\hline
\end{tabularx}
\end{center}

\step{4}{Vì sao (B) không thể}
Ở (B), độ thay đổi lẻ là $21$, nên phải có $\pm1\pm3\pm5\pm7\pm9=21$, tức $25-2s=21$, hay $s=2$: phải đổi dấu một nhóm số lẻ có tổng bằng $2$. Nhưng một số lẻ thì không bằng $2$ (chỉ có thể là $1,3,5,\dots$), còn hai số lẻ khác nhau trong nhóm có tổng ít nhất $1+3=4$. Không có cách nào! Vậy chấm không thể dừng ở $(30;40)$.

\step{5}{Một ví dụ đường đi minh họa (phương án D)}
Lần đầu đi lên: lên $1$, phải $2$, lên $3$, trái $4$, lên $5$, phải $6$, lên $7$, phải $8$, lên $9$, phải $10$. Chấm từ $(20;19)$ tới $(42;44)$.
\begin{figC}\begin{scope}[scale=0.17]
\draw[Line,line width=0.3pt] (-3,-1) grid (23,26);
\draw[tmBlue,line width=1.4pt,->] (0,0)--(0,1)--(2,1)--(2,4)--(-2,4)--(-2,9)--(4,9)--(4,16)--(12,16)--(12,25)--(22,25);
\fill[tmGreen] (0,0) circle (10pt); \fill[AnsRed] (22,25) circle (10pt);
\node[below right,font=\scriptsize,text=tmGreen!70!black] at (0,0) {$(20;19)$};
\node[below,font=\scriptsize,text=AnsRed] at (22,24.6) {$(42;44)$};
\end{scope}\end{figC}

\begin{warnbox}Chỉ kiểm tra ``điểm có quá xa không'' là chưa đủ: $(30;40)$ không hề xa hơn $(42;44)$. Điều quyết định là tính chẵn -- lẻ và các tổng có thể tạo được.\end{warnbox}

\verify{Mỗi đẳng thức trong bảng đều tính lại được, ví dụ $-1-3+5+7+9=17$. Duyệt bằng máy tính toàn bộ $2\cdot 2^{10}=2048$ cách đi cũng xác nhận: chỉ $(30;40)$ là không đạt được.}
Chọn phương án \textbf{(B)}.
\ans{(B) $(30;40)$}
}

\begin{probbox}{25}{\relax}
Trong trò chơi ``Ba đồng xu'', người chơi tung cùng lúc ba đồng xu. Người chơi thắng nếu cả ba đồng xu đều ngửa hoặc cả ba đều sấp. Nếu chỉ chơi một lần thì xác suất thắng là bao nhiêu?
\choices{$\dfrac16$}{$\dfrac14$}{$\dfrac{2}{27}$}{$\dfrac23$}{$\dfrac13$}
\end{probbox}
\probsol{
\textbf{Phân tích đề.} Mỗi đồng xu có $2$ mặt: ngửa (N) hoặc sấp (S), khả năng như nhau. Thắng khi ba đồng ra NNN hoặc SSS. Ta liệt kê mọi kết quả có thể rồi đếm.

\begin{recallbox}Khi các kết quả có khả năng xảy ra như nhau: \emph{xác suất} của một biến cố $=$ (số kết quả thuận lợi) $:$ (tổng số kết quả có thể).\end{recallbox}

\step{1}{Liệt kê mọi kết quả bằng sơ đồ cây}
\begin{figC}\begin{scope}[level distance=1.9cm,grow'=right,
  level 1/.style={sibling distance=2.4cm},level 2/.style={sibling distance=1.2cm},level 3/.style={sibling distance=0.6cm},
  every node/.style={font=\footnotesize\bfseries}, edge from parent/.style={draw=tmBlue!70,thick}]
\node[text=Ink] {Tung}
  child {node[text=tmOrange] {N}
    child {node[text=tmOrange] {N} child {node[text=AnsRed] {N \ $\to$ NNN (thắng)}} child {node[text=Slate] {S \ $\to$ NNS}}}
    child {node[text=tmOrange] {S} child {node[text=Slate] {N \ $\to$ NSN}} child {node[text=Slate] {S \ $\to$ NSS}}}}
  child {node[text=tmPurple] {S}
    child {node[text=tmPurple] {N} child {node[text=Slate] {N \ $\to$ SNN}} child {node[text=Slate] {S \ $\to$ SNS}}}
    child {node[text=tmPurple] {S} child {node[text=Slate] {N \ $\to$ SSN}} child {node[text=AnsRed] {S \ $\to$ SSS (thắng)}}}};
\node[font=\scriptsize,text=Slate] at (1.9,2.55) {đồng $1$};
\node[font=\scriptsize,text=Slate] at (3.8,2.55) {đồng $2$};
\node[font=\scriptsize,text=Slate] at (5.7,2.55) {đồng $3$};
\end{scope}\end{figC}
Đồng thứ nhất có $2$ khả năng, với mỗi khả năng đó đồng thứ hai có $2$ khả năng, rồi đồng thứ ba có $2$ khả năng: tổng cộng $2\cdot2\cdot2=8$ kết quả, có khả năng như nhau.

\step{2}{Đếm kết quả thuận lợi}
Chỉ có $2$ kết quả thắng: NNN và SSS.

\step{3}{Tính xác suất}
\[P(\text{thắng})=\frac{2}{8}=\frac14.\]

\begin{warnbox}Một lập luận \emph{sai} hay gặp: ``Số mặt ngửa có thể là $0,1,2,3$ -- bốn trường hợp, thắng ở $2$ trường hợp ($0$ và $3$), nên xác suất là $\dfrac24$.'' Sai vì bốn trường hợp này \emph{không} có khả năng như nhau: ``$1$ mặt ngửa'' xảy ra theo $3$ cách (NSS, SNS, SSN), còn ``$3$ mặt ngửa'' chỉ có $1$ cách.\end{warnbox}

\begin{tipbox}Cách nghĩ nhanh: đồng thứ nhất ra mặt gì cũng được. Để thắng, đồng thứ hai phải giống đồng thứ nhất (xác suất $\dfrac12$) và đồng thứ ba cũng phải giống (xác suất $\dfrac12$): $\dfrac12\cdot\dfrac12=\dfrac14$.\end{tipbox}

\verify{$8$ kết quả trong sơ đồ cây khác nhau và đầy đủ; đúng $2$ kết quả thắng. $\dfrac28=\dfrac14$.}
Chọn phương án \textbf{(B)}.
\ans{(B) $\dfrac14$}
}

\begin{probbox}{26}{\relax}
Cho tam giác $PQR$ vuông tại $Q$ và $PQ=QR$. Các đoạn thẳng $QS$, $TU$ và $VW$ vuông góc với $PR$; các đoạn thẳng $ST$ và $UV$ vuông góc với $QR$ (như hình vẽ). Phần tô đậm chiếm bao nhiêu phần diện tích tam giác $PQR$?
\begin{figC}\begin{scope}[scale=0.42]
\coordinate (Q) at (0,0); \coordinate (R) at (8,0); \coordinate (P) at (0,8);
\coordinate (S) at (4,4); \coordinate (T) at (4,0); \coordinate (U) at (6,2); \coordinate (V) at (6,0); \coordinate (W) at (7,1);
\fill[tmBlue!35] (S)--(T)--(U)--cycle; \fill[tmBlue!35] (U)--(V)--(W)--cycle;
\draw[Navy,line width=1pt] (P)--(Q)--(R)--cycle;
\draw[Navy,line width=0.8pt] (Q)--(S)--(T)--(U)--(V)--(W);
\draw[Navy,line width=0.5pt] (0,0.5)--(0.5,0.5)--(0.5,0);
\node[above] at (P) {$P$}; \node[below left] at (Q) {$Q$}; \node[below right] at (R) {$R$};
\node[above right] at (S) {$S$}; \node[below] at (T) {$T$}; \node[above right] at (U) {$U$}; \node[below] at (V) {$V$}; \node[above right] at (W) {$W$};
\end{scope}\end{figC}
\choices{$\dfrac{3}{16}$}{$\dfrac38$}{$\dfrac{5}{16}$}{$\dfrac{5}{32}$}{$\dfrac{7}{32}$}
\end{probbox}
\probsol{
\textbf{Phân tích đề.} Tam giác $PQR$ là tam giác \emph{vuông cân} tại $Q$. Hình vẽ được tạo ra bằng cách cứ ``hạ đường vuông góc'' liên tiếp. Phần tô đậm gồm hai tam giác: $STU$ và $UVW$. Ý tưởng chính: \emph{mỗi lần hạ đường vuông góc từ đỉnh góc vuông của một tam giác vuông cân, ta chia nó thành hai nửa bằng nhau, mỗi nửa lại là tam giác vuông cân.}

\begin{recallbox}Trong tam giác cân, đường cao kẻ từ đỉnh (xuống cạnh đáy) cũng là đường trung tuyến, chia tam giác thành hai tam giác bằng nhau. Tam giác vuông cân có hai góc nhọn bằng $45^\circ$.\end{recallbox}

\begin{proofbox}Xét tam giác $XYZ$ vuông cân tại $X$ ($XY=XZ$), kẻ $XH\perp YZ$. Hai tam giác vuông $XHY$ và $XHZ$ có cạnh huyền $XY=XZ$ và cạnh góc vuông $XH$ chung, nên bằng nhau (cạnh huyền -- cạnh góc vuông). Vậy mỗi tam giác có diện tích bằng một nửa tam giác $XYZ$. Hơn nữa, $\widehat{Y}=45^\circ$ và $\widehat{XHY}=90^\circ$ nên $\widehat{HXY}=45^\circ$: tam giác $XHY$ vuông cân tại $H$. Tương tự với tam giác $XHZ$.\end{proofbox}

\step{1}{Chia đôi liên tiếp}
Gọi diện tích tam giác $PQR$ là $1$ (cả tam giác).
\begin{itemize}[leftmargin=1.4em,itemsep=1pt]
\item $QS\perp PR$ chia $PQR$ thành hai nửa: tam giác $QSR$ có diện tích $\dfrac12$, vuông cân tại $S$.
\item $ST\perp QR$ chia $QSR$ thành hai nửa: tam giác $STR$ có diện tích $\dfrac14$, vuông cân tại $T$.
\item $TU\perp SR$ chia $STR$ thành hai nửa: tam giác $STU$ \textbf{(tô đậm)} có diện tích $\dfrac18$; tam giác $TUR$ có diện tích $\dfrac18$, vuông cân tại $U$.
\item $UV\perp TR$ chia $TUR$ thành hai nửa: tam giác $UVR$ có diện tích $\dfrac1{16}$, vuông cân tại $V$.
\item $VW\perp UR$ chia $UVR$ thành hai nửa: tam giác $UVW$ \textbf{(tô đậm)} có diện tích $\dfrac1{32}$.
\end{itemize}

\step{2}{Cộng diện tích phần tô đậm}
\[\frac18+\frac1{32}=\frac{4}{32}+\frac{1}{32}=\frac{5}{32}.\]

\step{3}{Kiểm tra bằng lưới ô vuông}
Đặt $Q$ ở gốc, mỗi cạnh góc vuông dài $8$ ô: $Q(0;0)$, $R(8;0)$, $P(0;8)$. Khi đó $S(4;4)$, $T(4;0)$, $U(6;2)$, $V(6;0)$, $W(7;1)$ (như hình). Tam giác $PQR$ có diện tích $\dfrac{8\cdot 8}{2}=32$ ô. Tam giác $STU$ có cạnh $ST=4$ (thẳng đứng) và đỉnh $U$ cách $ST$ một đoạn $2$, nên diện tích $\dfrac{4\cdot 2}{2}=4$ ô. Tam giác $UVW$ có $UV=2$, $W$ cách $UV$ một đoạn $1$, diện tích $\dfrac{2\cdot1}{2}=1$ ô. Tô đậm: $4+1=5$ ô trên $32$ ô, tức $\dfrac5{32}$.

\begin{recallbox}Diện tích tam giác $=\dfrac{\text{cạnh đáy}\times\text{chiều cao}}{2}$.\end{recallbox}

\begin{warnbox}Nhiều bạn chọn $\dfrac{3}{16}$ hay $\dfrac38$ vì ``nhìn hình ước lượng''. Hãy luôn tính bằng chuỗi chia đôi: mỗi tam giác nhỏ hơn bằng một nửa tam giác chứa nó.\end{warnbox}

\verify{Hai cách (chia đôi và đếm ô) cùng cho $\dfrac{5}{32}$.}
Chọn phương án \textbf{(D)}.
\ans{(D) $\dfrac{5}{32}$}
}

\begin{probbox}{27}{\relax}
Chị Mai đến một phiên chợ đồ cũ. Ở đó, cứ sau mỗi $15$ phút, giá của mọi món hàng lại được giảm $10\%$ so với giá \emph{hiện tại} của nó. Lúc $9\,$h$\,00$, một tấm thảm có giá $100\,000$ đồng. Lúc $9\,$h$\,15$, giá tấm thảm giảm còn $90\,000$ đồng. Ngay khi giá tấm thảm xuống \emph{dưới} $80\,000$ đồng, chị Mai mua nó. Chị Mai mua tấm thảm lúc mấy giờ?
\choicesII{$9\,$h$\,45$}{$9\,$h$\,15$}{$9\,$h$\,30$}{$10\,$h$\,15$}{$10\,$h$\,00$}
\end{probbox}
\probsol{
\textbf{Phân tích đề.} Mỗi lần giảm giá là giảm $10\%$ của \emph{giá đang có} (không phải của giá ban đầu $100\,000$ đồng). Ta tính giá sau từng mốc $15$ phút cho đến khi giá nhỏ hơn $80\,000$ đồng.

\begin{recallbox}Giảm $10\%$ một số nghĩa là bớt đi $\dfrac{10}{100}=\dfrac1{10}$ của số đó; phần còn lại là $90\%$ số đó. Vậy: giá mới $=$ giá cũ $\times 0{,}9$ (hoặc giá cũ $-$ giá cũ $:10$).\end{recallbox}

\step{1}{Lập bảng giá theo thời gian}
\begin{center}
\begin{tabular}{|c|c|c|c|}
\hline
\rowcolor{Tint}\textbf{Thời điểm} & \textbf{Giá trước khi giảm} & \textbf{Số tiền giảm ($10\%$)} & \textbf{Giá mới}\\
\hline
$9\,$h$\,00$ & -- & -- & $100\,000$\\
$9\,$h$\,15$ & $100\,000$ & $10\,000$ & $90\,000$\\
$9\,$h$\,30$ & $90\,000$ & $9\,000$ & $81\,000$\\
$9\,$h$\,45$ & $81\,000$ & $8\,100$ & $\mathbf{72\,900}$\\
\hline
\end{tabular}
\end{center}

\step{2}{Tìm lần đầu tiên giá dưới $80\,000$ đồng}
Lúc $9\,$h$\,30$ giá là $81\,000$ đồng -- vẫn \emph{lớn hơn} $80\,000$ đồng, chị Mai chưa mua. Lúc $9\,$h$\,45$ giá là $72\,900$ đồng $<80\,000$ đồng: đây là lần đầu giá xuống dưới $80\,000$ đồng. Chị Mai mua lúc $9\,$h$\,45$.

\begin{warnbox}(1) Giảm $10\%$ \emph{của giá hiện tại}, nên số tiền giảm mỗi lần nhỏ dần ($10\,000$, rồi $9\,000$, rồi $8\,100$\dots), không phải lần nào cũng giảm $10\,000$ đồng. (2) Lúc $9\,$h$\,30$ giá $81\,000$ đồng ``rất gần'' $80\,000$ nhưng chưa \emph{dưới} $80\,000$, nên chưa phải lúc mua.\end{warnbox}

\verify{$100\,000\times 0{,}9\times 0{,}9\times 0{,}9=72\,900$ (sau $3$ lần giảm, tức $3\times 15=45$ phút kể từ $9\,$h$\,00$); còn sau $2$ lần giảm: $100\,000\times0{,}81=81\,000>80\,000$.}
Chọn phương án \textbf{(A)}.
\ans{(A) $9\,$h$\,45$}
}

\begin{probbox}{28}{\relax}
Người ta dùng số $2018$ để tạo ra các số nguyên dương có sáu chữ số: mỗi số đó phải chứa bốn chữ số $2,0,1,8$ \emph{đứng liền nhau và theo đúng thứ tự này}. Ví dụ, số $720\,186$ thỏa mãn, còn các số $209\,318$ và $210\,893$ thì không. Có bao nhiêu số có sáu chữ số như vậy chia hết cho $9$?
\choices{$28$}{$27$}{$31$}{$34$}{$22$}
\end{probbox}
\probsol{
\textbf{Phân tích đề.} Số có sáu chữ số chứa khối ``$2018$'' liền nhau, cộng thêm hai chữ số nữa (gọi là $a$ và $b$). Cần đếm các số chia hết cho $9$.

\begin{recallbox}\textbf{Dấu hiệu chia hết cho $9$:} một số chia hết cho $9$ khi và chỉ khi tổng các chữ số của nó chia hết cho $9$. Chữ số đầu tiên của một số có sáu chữ số phải khác $0$.\end{recallbox}

\step{1}{Điều kiện về tổng chữ số}
Tổng bốn chữ số của khối là $2+0+1+8=11$. Cả số có tổng chữ số $11+a+b$. Vì $a,b$ là chữ số nên $0\le a+b\le 18$, do đó $11\le 11+a+b\le 29$. Các bội của $9$ trong khoảng này là $18$ và $27$. Vậy
\[a+b=7\qquad\text{hoặc}\qquad a+b=16.\]
Các cặp có tổng $16$ chỉ có: $(7;9),\ (8;8),\ (9;7)$.

\step{2}{Ba vị trí có thể đặt khối $2018$}
\begin{center}
\begin{tabular}{|c|c|c|c|c|}
\hline
\rowcolor{Tint}\textbf{Dạng số} & \textbf{Điều kiện thêm} & \textbf{$a+b=7$} & \textbf{$a+b=16$} & \textbf{Cộng}\\
\hline
$\nd{2018ab}$ & không có ($a$ được bằng $0$) & $8$ cặp: $a=0,1,\dots,7$ & $3$ & $11$\\
$\nd{a2018b}$ & $a\neq 0$ & $7$ cặp: $a=1,\dots,7$ & $3$ & $10$\\
$\nd{ab2018}$ & $a\neq 0$ & $7$ cặp: $a=1,\dots,7$ & $3$ & $10$\\
\hline
\multicolumn{4}{|r|}{\textbf{Tổng}} & $\mathbf{31}$\\
\hline
\end{tabular}
\end{center}
Giải thích cột ``$a+b=7$'': các cặp là $(0;7),(1;6),(2;5),(3;4),(4;3),(5;2),(6;1),(7;0)$ -- có $8$ cặp. Khi $a$ đứng đầu số thì bỏ cặp $(0;7)$, còn $7$ cặp. Các cặp có tổng $16$ đều có $a\ge 7$ nên không bị loại.

\step{3}{Có số nào bị đếm hai lần không?}
Một số bị đếm hai lần chỉ khi khối $2018$ xuất hiện ở hai vị trí khác nhau trong cùng một số có sáu chữ số. Hai khối $4$ chữ số trong $6$ chữ số buộc phải ``chồng lên nhau'', mà khối $2018$ không thể chồng lên chính nó (ví dụ dịch $1$ vị trí thì chữ số $0$ phải trùng chữ số $2$\dots). Vậy không có số nào bị đếm hai lần.

\begin{warnbox}Hay quên một trong hai trường hợp: nhiều bạn chỉ xét $a+b=7$ (được $8+7+7=22$, đúng là phương án nhiễu (E)) mà quên $a+b=16$. Cũng hay quên điều kiện chữ số đầu khác $0$ (khi đó sẽ đếm thành $34$).\end{warnbox}

\verify{Vài ví dụ: $201\,870$ có tổng chữ số $18$; $720\,186$ có tổng $24$ -- không chia hết cho $9$ (đúng, vì $7+6=13$ không bằng $7$ hay $16$); $920\,187$ có tổng $27$, chia hết cho $9$. Duyệt bằng máy tính mọi số từ $100\,000$ đến $999\,999$ cũng được đúng $31$ số.}
Chọn phương án \textbf{(C)}.
\ans{(C) $31$}
}

\begin{probbox}{29}{\relax}
Gọi $N$ là tích của tất cả các số nguyên dương lẻ từ $1$ đến $99$ mà \emph{không} có chữ số tận cùng là $5$, tức là
\[N=1\times 3\times 7\times 9\times 11\times 13\times 17\times 19\times\cdots\times 91\times 93\times 97\times 99.\]
Chữ số hàng đơn vị của $N$ là
\choices{$1$}{$3$}{$5$}{$7$}{$9$}
\end{probbox}
\probsol{
\textbf{Phân tích đề.} $N$ là một số khổng lồ, không thể tính trực tiếp. Nhưng ta chỉ cần \emph{chữ số tận cùng}, và chữ số tận cùng của một tích chỉ phụ thuộc vào chữ số tận cùng của các thừa số.

\begin{recallbox}Chữ số tận cùng của một tích bằng chữ số tận cùng của tích các chữ số tận cùng của các thừa số. Ví dụ: $13\times 17=221$ có tận cùng $1$, cũng như $3\times 7=21$.\end{recallbox}

\step{1}{Xem các thừa số có chữ số tận cùng nào}
Các số lẻ không tận cùng bằng $5$ chỉ tận cùng bằng $1,3,7,9$. Chia các thừa số thành $10$ nhóm theo từng chục:
\[(1\cdot3\cdot7\cdot9),\ (11\cdot13\cdot17\cdot19),\ (21\cdot23\cdot27\cdot29),\ \dots,\ (91\cdot93\cdot97\cdot99).\]
Mỗi nhóm có $4$ thừa số với chữ số tận cùng lần lượt $1,3,7,9$; có $10$ nhóm (từ chục $0$ đến chục $9$), tổng cộng $40$ thừa số.

\step{2}{Chữ số tận cùng của mỗi nhóm}
$1\cdot3\cdot7\cdot9=189$, tận cùng là $9$. Vậy tích của mỗi nhóm tận cùng bằng $9$.

\step{3}{Nhân $10$ nhóm lại}
Chữ số tận cùng của $N$ bằng chữ số tận cùng của $\underbrace{9\cdot 9\cdots 9}_{10\text{ thừa số}}=9^{10}$. Ta có $9\cdot 9=81$ tận cùng $1$. Ghép $10$ thừa số $9$ thành $5$ cặp, mỗi cặp tận cùng $1$, nên $9^{10}$ tận cùng bằng $1\cdot1\cdot1\cdot1\cdot1=1$.

\begin{tipbox}Lũy thừa của $9$ có chữ số tận cùng luân phiên: $9^1\to 9$, $9^2\to 1$, $9^3\to 9$, $9^4\to 1$,\dots: số mũ lẻ tận cùng $9$, số mũ chẵn tận cùng $1$.\end{tipbox}

\begin{warnbox}Đừng nghĩ ``tích của toàn số lẻ thì tận cùng là $9$'' hay chọn $5$ theo phản xạ: các số tận cùng bằng $5$ đã bị loại khỏi tích. Cũng cần đếm đúng số nhóm: từ $1$ đến $99$ có $10$ nhóm (nhóm đầu là $1,3,7,9$).\end{warnbox}

\verify{Tính dần: sau nhóm $1$ tận cùng $9$, sau nhóm $2$ tận cùng $9\cdot9\to1$, sau nhóm $3$ tận cùng $9$,\dots, sau nhóm $10$ (số chẵn) tận cùng $1$. Máy tính tính đúng tích $N$ cũng cho chữ số tận cùng $1$.}
Chọn phương án \textbf{(A)}.
\ans{(A) $1$}
}

\begin{probbox}{30}{\relax}
Năm bạn học sinh làm một bài kiểm tra có điểm tối đa là $50$. Điểm của bốn bạn là $42$, $43$, $46$ và $49$. Điểm của bạn thứ năm là $N$. Trung bình cộng điểm của năm bạn bằng trung vị điểm của năm bạn. Số giá trị có thể có của $N$ là
\choices{$3$}{$4$}{$1$}{$0$}{$2$}
\end{probbox}
\probsol{
\textbf{Phân tích đề.} $N$ là một điểm từ $0$ đến $50$. Trung bình cộng phụ thuộc vào $N$; trung vị cũng phụ thuộc vào vị trí của $N$ khi xếp năm điểm theo thứ tự. Ta xét từng trường hợp vị trí của $N$.

\begin{recallbox}\textbf{Trung vị} của năm số là số đứng \emph{chính giữa} (số thứ ba) khi xếp chúng từ nhỏ đến lớn. \textbf{Trung bình cộng} $=$ tổng $:$ số các số.\end{recallbox}

\step{1}{Biểu diễn trung bình cộng theo $N$}
Tổng bốn điểm đã biết: $42+43+46+49=180$. Trung bình cộng của năm bạn là $\dfrac{180+N}{5}$.

\step{2}{Trung vị là bao nhiêu? -- Ba trường hợp}
Xếp $42,43,46,49$ cùng với $N$ theo thứ tự tăng dần:
\begin{itemize}[leftmargin=1.4em,itemsep=1pt]
\item Nếu $N\le 43$: thứ tự là $\dots,43,46,49$ với $N$ ở hai vị trí đầu, số thứ ba là $43$. Trung vị $=43$.
\item Nếu $43\le N\le 46$: thứ tự $42,43,N,46,49$. Trung vị $=N$.
\item Nếu $N\ge 46$: thứ tự $42,43,46,\dots$ với $N$ ở hai vị trí cuối. Trung vị $=46$.
\end{itemize}
\begin{figC}
\draw[->,thick,Navy] (0,0)--(10.6,0);
\foreach \x/\lab in {0/0,6/42,6.6/43,8.4/46,10.2/50}{\draw[Navy] (\x,0.1)--(\x,-0.1); \node[below,font=\footnotesize] at (\x,-0.1) {$\lab$};}
\draw[tmBlue,line width=3pt,opacity=0.5] (0,0.35)--(6.6,0.35); \node[above,font=\scriptsize,text=tmBlue] at (3.3,0.4) {trung vị $=43$};
\draw[tmOrange,line width=3pt,opacity=0.6] (6.6,0.35)--(8.4,0.35); \node[above,font=\scriptsize,text=tmOrange!80!black] at (7.5,0.4) {trung vị $=N$};
\draw[tmGreen,line width=3pt,opacity=0.5] (8.4,0.35)--(10.2,0.35); \node[above,font=\scriptsize,text=tmGreen!70!black] at (9.4,0.85) {trung vị $=46$};
\end{figC}

\step{3}{Giải từng trường hợp}
\begin{itemize}[leftmargin=1.4em,itemsep=1pt]
\item \textbf{Trường hợp 1} (trung vị $43$): $\dfrac{180+N}{5}=43\Rightarrow 180+N=215\Rightarrow N=35$. Kiểm tra: $35\le 43$ -- đúng trường hợp đang xét. \textbf{Nhận.}
\item \textbf{Trường hợp 2} (trung vị $N$): $\dfrac{180+N}{5}=N\Rightarrow 180+N=5N\Rightarrow 4N=180\Rightarrow N=45$. Kiểm tra: $43\le 45\le 46$. \textbf{Nhận.}
\item \textbf{Trường hợp 3} (trung vị $46$): $\dfrac{180+N}{5}=46\Rightarrow 180+N=230\Rightarrow N=50$. Kiểm tra: $50\ge 46$ và $50$ không vượt điểm tối đa. \textbf{Nhận.}
\end{itemize}

\begin{warnbox}Sau khi giải ra $N$ trong mỗi trường hợp, \emph{phải kiểm tra} $N$ có thực sự nằm trong trường hợp đó không (và có nằm trong khoảng $0$ đến $50$ không). Ở bài này cả ba giá trị đều hợp lệ, nhưng ở nhiều bài khác sẽ có giá trị bị loại. Đừng dừng lại sau khi tìm được một giá trị!\end{warnbox}

\verify{
$N=35$: điểm $35,42,43,46,49$, tổng $215$, trung bình $43$, trung vị $43$.
$N=45$: điểm $42,43,45,46,49$, tổng $225$, trung bình $45$, trung vị $45$.
$N=50$: điểm $42,43,46,49,50$, tổng $230$, trung bình $46$, trung vị $46$.}
Có $3$ giá trị của $N$, ứng với phương án \textbf{(A)}.
\ans{(A) $3$}
}

% ---------------- part c ----------------
% =====================================================================
% Chủ đề 4 — Số học và Tư duy số — phần c: bài 31–40 (mở rộng từ bài 1–15)
% =====================================================================
\clearpage
\dividerbar{Mở rộng -- Nâng cao}

% ---------------------------------------------------------------- Bài 31
\begin{probbox}{31}{\relax}
Cho các số: $-1{,}5;\quad 0{,}25;\quad -0{,}75;\quad 1{,}2;\quad -2$.
\begin{enumerate}[label=\alph*),leftmargin=1.6em,itemsep=0pt,topsep=2pt]
  \item Sắp xếp các số theo thứ tự tăng dần.
  \item Những số nào nằm giữa $-1$ và $1$ trên trục số?
  \item Tính tổng của số lớn nhất và số nhỏ nhất.
\end{enumerate}
\end{probbox}
\probsol{
\textbf{Phân tích đề.} Bài này kết hợp hai kỹ năng của Bài 1 và Bài 3: so sánh số thập phân và so sánh số âm.
\begin{recallbox}Số âm luôn nhỏ hơn số dương. Trong hai số âm, số nào xa $0$ hơn thì nhỏ hơn. Trong hai số dương,
so sánh phần nguyên rồi từng hàng thập phân.\end{recallbox}
\step{1}{Chia thành hai nhóm: số âm và số dương}
Số âm: $-1{,}5;\ -0{,}75;\ -2$. \quad Số dương: $0{,}25;\ 1{,}2$.
\step{2}{Sắp xếp nhóm số âm}
Khoảng cách đến $0$: $-2$ cách $2$; \ $-1{,}5$ cách $1{,}5$; \ $-0{,}75$ cách $0{,}75$. Số xa $0$ nhất là nhỏ nhất:
\[ -2 < -1{,}5 < -0{,}75. \]
\step{3}{Sắp xếp nhóm số dương rồi ghép lại}
$0{,}25 < 1{,}2$ (phần nguyên $0 < 1$). Ghép: số âm đứng trước, số dương đứng sau:
\[ -2 < -1{,}5 < -0{,}75 < 0{,}25 < 1{,}2. \]
\begin{figC}\begin{scope}[x=1.6cm]
  \draw[->,thick,Navy] (-2.5,0)--(1.7,0);
  \foreach \x in {-2,-1,0,1}{\draw[Navy] (\x,0.08)--(\x,-0.08); \node[below,font=\scriptsize] at (\x,-0.1) {$\x$};}
  \foreach \x/\l in {-2/{-2},-1.5/{-1{,}5},-0.75/{-0{,}75},0.25/{0{,}25},1.2/{1{,}2}}{\fill[tmOrange] (\x,0) circle (2.6pt); \node[above,font=\scriptsize] at (\x,0.1) {$\l$};}
  \draw[tmGreen!70!black,line width=2pt,opacity=0.35] (-1,0)--(1,0);
\end{scope}\end{figC}
\step{4}{Các số nằm giữa $-1$ và $1$}
Trên trục số, đoạn tô xanh từ $-1$ đến $1$ chứa $-0{,}75$ và $0{,}25$. (Số $-1{,}5$ ở bên trái $-1$; số $1{,}2$ ở bên phải $1$.)
\step{5}{Tổng của số lớn nhất và số nhỏ nhất}
Số lớn nhất là $1{,}2$, số nhỏ nhất là $-2$:
\[ 1{,}2 + (-2) = -(2 - 1{,}2) = -0{,}8. \]
\begin{recallbox}Cộng hai số trái dấu: lấy số lớn trừ số bé (không kể dấu), rồi đặt dấu của số có khoảng cách đến $0$ lớn hơn.\end{recallbox}
\begin{warnbox}Nhiều bạn xếp $-0{,}75$ trước $-1{,}5$ vì ``$0{,}75 < 1{,}5$''. Với số âm thì ngược lại: $-1{,}5 < -0{,}75$.\end{warnbox}
\verify{$-0{,}8 - 1{,}2 = -2$. Đúng.}
\ans{a) $-2 < -1{,}5 < -0{,}75 < 0{,}25 < 1{,}2$;\quad b) $-0{,}75$ và $0{,}25$;\quad c) $-0{,}8$}
}

% ---------------------------------------------------------------- Bài 32
\begin{probbox}{32}{\relax}
Cho biểu thức $2 + 3 \times 4 - 1$.
\begin{enumerate}[label=\alph*),leftmargin=1.6em,itemsep=0pt,topsep=2pt]
  \item Tính giá trị của biểu thức.
  \item Đặt thêm \emph{một} cặp dấu ngoặc $(\ )$ vào biểu thức (không đổi thứ tự các số và các dấu) để được giá trị
  lớn nhất có thể. Giá trị lớn nhất đó là bao nhiêu?
  \item Cũng hỏi như câu b) nhưng với giá trị nhỏ nhất.
\end{enumerate}
\end{probbox}
\probsol{
\textbf{Phân tích đề.} Bài này mở rộng Bài 4 và Bài 7: vị trí dấu ngoặc làm thay đổi thứ tự tính, nên làm thay đổi kết quả.
\begin{recallbox}Trong ngoặc trước; rồi nhân, chia; cuối cùng cộng, trừ (từ trái sang phải).\end{recallbox}
\step{1}{Giá trị khi chưa có ngoặc}
$2 + 3\times 4 - 1 = 2 + 12 - 1 = 13$.
\step{2}{Liệt kê có hệ thống mọi cách đặt một cặp ngoặc}
Một cặp ngoặc phải bao ít nhất hai số liền nhau. Gọi bốn số theo thứ tự là số thứ $1, 2, 3, 4$ ($2;\ 3;\ 4;\ 1$).
Ngoặc có thể bắt đầu ở số thứ $i$ và kết thúc ở số thứ $j$ với $i < j$: có $6$ cách.
\begin{center}
\begin{tabular}{|c|l|c|}
\hline
\rowcolor{tmBlue!10} Ngoặc bao & Biểu thức & Giá trị\\ \hline
số 1 -- 2 & $(2+3)\times 4 - 1 = 5\times 4 - 1$ & $\mathbf{19}$\\ \hline
số 1 -- 3 & $(2+3\times 4) - 1 = 14 - 1$ & $13$\\ \hline
số 1 -- 4 & $(2+3\times 4-1)$ & $13$\\ \hline
số 2 -- 3 & $2 + (3\times 4) - 1$ & $13$\\ \hline
số 2 -- 4 & $2 + (3\times 4 - 1) = 2 + 11$ & $13$\\ \hline
số 3 -- 4 & $2 + 3\times(4 - 1) = 2 + 9$ & $\mathbf{11}$\\ \hline
\end{tabular}
\end{center}
\step{3}{Đọc kết quả}
Lớn nhất là $19$, khi viết $(2+3)\times 4 - 1$: ta ``phóng to'' tổng $2 + 3$ lên $4$ lần.\\
Nhỏ nhất là $11$, khi viết $2 + 3\times(4-1)$: ta nhân $3$ với một số nhỏ hơn ($3$ thay vì $4$).
\begin{tipbox}Muốn kết quả lớn, hãy để phép nhân tác động lên một tổng lớn; muốn kết quả nhỏ, hãy để phép nhân tác động
lên một hiệu. Nhưng để chắc chắn, cách tốt nhất vẫn là liệt kê đủ mọi trường hợp như trong bảng.\end{tipbox}
\verify{$(2+3)\times 4 - 1 = 20 - 1 = 19$ và $2 + 3\times 3 = 11$.}
\ans{a) $13$;\quad b) $(2+3)\times 4 - 1 = 19$;\quad c) $2 + 3\times(4-1) = 11$}
}

% ---------------------------------------------------------------- Bài 33
\begin{probbox}{33}{\relax}
Cuốn truyện của bạn An dày $360$ trang. Ngày thứ nhất An đọc $\dfrac{1}{4}$ cuốn truyện. Ngày thứ hai An đọc
$\dfrac{1}{3}$ \emph{số trang còn lại}. Hỏi sau hai ngày, An còn bao nhiêu trang chưa đọc, và số trang đó bằng mấy phần
cuốn truyện?
\end{probbox}
\probsol{
\textbf{Phân tích đề.} Cần chú ý: phân số $\dfrac13$ của ngày thứ hai được tính trên \emph{số trang còn lại}, không phải
trên cả cuốn truyện. Đây là phiên bản ``hai bước'' của Bài 9.
\begin{recallbox}$\dfrac{m}{n}$ của $a$ bằng $a : n \times m$.\end{recallbox}
\step{1}{Ngày thứ nhất}
An đọc $360 : 4 = 90$ (trang). Còn lại $360 - 90 = 270$ (trang).
\step{2}{Ngày thứ hai (tính trên $270$ trang)}
An đọc $270 : 3 = 90$ (trang). Còn lại $270 - 90 = 180$ (trang).
\step{3}{Phần cuốn truyện còn lại}
$\dfrac{180}{360} = \dfrac12$.
\begin{barfigC}
  \blk{0}{2}{0}{0.6}{tmBlue}{$90$}
  \blk{2}{4}{0}{0.6}{tmGreen}{$90$}
  \blk{4}{6}{0}{0.6}{tmOrange}{$90$}
  \blk{6}{8}{0}{0.6}{tmOrange}{$90$}
  \brcup{0}{2}{0.5}{ngày 1}
  \brcup{2}{4}{0.5}{ngày 2}
  \brcup{4}{8}{0.5}{chưa đọc}
  \brcdn{2}{8}{-0.5}{còn lại sau ngày 1: $270$ trang, chia $3$ phần}
\end{barfigC}
\begin{warnbox}Sai lầm phổ biến: cộng $\dfrac14 + \dfrac13 = \dfrac{7}{12}$ rồi tính phần còn lại là $\dfrac{5}{12}$ cuốn
($150$ trang). Sai, vì $\dfrac13$ là của \emph{phần còn lại} ($270$ trang), không phải của cả cuốn ($360$ trang).\end{warnbox}
\step{4}{Cách khác: nhân các phân số ``còn lại''}
Sau ngày 1 còn $\dfrac34$ cuốn; sau ngày 2 còn $\dfrac23$ của phần đó, tức $\dfrac23\times\dfrac34 = \dfrac{6}{12} = \dfrac12$ cuốn,
bằng $360 : 2 = 180$ trang.
\verify{$90 + 90 + 180 = 360$. Đúng.}
\ans{$180$ trang, bằng $\dfrac12$ cuốn truyện}
}

% ---------------------------------------------------------------- Bài 34
\begin{probbox}{34}{\relax}
Bài kiểm tra thứ nhất có $40$ câu hỏi, bạn Bình làm đúng $32$ câu.
\begin{enumerate}[label=\alph*),leftmargin=1.6em,itemsep=0pt,topsep=2pt]
  \item Bình làm đúng bao nhiêu phần trăm số câu của bài thứ nhất?
  \item Bài kiểm tra thứ hai có $60$ câu hỏi. Bình phải làm đúng ít nhất bao nhiêu câu ở bài thứ hai để tính chung cả
  hai bài, số câu đúng đạt ít nhất $85\%$ tổng số câu?
\end{enumerate}
\end{probbox}
\probsol{
\textbf{Phân tích đề.} Câu a) là tìm tỉ số phần trăm (ngược với Bài 5). Câu b) là tìm phần trăm của một tổng.
\begin{recallbox}Tỉ số phần trăm của $a$ so với $b$ là $\dfrac{a}{b}\times 100\%$. \ $p\%$ của $a$ là $a\times\dfrac{p}{100}$.\end{recallbox}
\step{1}{Tỉ số phần trăm của bài thứ nhất}
$\dfrac{32}{40}\times 100\% = 0{,}8\times 100\% = 80\%$.
\step{2}{Tổng số câu của cả hai bài}
$40 + 60 = 100$ (câu).
\step{3}{Số câu đúng cần có tính chung}
$85\%$ của $100$ câu là $100\times\dfrac{85}{100} = 85$ (câu).
\step{4}{Số câu đúng cần ở bài thứ hai}
Bình đã có $32$ câu đúng, nên cần thêm ít nhất $85 - 32 = 53$ (câu) ở bài thứ hai.
\begin{warnbox}Không được lấy ``trung bình của hai tỉ lệ phần trăm'': nếu tìm $x$ sao cho $(80\% + x) : 2 = 85\%$ thì được
$x = 90\%$ của $60$ câu là $54$ câu -- sai! Hai bài có số câu khác nhau, nên phải tính trên \emph{tổng số câu}.\end{warnbox}
\verify{Với $53$ câu đúng: $(32 + 53) : 100 = 85\%$ (vừa đủ). Với $52$ câu: $84 : 100 = 84\% < 85\%$ (không đủ).}
\ans{a) $80\%$;\quad b) ít nhất $53$ câu}
}

% ---------------------------------------------------------------- Bài 35
\begin{probbox}{35}{\relax}
Một cửa hàng bán hạt điều đóng gói: gói $3$ kg giá $255$ nghìn đồng, gói $5$ kg giá $410$ nghìn đồng (không bán lẻ).
\begin{enumerate}[label=\alph*),leftmargin=1.6em,itemsep=0pt,topsep=2pt]
  \item Loại gói nào rẻ hơn tính theo mỗi ki-lô-gam?
  \item Cô Hằng cần mua \emph{đúng} $15$ kg. Cô nên mua thế nào để tốn ít tiền nhất? Hết bao nhiêu tiền?
  \item Nếu cần mua \emph{đúng} $13$ kg thì hết bao nhiêu tiền?
\end{enumerate}
\end{probbox}
\probsol{
\textbf{Phân tích đề.} Mở rộng Bài 6: so sánh đơn giá, rồi tìm cách ghép các gói cho đúng khối lượng.
\begin{recallbox}Đơn giá $=$ thành tiền $:$ số lượng.\end{recallbox}
\step{1}{Đơn giá của từng loại gói}
Gói $3$ kg: $255 : 3 = 85$ (nghìn đồng/kg). \quad Gói $5$ kg: $410 : 5 = 82$ (nghìn đồng/kg).\\
Vậy gói $5$ kg rẻ hơn ($82 < 85$).
\step{2}{Các cách ghép đúng $15$ kg}
Gọi $a$ là số gói $3$ kg, $b$ là số gói $5$ kg: $3a + 5b = 15$. Thử $b = 0, 1, 2, 3$:
\begin{center}
\begin{tabular}{|c|c|c|c|}
\hline
\rowcolor{tmBlue!10} $b$ & $15 - 5b$ & $a$ & Thành tiền (nghìn đồng)\\ \hline
$0$ & $15$ & $5$ & $5\times 255 = 1275$\\ \hline
$1$ & $10$ & không chia hết cho $3$ & --\\ \hline
$2$ & $5$ & không chia hết cho $3$ & --\\ \hline
$3$ & $0$ & $0$ & $3\times 410 = \mathbf{1230}$\\ \hline
\end{tabular}
\end{center}
Mua $3$ gói $5$ kg là rẻ nhất: $1\,230$ nghìn đồng $= 1\,230\,000$ đồng.
\step{3}{Đúng $13$ kg}
$3a + 5b = 13$: \ $b = 0 \Rightarrow 13$ (không chia hết cho $3$); \ $b = 1 \Rightarrow 8$ (không chia hết cho $3$);
\ $b = 2 \Rightarrow 3$, $a = 1$. Chỉ có một cách: $1$ gói $3$ kg và $2$ gói $5$ kg.\\
Thành tiền: $255 + 2\times 410 = 255 + 820 = 1075$ (nghìn đồng) $= 1\,075\,000$ đồng.
\begin{warnbox}Không thể trả lời câu c) bằng cách lấy $13\times 82$, vì cửa hàng không bán lẻ: phải ghép các gói sao cho
đúng $13$ kg.\end{warnbox}
\verify{$3 + 5 + 5 = 13$ kg; $5 + 5 + 5 = 15$ kg.}
\ans{a) gói $5$ kg ($82$ nghìn đồng/kg);\quad b) $3$ gói $5$ kg, $1\,230\,000$ đồng;\quad c) $1\,075\,000$ đồng}
}

% ---------------------------------------------------------------- Bài 36
\begin{probbox}{36}{\relax}
Năm bạn trong một nhóm có cân nặng (kg) là $38;\ 41;\ 40;\ 37$ và $x$. Cân nặng trung bình của nhóm là $40$ kg.
\begin{enumerate}[label=\alph*),leftmargin=1.6em,itemsep=0pt,topsep=2pt]
  \item Tìm $x$. Bạn thứ năm nặng hơn hay nhẹ hơn mức trung bình bao nhiêu ki-lô-gam?
  \item Tìm trung vị cân nặng của nhóm.
  \item Một bạn thứ sáu gia nhập nhóm, và cân nặng trung bình của cả sáu bạn là $41$ kg. Bạn thứ sáu nặng bao nhiêu?
\end{enumerate}
\end{probbox}
\probsol{
\textbf{Phân tích đề.} Mở rộng Bài 14: dùng tổng và độ lệch để tìm số còn thiếu; thêm trung vị và bài toán ``thêm một người''.
\begin{recallbox}Tổng $=$ trung bình cộng $\times$ số các số. Tổng các độ lệch so với trung bình cộng bằng $0$.\end{recallbox}
\step{1}{Tìm $x$ bằng độ lệch}
So với $40$: $38 \to -2$; \ $41 \to +1$; \ $40 \to 0$; \ $37 \to -3$. Tổng: $-2 + 1 + 0 - 3 = -4$.\\
Để tổng các độ lệch bằng $0$, bạn thứ năm phải lệch $+4$: $x = 40 + 4 = 44$ (kg), nặng hơn mức trung bình $4$ kg.\\
Kiểm tra bằng tổng: $5\times 40 = 200$; \ $200 - (38 + 41 + 40 + 37) = 200 - 156 = 44$.
\step{2}{Trung vị}
Sắp xếp tăng dần: $37;\ 38;\ \mathbf{40};\ 41;\ 44$. Số đứng giữa (thứ ba) là $40$. Trung vị là $40$ kg.
\step{3}{Bạn thứ sáu}
Tổng cân nặng sáu bạn: $6\times 41 = 246$ (kg). Năm bạn đầu nặng tổng cộng $200$ kg. Bạn thứ sáu nặng $246 - 200 = 46$ (kg).
\begin{tipbox}Có thể nghĩ bằng độ lệch: trung bình tăng $1$ kg cho cả $6$ người, tức cần thêm $6$ kg so với mức $40$;
bạn mới phải nặng $40 + 6 = 46$ kg.\end{tipbox}
\begin{warnbox}Không phải ``bạn thứ sáu nặng $41$ kg'' -- nếu thế, trung bình chỉ tăng một chút: $241 : 6 \approx 40{,}17$ kg.\end{warnbox}
\verify{$(37 + 38 + 40 + 41 + 44 + 46) : 6 = 246 : 6 = 41$. Đúng.}
\ans{a) $x = 44$ kg, nặng hơn trung bình $4$ kg;\quad b) $40$ kg;\quad c) $46$ kg}
}

% ---------------------------------------------------------------- Bài 37
\begin{probbox}{37}{\stars{1}}
Một đường chạy vòng quanh sân trường dài $200$ m. Mai chạy $6$ vòng hết $4$ phút. Tú chạy đều với vận tốc $18$ km/h.
\begin{enumerate}[label=\alph*),leftmargin=1.6em,itemsep=0pt,topsep=2pt]
  \item Ai chạy nhanh hơn?
  \item Nếu giữ nguyên vận tốc, mỗi bạn chạy $3$ km thì hết bao nhiêu phút?
  \item Hai bạn xuất phát cùng lúc, từ cùng một điểm trên đường chạy, nhưng chạy \emph{ngược chiều} nhau. Sau bao nhiêu
  giây thì hai bạn gặp nhau lần đầu tiên?
  \item Một hôm Tú tăng vận tốc lên $21{,}6$ km/h, còn Mai vẫn chạy như cũ. Hai bạn xuất phát cùng lúc, cùng một điểm, chạy
  \emph{cùng chiều}. Sau bao lâu thì Tú vượt Mai đúng một vòng?
\end{enumerate}
\end{probbox}
\probsol{
\textbf{Phân tích đề.} Mở rộng Bài 12: tính vận tốc từ số vòng chạy, đổi km/h sang m/s, rồi xét hai người chuyển động
trên cùng một đường vòng.
\begin{recallbox}$v = s : t$, \ $t = s : v$. \ $1$ km/h $= \dfrac{5}{18}$ m/s, tức muốn đổi km/h sang m/s ta nhân với $5$ rồi chia cho $18$.\end{recallbox}
\step{1}{Vận tốc của Mai}
Quãng đường: $6\times 200 = 1200$ (m). Thời gian: $4$ phút $= 240$ giây. \ $v_{\text{Mai}} = 1200 : 240 = 5$ (m/s).
\step{2}{Vận tốc của Tú (đổi ra m/s)}
$v_{\text{Tú}} = 18\times\dfrac{5}{18} = 5$ (m/s). Hai bạn chạy \textbf{nhanh như nhau}!
\begin{warnbox}So sánh trực tiếp ``$5$'' với ``$18$'' là vô nghĩa vì khác đơn vị (m/s và km/h). Phải đổi về cùng đơn vị
trước.\end{warnbox}
\step{3}{Thời gian chạy $3$ km}
$3$ km $= 3000$ m. \ $t = 3000 : 5 = 600$ (giây) $= 600 : 60 = 10$ (phút). Cả hai bạn đều mất $10$ phút.
\step{4}{Chạy ngược chiều}
Mỗi giây Mai chạy $5$ m, Tú chạy $5$ m theo hướng ngược lại, nên quãng đường ``chung'' hai bạn chạy được tăng
$5 + 5 = 10$ m mỗi giây. Hai bạn gặp nhau lần đầu khi \emph{tổng} quãng đường của hai bạn bằng đúng một vòng $200$ m:
\[ t = 200 : 10 = 20 \text{ (giây)}. \]
\begin{figC}
  \draw[Slate!60,line width=4pt] (0,0) circle (1.1);
  \fill[Navy] (0,1.1) circle (2.5pt); \node[above=4pt,font=\scriptsize] at (0,1.35) {xuất phát};
  \draw[->,tmBlue,thick] (95:1.4) arc (95:265:1.4);
  \draw[->,tmOrange,thick] (85:1.4) arc (85:-85:1.4);
  \node[left,font=\scriptsize,text=tmBlue!80!black] at (-1.5,0) {Mai: $100$ m};
  \node[right,font=\scriptsize,text=tmOrange!80!black] at (1.5,0) {Tú: $100$ m};
  \fill[AnsRed] (0,-1.1) circle (2.5pt); \node[below,font=\scriptsize,text=AnsRed] at (0,-1.5) {gặp nhau sau $20$ s};
\end{figC}
Sau $20$ giây, mỗi bạn chạy $5\times 20 = 100$ m, đúng nửa vòng: họ gặp nhau ở điểm đối diện với điểm xuất phát.
\step{5}{Chạy cùng chiều, Tú nhanh hơn}
$v_{\text{Tú}} = 21{,}6\times\dfrac{5}{18} = \dfrac{108}{18} = 6$ (m/s). Mỗi giây, Tú chạy hơn Mai $6 - 5 = 1$ m. Tú vượt Mai
đúng một vòng khi Tú chạy hơn Mai $200$ m:
\[ t = 200 : 1 = 200 \text{ (giây)} = 3 \text{ phút } 20 \text{ giây}. \]
\verify{Sau $200$ giây: Tú chạy $6\times 200 = 1200$ m ($6$ vòng), Mai chạy $5\times 200 = 1000$ m ($5$ vòng). Tú hơn đúng $1$ vòng.}
\begin{parentbox}Có thể diễn lại câu c) và d) bằng hai đồng xu trên một chiếc đĩa tròn: ngược chiều thì ``khoảng cách cần
xóa'' giảm theo \emph{tổng} hai vận tốc; cùng chiều thì giảm theo \emph{hiệu} hai vận tốc.\end{parentbox}
\ans{a) nhanh như nhau ($5$ m/s);\quad b) $10$ phút;\quad c) $20$ giây;\quad d) $200$ giây ($3$ phút $20$ giây)}
}

% ---------------------------------------------------------------- Bài 38
\begin{probbox}{38}{\stars{1}}
\begin{enumerate}[label=\alph*),leftmargin=1.6em,itemsep=0pt,topsep=0pt]
  \item Tìm các chữ số $a$, $b$ để số $\nd{4a7b}$ chia hết cho cả $2$, $5$ và $9$.
  \item Có bao nhiêu số có bốn chữ số dạng $\nd{5a3b}$ chia hết cho $9$? Hãy liệt kê chúng.
\end{enumerate}
\end{probbox}
\probsol{
\textbf{Phân tích đề.} Mở rộng Bài 10: dùng đồng thời nhiều dấu hiệu chia hết và đếm có hệ thống.
\begin{recallbox}Chia hết cho $2$ và $5$ $\iff$ chữ số tận cùng là $0$. \ Chia hết cho $9$ $\iff$ tổng các chữ số chia hết cho $9$.\end{recallbox}
\begin{proofbox}Vì sao ``chia hết cho cả $2$ và $5$'' lại tương đương với ``tận cùng là $0$''? Chia hết cho $2$ cần chữ số
tận cùng chẵn; chia hết cho $5$ cần chữ số tận cùng là $0$ hoặc $5$. Chữ số vừa chẵn vừa thuộc $\{0;5\}$ chỉ có $0$.\end{proofbox}
\step{1}{Câu a) -- tìm $b$}
$\nd{4a7b}$ chia hết cho $2$ và $5$ nên $b = 0$.
\step{2}{Câu a) -- tìm $a$}
Tổng các chữ số: $4 + a + 7 + 0 = 11 + a$. Vì $0 \le a \le 9$ nên $11 \le 11 + a \le 20$. Số chia hết cho $9$ trong khoảng
này chỉ có $18$. Vậy $11 + a = 18$, $a = 7$. Số cần tìm: $4770$.
\verify{$4770 = 9\times 530$, tận cùng là $0$ nên chia hết cho $2$ và $5$.}
\step{3}{Câu b) -- điều kiện}
Tổng các chữ số của $\nd{5a3b}$ là $5 + a + 3 + b = 8 + a + b$. Vì $0 \le a + b \le 18$ nên $8 \le 8 + a + b \le 26$. Các bội
của $9$ trong khoảng này là $9$ và $18$. Do đó $a + b = 1$ hoặc $a + b = 10$.
\step{4}{Câu b) -- đếm có hệ thống}
$a + b = 1$: $(a;b) = (0;1)$ hoặc $(1;0)$ \ -- \ $2$ số.\\
$a + b = 10$: $a$ chạy từ $1$ đến $9$, $b = 10 - a$ (từ $9$ xuống $1$) \ -- \ $9$ số. ($a = 0$ thì $b = 10$, không phải chữ số.)\\
Tổng cộng: $2 + 9 = 11$ số:
\[ 5031,\ 5130,\ 5139,\ 5238,\ 5337,\ 5436,\ 5535,\ 5634,\ 5733,\ 5832,\ 5931. \]
\begin{warnbox}Đừng quên $a$ được phép bằng $0$ (vì chữ số đầu tiên là $5$), và đừng bỏ sót trường hợp tổng chữ số bằng $9$
(tức $a + b = 1$).\end{warnbox}
\verify{Mỗi số trên có tổng chữ số bằng $9$ hoặc $18$, ví dụ $5\,436$: $5+4+3+6 = 18$.}
\ans{a) $a = 7$, $b = 0$ (số $4770$);\quad b) $11$ số}
}

% ---------------------------------------------------------------- Bài 39
\begin{probbox}{39}{\stars{1}}
Trong phép cộng dưới đây, $M$ và $N$ là hai chữ số khác nhau và khác $0$. Tìm tất cả các cặp $(M;N)$ thỏa mãn.
\[ \begin{array}{r@{\;}c@{\;}c} & M & N\\ & N & M\\ + & N & N\\ \hline 1 & 8 & 7 \end{array} \]
\end{probbox}
\probsol{
\textbf{Phân tích đề.} Mở rộng Bài 13. Lần này thay vì thử từng hàng, ta dùng cấu tạo số để đưa về một đẳng thức đơn giản,
rồi tìm \emph{tất cả} các nghiệm.
\begin{recallbox}$\nd{MN} = 10M + N$, \ $\nd{NM} = 10N + M$, \ $\nd{NN} = 10N + N = 11N$.\end{recallbox}
\step{1}{Viết tổng bằng cấu tạo số}
\[ \nd{MN} + \nd{NM} + \nd{NN} = (10M + N) + (10N + M) + 11N = 11M + 22N = 11\times(M + 2N). \]
\step{2}{Lập đẳng thức}
$11\times(M + 2N) = 187$, nên $M + 2N = 187 : 11 = 17$.
\step{3}{Tìm mọi cặp chữ số}
$M = 17 - 2N$ phải là chữ số từ $1$ đến $9$, nên $17 - 2N \le 9$, tức $N \ge 4$; và $17 - 2N \ge 1$, tức $N \le 8$.
\begin{center}
\begin{tabular}{|c|c|c|c|c|c|}
\hline
\rowcolor{tmBlue!10} $N$ & $4$ & $5$ & $6$ & $7$ & $8$\\ \hline
$M = 17 - 2N$ & $9$ & $7$ & $5$ & $3$ & $1$\\ \hline
Phép cộng & $94+49+44$ & $75+57+55$ & $56+65+66$ & $37+73+77$ & $18+81+88$\\ \hline
\end{tabular}
\end{center}
Cả năm cặp đều có $M \ne N$, và mỗi tổng đều bằng $187$.
\begin{tipbox}Khi các chữ cái xuất hiện ``đối xứng'' như $\nd{MN}$ và $\nd{NM}$, hãy thử cấu tạo số: thường có một thừa số
chung ($11$ ở đây) giúp bài toán gọn hẳn lại.\end{tipbox}
\begin{warnbox}Đề hỏi \emph{tất cả} các cặp. Tìm được một cặp (chẳng hạn $M = 9$, $N = 4$) rồi dừng là chưa đủ.\end{warnbox}
\verify{$94 + 49 + 44 = 187$; \ $75 + 57 + 55 = 187$; \ $56 + 65 + 66 = 187$; \ $37 + 73 + 77 = 187$; \ $18 + 81 + 88 = 187$.}
\ans{$(M;N) \in \{(9;4),\ (7;5),\ (5;6),\ (3;7),\ (1;8)\}$}
}

% ---------------------------------------------------------------- Bài 40
\begin{probbox}{40}{\stars{2}}
Một dãy số được tạo theo quy luật Fibonacci: kể từ số thứ ba, mỗi số bằng tổng của hai số đứng ngay trước nó.
\begin{enumerate}[label=\alph*),leftmargin=1.6em,itemsep=0pt,topsep=2pt]
  \item Biết số thứ nhất là $2$ và số thứ năm là $19$. Tìm số thứ hai và số thứ tám của dãy.
  \item Tính trung bình cộng của năm số đầu tiên của dãy ở câu a).
  \item Trong dãy Fibonacci quen thuộc $1, 1, 2, 3, 5, 8, \ldots$, số thứ $2026$ là số chẵn hay số lẻ? Giải thích.
\end{enumerate}
\end{probbox}
\probsol{
\textbf{Phân tích đề.} Mở rộng Bài 15. Câu a) chưa biết số thứ hai, ta gọi nó là $b$ và biểu diễn các số sau theo $b$.
Câu c) không thể viết ra $2026$ số, ta cần tìm quy luật chẵn -- lẻ.
\step{1}{Biểu diễn các số theo $b$}
\begin{center}
\begin{tabular}{|c|c|c|c|c|c|}
\hline
\rowcolor{tmBlue!10} Thứ tự & 1 & 2 & 3 & 4 & 5\\ \hline
Số & $2$ & $b$ & $2 + b$ & $b + (2 + b) = 2 + 2b$ & $(2 + b) + (2 + 2b) = 4 + 3b$\\ \hline
\end{tabular}
\end{center}
\step{2}{Tìm $b$}
$4 + 3b = 19 \Rightarrow 3b = 15 \Rightarrow b = 5$. Số thứ hai là $5$.
\step{3}{Viết tiếp đến số thứ tám}
$2,\ 5,\ 7,\ 12,\ 19,\ 31,\ 50,\ 81$ \quad ($12 + 19 = 31$; \ $19 + 31 = 50$; \ $31 + 50 = 81$). Số thứ tám là $81$.
\step{4}{Trung bình cộng của năm số đầu}
$(2 + 5 + 7 + 12 + 19) : 5 = 45 : 5 = 9$.
\step{5}{Quy luật chẵn -- lẻ của dãy Fibonacci}
\begin{recallbox}lẻ $+$ lẻ $=$ chẵn; \ lẻ $+$ chẵn $=$ lẻ; \ chẵn $+$ lẻ $=$ lẻ.\end{recallbox}
Viết L (lẻ), C (chẵn) cho các số $1, 1, 2, 3, 5, 8, 13, 21, 34, \ldots$:
\begin{center}
\begin{tabular}{|c|c|c|c|c|c|c|c|c|c|}
\hline
\rowcolor{tmBlue!10} Thứ tự & 1 & 2 & 3 & 4 & 5 & 6 & 7 & 8 & 9\\ \hline
Số & 1 & 1 & 2 & 3 & 5 & 8 & 13 & 21 & 34\\ \hline
Chẵn/lẻ & L & L & \textbf{C} & L & L & \textbf{C} & L & L & \textbf{C}\\ \hline
\end{tabular}
\end{center}
Mẫu ``L, L, C'' lặp lại mãi mãi.
\begin{proofbox}Giả sử tới một chỗ nào đó ta có hai số liên tiếp là L, L. Khi đó số tiếp theo là L $+$ L $=$ C; số sau nữa là
L $+$ C $=$ L; số sau nữa là C $+$ L $=$ L. Ta lại được hai số liên tiếp L, L, nên chu kì ``L, L, C'' bắt đầu lại. Dãy mở đầu
bằng $1, 1$ (L, L), nên mẫu này lặp lại từ đầu đến cuối: các số ở vị trí chia hết cho $3$ là số chẵn, các số còn lại là số lẻ.\end{proofbox}
\step{6}{Xét vị trí $2026$}
$2026 = 3\times 675 + 1$, tức $2026$ chia cho $3$ dư $1$. Vị trí này ứng với chữ ``L'' đầu tiên trong một chu kì. Vậy số thứ
$2026$ là \textbf{số lẻ}.
\begin{warnbox}Không thể kết luận ``$2026$ là số chẵn nên số thứ $2026$ cũng chẵn''. Tính chẵn -- lẻ của \emph{vị trí} khác
với tính chẵn -- lẻ của \emph{số} ở vị trí đó; ví dụ số thứ $2$ là $1$ (lẻ).\end{warnbox}
\verify{Câu a): $2 + 5 = 7$, $5 + 7 = 12$, $7 + 12 = 19$ -- đúng số thứ năm là $19$. Câu c): kiểm tra mẫu trên các vị trí
nhỏ: vị trí $4 = 3\times 1 + 1$ là số $3$ (lẻ); vị trí $7 = 3\times 2 + 1$ là số $13$ (lẻ).}
\ans{a) số thứ hai là $5$, số thứ tám là $81$;\quad b) $9$;\quad c) số lẻ}
}

% ---------------- part d ----------------
% ============================================================
% Topic 4 -- Số học và Tư duy số -- part d: problems 41--50 (VI)
% Advanced extensions of problems 16--30
% ============================================================

\begin{probbox}{41}{\stars{1}}
\textbf{Tổng các số liên tiếp} (mở rộng Bài 16).
\begin{enumerate}[label=\alph*),leftmargin=1.6em,itemsep=1pt]
\item Hãy tìm \emph{tất cả} các cách viết số $100$ thành tổng của ít nhất hai số nguyên dương liên tiếp.
\item Chứng minh rằng số $64$ không thể viết được thành tổng của ít nhất hai số nguyên dương liên tiếp.
\end{enumerate}
\end{probbox}
\probsol{
\textbf{Phân tích đề.} Một tổng các số nguyên dương liên tiếp được xác định bởi hai thông tin: \emph{có bao nhiêu số} (gọi là $k$, với $k\ge 2$) và \emph{số đầu tiên} (gọi là $a$, với $a\ge 1$). Ta cần một công thức tính nhanh tổng như vậy, rồi dùng nó để ``thu hẹp'' các khả năng.

\begin{recallbox}\textbf{Tổng các số cách đều} (cách tính của nhà toán học Gauss): tổng $=\dfrac{(\text{số đầu}+\text{số cuối})\times\text{số số hạng}}{2}$.\end{recallbox}

\begin{proofbox}Viết tổng hai lần, một lần xuôi, một lần ngược, rồi cộng theo cột: ví dụ $S=9+10+\dots+16$ và $S=16+15+\dots+9$. Mỗi cột có tổng $9+16=25$ và có $8$ cột, nên $2S=25\cdot 8$, $S=100$. Lí luận này đúng với mọi dãy cách đều.\end{proofbox}

\step{1}{Lập đẳng thức chính}
Với $k$ số liên tiếp bắt đầu từ $a$: số cuối là $a+k-1$. Gọi $m=(\text{số đầu}+\text{số cuối})=2a+k-1$. Theo công thức Gauss: tổng $=\dfrac{m\cdot k}{2}$. Nếu tổng bằng $100$ thì
\[k\cdot m=200.\]

\step{2}{Hai nhận xét quan trọng về $k$ và $m$}
\begin{itemize}[leftmargin=1.4em,itemsep=1pt]
\item $m-k=2a-1$ là số lẻ, nên \textbf{$k$ và $m$ khác tính chẵn lẻ} (một số chẵn, một số lẻ).
\item $m-k=2a-1\ge 1$, nên \textbf{$m>k$}.
\end{itemize}

\step{3}{Liệt kê các cách phân tích $200=k\cdot m$ với $2\le k<m$}
Vì $k<m$ và $k\cdot m=200$ nên $k\cdot k<200$, tức $k\le 14$. Các ước của $200$ từ $2$ đến $14$ là $2,4,5,8,10$.
\begin{center}
\begin{tabular}{|c|c|c|c|}
\hline
\rowcolor{Tint}$k$ & $m=200:k$ & Khác tính chẵn lẻ? & Kết quả\\
\hline
$2$ & $100$ & không (cùng chẵn) & loại\\
$4$ & $50$ & không & loại\\
$5$ & $40$ & có & $2a+4=40\Rightarrow a=18$\\
$8$ & $25$ & có & $2a+7=25\Rightarrow a=9$\\
$10$ & $20$ & không & loại\\
\hline
\end{tabular}
\end{center}

\step{4}{Viết ra các tổng}
\[100=18+19+20+21+22\qquad\text{và}\qquad 100=9+10+11+12+13+14+15+16.\]
Có đúng \textbf{hai} cách.

\begin{tipbox}Nhớ Bài 16: với $5$ số liên tiếp (số lẻ số hạng), số ở giữa bằng trung bình cộng $=100:5=20$, nên năm số là $18,19,20,21,22$. Ý tưởng ``số ở giữa'' giúp kiểm tra rất nhanh.\end{tipbox}

\step{5}{Câu b): vì sao $64$ không viết được}
Giả sử $64$ là tổng của $k\ge 2$ số nguyên dương liên tiếp. Lập luận như trên: $k\cdot m=2\cdot 64=128$ với $m>k\ge 2$ và $k$, $m$ khác tính chẵn lẻ. Như vậy một trong hai số $k$, $m$ là số lẻ và là ước của $128$. Nhưng $128=2^7$ chỉ có \emph{một} ước lẻ là $1$, trong khi cả $k$ và $m$ đều $\ge 2$. Mâu thuẫn! Vậy $64$ không viết được như yêu cầu.

\begin{proofbox}Cùng lí luận đó chứng tỏ: \emph{mọi lũy thừa của $2$} ($2,4,8,16,32,64,\dots$) đều không viết được thành tổng của ít nhất hai số nguyên dương liên tiếp. Ngược lại, số nào có ước lẻ lớn hơn $1$ thì đều viết được.\end{proofbox}

\begin{warnbox}Đừng quên điều kiện ``nguyên dương'': nếu cho phép số $0$ và số âm thì có thêm những cách ``giả'', ví dụ $(-17)+(-16)+\dots+21+22=100$ (các số từ $-17$ đến $17$ triệt tiêu nhau, còn lại $18+19+20+21+22$). Trong bài này mọi số phải $\ge 1$, nên $a\ge 1$, dẫn tới $m>k$.\end{warnbox}

\verify{$18+19+20+21+22=100$; $9+10+\dots+16=\dfrac{(9+16)\cdot 8}{2}=100$. Thử tay với $64$: $2$ số liên tiếp có tổng lẻ; $3$ số liên tiếp có tổng chia hết cho $3$; \dots\ không cách nào ra $64$.}
\ans{a) Hai cách: $18+19+20+21+22$ và $9+10+\dots+16$. \ b) Xem chứng minh.}
}

\begin{probbox}{42}{\stars{1}}
\textbf{Ước chung và số ước} (mở rộng Bài 17).
\begin{enumerate}[label=\alph*),leftmargin=1.6em,itemsep=1pt]
\item Có bao nhiêu số nguyên dương là ước của cả hai số $720$ và $1080$?
\item Tìm số nguyên dương nhỏ nhất có đúng $12$ ước.
\end{enumerate}
\end{probbox}
\probsol{
\textbf{Phân tích đề.} Liệt kê ước của $720$ và $1080$ thì rất dài (mỗi số có vài chục ước). Ta dùng hai công cụ mạnh: (1) ước chung của hai số chính là ước của ƯCLN; (2) công thức đếm số ước từ phân tích ra thừa số nguyên tố.

\begin{recallbox}\textbf{Đếm số ước.} Nếu $n=p^{x}\cdot q^{y}\cdot r^{z}$ ($p,q,r$ là các số nguyên tố khác nhau) thì $n$ có $(x+1)(y+1)(z+1)$ ước.\end{recallbox}

\begin{proofbox}Mỗi ước của $n$ có dạng $p^{i}\cdot q^{j}\cdot r^{l}$ với $0\le i\le x$, $0\le j\le y$, $0\le l\le z$. Số mũ $i$ có $x+1$ cách chọn ($0,1,\dots,x$), $j$ có $y+1$ cách, $l$ có $z+1$ cách. Theo quy tắc nhân có $(x+1)(y+1)(z+1)$ ước, và các lựa chọn khác nhau cho các ước khác nhau (vì phân tích ra thừa số nguyên tố là duy nhất).\end{proofbox}

\par\smallskip\textbf{Câu a).}
\step{1}{Phân tích ra thừa số nguyên tố}
$720=2^4\cdot 3^2\cdot 5$ (vì $720=16\cdot 45=16\cdot 9\cdot 5$); \quad $1080=2^3\cdot 3^3\cdot 5$ (vì $1080=8\cdot 135=8\cdot 27\cdot 5$).

\step{2}{Tìm ƯCLN}
Lấy mỗi thừa số nguyên tố chung với số mũ \emph{nhỏ nhất}: $\text{ƯCLN}(720,1080)=2^3\cdot 3^2\cdot 5=8\cdot 9\cdot 5=360$.
(Kiểm tra bằng cách khác: $1080-720=360$ và $720=2\cdot 360$, nên ƯCLN là $360$.)

\step{3}{Đếm ước của $360$}
$360=2^3\cdot 3^2\cdot 5^1$ có $(3+1)(2+1)(1+1)=4\cdot 3\cdot 2=24$ ước. Vậy có $24$ số là ước của cả $720$ và $1080$.

\par\smallskip\textbf{Câu b).}
\step{4}{Các ``dạng'' của số có đúng $12$ ước}
Số ước là tích các thừa số (số mũ $+1$). Ta cần viết $12$ thành tích các số $\ge 2$:
\[12=12,\qquad 12=6\cdot 2,\qquad 12=4\cdot 3,\qquad 12=3\cdot 2\cdot 2.\]
Tương ứng với bốn dạng: $p^{11}$, \ $p^{5}q$, \ $p^{3}q^{2}$, \ $p^{2}qr$.

\step{5}{Số nhỏ nhất của mỗi dạng}
Để số nhỏ nhất, dùng các số nguyên tố nhỏ nhất $2,3,5$, và đặt số mũ \emph{lớn} cho số nguyên tố \emph{nhỏ}:
\begin{center}
\begin{tabular}{|c|c|c|}
\hline
\rowcolor{Tint}Dạng & Số nhỏ nhất & Giá trị\\
\hline
$p^{11}$ & $2^{11}$ & $2048$\\
$p^{5}q$ & $2^5\cdot 3$ & $96$\\
$p^{3}q^{2}$ & $2^3\cdot 3^2$ & $72$\\
$p^{2}qr$ & $2^2\cdot 3\cdot 5$ & $\mathbf{60}$\\
\hline
\end{tabular}
\end{center}
Nhỏ nhất là $60$.

\begin{warnbox}Đừng đặt số mũ lớn cho số nguyên tố lớn: ví dụ $2\cdot 3^5=486$ lớn hơn nhiều so với $2^5\cdot 3=96$. Cũng đừng quên dạng có \emph{ba} thừa số nguyên tố ($p^2qr$) -- chính dạng này cho đáp số.\end{warnbox}

\verify{Các ước của $60$: $1,2,3,4,5,6,10,12,15,20,30,60$ -- đúng $12$ ước. Duyệt bằng máy tính các số từ $1$ đến $59$: không số nào có $12$ ước.}
\ans{a) $24$ số. \ b) $60$.}
}

\begin{probbox}{43}{\stars{1}}
\textbf{Phân số thay đổi thế nào?} (mở rộng Bài 18 và 19).
\begin{enumerate}[label=\alph*),leftmargin=1.6em,itemsep=1pt]
\item Khi cộng cùng một số tự nhiên $n$ vào cả tử số và mẫu số của phân số $\dfrac58$, ta được một phân số bằng $\dfrac45$. Tìm $n$.
\item Cho phân số $\dfrac ab$ với $0<a<b$. Chứng minh rằng với mọi số nguyên dương $n$ thì $\dfrac{a+n}{b+n}>\dfrac ab$.
\item Có bao nhiêu phân số bằng $\dfrac37$ mà mẫu số là một số có hai chữ số?
\end{enumerate}
\end{probbox}
\probsol{
\textbf{Phân tích đề.} Câu a) và b) cùng nói về việc ``cộng thêm cùng một số vào tử và mẫu''. Điều \emph{không thay đổi} khi làm vậy là \emph{hiệu} giữa mẫu và tử -- đây là chìa khóa.

\par\smallskip\textbf{Câu a).}
\step{1}{Tìm đại lượng không đổi}
Mẫu $-$ tử của $\dfrac58$ là $8-5=3$. Sau khi cộng $n$ vào cả hai: $(8+n)-(5+n)=3$, vẫn bằng $3$.

\step{2}{Dùng sơ đồ ``phần'' như Bài 19}
Phân số mới bằng $\dfrac45$, nên tử mới gồm $4$ phần, mẫu mới gồm $5$ phần; hiệu là $1$ phần. Mà hiệu bằng $3$, nên $1$ phần $=3$. Tử mới $=4\cdot 3=12$, mẫu mới $=5\cdot 3=15$.

\step{3}{Tìm $n$}
$n=12-5=7$ (và cũng bằng $15-8=7$).

\verify{$\dfrac{5+7}{8+7}=\dfrac{12}{15}=\dfrac{12:3}{15:3}=\dfrac45$. Đúng.}

\par\smallskip\textbf{Câu b).}
\step{4}{So sánh ``phần còn thiếu để bằng $1$''}
Vì $a<b$ nên cả hai phân số đều nhỏ hơn $1$. Phần còn thiếu để đủ $1$ là
\[1-\frac ab=\frac{b-a}{b},\qquad 1-\frac{a+n}{b+n}=\frac{(b+n)-(a+n)}{b+n}=\frac{b-a}{b+n}.\]
Hai phân số $\dfrac{b-a}{b}$ và $\dfrac{b-a}{b+n}$ có \emph{cùng tử số dương}, còn mẫu $b+n>b$. Phân số có mẫu lớn hơn thì nhỏ hơn, nên $\dfrac{b-a}{b+n}<\dfrac{b-a}{b}$.

\step{5}{Kết luận}
Phân số $\dfrac{a+n}{b+n}$ ``thiếu ít hơn'' để bằng $1$, tức là nó lớn hơn: $\dfrac{a+n}{b+n}>\dfrac ab$.

\begin{recallbox}Hai phân số dương có cùng tử số: phân số nào có mẫu lớn hơn thì nhỏ hơn (cùng một cái bánh chia cho nhiều người hơn thì mỗi phần nhỏ hơn).\end{recallbox}

\begin{proofbox}Cách khác (nhân chéo): với các mẫu dương, $\dfrac{a+n}{b+n}>\dfrac ab$ khi và chỉ khi $(a+n)\cdot b>a\cdot(b+n)$, tức $ab+nb>ab+na$, tức $nb>na$. Điều này đúng vì $n>0$ và $b>a$.\end{proofbox}

\begin{warnbox}Đây chính là lí do ở Bài 18 ta \emph{không được} ``cộng thêm'' vào tử và mẫu: cộng thêm làm phân số thay đổi (nhỏ hơn $1$ thì tăng lên). Chỉ \emph{nhân} (hoặc chia) cả tử và mẫu với cùng một số khác $0$ mới giữ nguyên phân số.\end{warnbox}

\par\smallskip\textbf{Câu c).}
\step{6}{Dạng chung và điều kiện}
Các phân số bằng $\dfrac37$ có dạng $\dfrac{3k}{7k}$ ($k$ nguyên dương). Mẫu $7k$ có hai chữ số: $10\le 7k\le 99$. Ta có $7\cdot 1=7<10$, $7\cdot 2=14\ge 10$; $7\cdot 14=98\le 99$, $7\cdot 15=105>99$. Vậy $k=2,3,\dots,14$: có $14-2+1=13$ giá trị.
Các phân số là $\dfrac{6}{14},\dfrac{9}{21},\dots,\dfrac{42}{98}$.

\ans{a) $n=7$. \ b) Xem chứng minh. \ c) $13$ phân số.}
}

\begin{probbox}{44}{\stars{1}}
\textbf{Cột điện, xe đạp và đoàn tàu} (mở rộng Bài 20).
\begin{enumerate}[label=\alph*),leftmargin=1.6em,itemsep=1pt]
\item Bạn An đạp xe đều với vận tốc $18$ km/h dọc một con đường có các cột điện cách đều nhau $50$ m. Hỏi từ lúc An đi qua cột điện thứ nhất đến lúc An đi qua cột điện thứ $25$ mất bao nhiêu phút?
\item Một đoàn tàu chạy đều đi qua một cột điện hết $9$ giây và đi qua một cây cầu dài $350$ m hết $30$ giây (tính từ lúc đầu tàu bắt đầu lên cầu đến lúc đuôi tàu vừa rời khỏi cầu). Tính chiều dài đoàn tàu và vận tốc của tàu theo km/h.
\end{enumerate}
\end{probbox}
\probsol{
\textbf{Phân tích đề.} Câu a) có ``bẫy trồng cây'': $25$ cột thì có bao nhiêu khoảng? Câu b) cần hiểu rõ khi ``đi qua'' một vật, đoàn tàu (một vật có chiều dài) phải đi được quãng đường bao nhiêu.

\begin{recallbox}Quãng đường $=$ vận tốc $\times$ thời gian. Đổi $1$ km/h: $1000$ m trong $3600$ giây, nên $1$ km/h $=\dfrac{1000}{3600}=\dfrac{5}{18}$ m/giây; ngược lại $1$ m/giây $=3{,}6$ km/h.\end{recallbox}

\par\smallskip\textbf{Câu a).}
\step{1}{Đếm số khoảng giữa các cột}
\begin{figC}
\foreach \x in {0,1,2,3}{\draw[Navy,line width=1.5pt] (\x*1.2,0)--(\x*1.2,0.8); \node[below,font=\scriptsize] at (\x*1.2,0) {cột \pgfmathparse{int(\x+1)}\pgfmathresult};}
\foreach \x in {0,1,2}{\draw[<->,tmOrange] (\x*1.2+0.05,0.4)--(\x*1.2+1.15,0.4); \node[above,font=\scriptsize,text=tmOrange!80!black] at (\x*1.2+0.6,0.4) {$50$ m};}
\node at (4.35,0.4) {$\cdots$};
\node[anchor=west,font=\footnotesize,text=Ink] at (5,0.4) {$4$ cột $\to$ $3$ khoảng};
\end{figC}
Giữa cột thứ nhất và cột thứ $25$ có $25-1=24$ khoảng, tổng chiều dài $24\cdot 50=1200$ m.

\step{2}{Đổi vận tốc và tính thời gian}
$18$ km/h $=18\cdot\dfrac{5}{18}=5$ m/giây. Thời gian: $1200:5=240$ giây $=240:60=4$ phút.

\begin{warnbox}Nếu tính $25$ khoảng ($25\cdot 50=1250$ m) sẽ được $250$ giây -- sai. Số khoảng luôn ít hơn số cột $1$ đơn vị (khi tính từ cột đầu đến cột cuối).\end{warnbox}

\par\smallskip\textbf{Câu b).}
\step{3}{Hiểu ``đi qua cột điện'' và ``đi qua cây cầu''}
\begin{itemize}[leftmargin=1.4em,itemsep=1pt]
\item Đi qua cột điện: từ lúc đầu tàu ngang cột đến lúc đuôi tàu ngang cột, tàu đi được đúng \emph{một chiều dài tàu} $L$.
\item Đi qua cầu: từ lúc đầu tàu ở đầu cầu đến lúc đuôi tàu rời cuối cầu, tàu đi được \emph{chiều dài cầu cộng chiều dài tàu}: $350+L$.
\end{itemize}
\begin{figC}
\fill[Slate!25] (2,-0.15) rectangle (7,0); \node[below,font=\scriptsize] at (4.5,-0.15) {cầu $350$ m};
\draw[tmBlue,fill=tmBlue!20,rounded corners=2pt] (0,0.05) rectangle (2,0.45); \node[font=\scriptsize] at (1,0.25) {tàu (lúc đầu)};
\draw[tmOrange,fill=tmOrange!20,rounded corners=2pt] (7,0.05) rectangle (9,0.45); \node[font=\scriptsize] at (8,0.25) {tàu (lúc cuối)};
\draw[<->,Navy] (2,0.8)--(9,0.8) node[midway,above,font=\scriptsize] {đầu tàu đi được $350+L$};
\draw[dashed,Slate] (2,0.05)--(2,0.8); \draw[dashed,Slate] (9,0.05)--(9,0.8);
\end{figC}

\step{4}{So sánh hai lần}
Qua cầu mất $30$ giây, qua cột mất $9$ giây. Thời gian chênh lệch $30-9=21$ giây chính là thời gian tàu đi thêm $350$ m (chiều dài cầu). Vận tốc tàu:
\[v=350:21=\frac{350}{21}=\frac{50}{3}\ \text{m/giây}.\]

\step{5}{Tính chiều dài tàu và đổi vận tốc}
Chiều dài tàu $L=v\cdot 9=\dfrac{50}{3}\cdot 9=150$ m. Vận tốc: $\dfrac{50}{3}\cdot 3{,}6=\dfrac{50\cdot 3{,}6}{3}=60$ km/h.

\begin{warnbox}Lỗi hay gặp: cho rằng qua cầu tàu chỉ đi $350$ m (quên cộng chiều dài tàu), khi đó vận tốc là $350:30\approx 11{,}7$ m/giây và kết quả sai. Hãy theo dõi \emph{một điểm cố định} trên tàu (ví dụ đầu tàu) để biết tàu đã đi bao xa.\end{warnbox}

\verify{Với $v=\dfrac{50}{3}$ m/giây và $L=150$ m: qua cột mất $150:\dfrac{50}{3}=9$ giây; qua cầu mất $(350+150):\dfrac{50}{3}=500\cdot\dfrac{3}{50}=30$ giây. Đúng cả hai dữ kiện.}
\ans{a) $4$ phút. \ b) Tàu dài $150$ m, vận tốc $60$ km/h.}
}

\begin{probbox}{45}{\stars{2}}
\textbf{Tích các chữ số} (mở rộng Bài 21).
\begin{enumerate}[label=\alph*),leftmargin=1.6em,itemsep=1pt]
\item Có bao nhiêu số có bốn chữ số mà tích các chữ số bằng $36$?
\item Tìm số nguyên dương nhỏ nhất có tích các chữ số bằng $360$.
\item Tìm số nguyên dương lớn nhất có các chữ số \emph{đôi một khác nhau} và tích các chữ số bằng $360$.
\end{enumerate}
\end{probbox}
\probsol{
\textbf{Phân tích đề.} Cả ba câu đều bắt đầu bằng việc \emph{phân tích số ra thừa số nguyên tố}, rồi ``chia'' các thừa số đó vào các chữ số (mỗi chữ số từ $1$ đến $9$). Chú ý: các chữ số đều khác $0$ (nếu không tích bằng $0$).

\begin{recallbox}$36=2^2\cdot 3^2$; \ $360=2^3\cdot 3^2\cdot 5$. Các chữ số và phân tích của chúng: $1$; $2$; $3$; $4=2^2$; $5$; $6=2\cdot 3$; $7$; $8=2^3$; $9=3^2$.\end{recallbox}

\par\smallskip\textbf{Câu a).}
\step{1}{Tìm các bộ bốn chữ số (không kể thứ tự) có tích $36$}
Làm như Bài 21, xếp $a\le b\le c\le d$:
\begin{itemize}[leftmargin=1.4em,itemsep=1pt]
\item Có hai chữ số $1$: còn lại tích $36$ với hai chữ số: $\{1,1,4,9\}$, $\{1,1,6,6\}$.
\item Có đúng một chữ số $1$: ba chữ số còn lại ($\ge 2$) có tích $36$: $\{1,2,2,9\}$, $\{1,2,3,6\}$, $\{1,3,3,4\}$.
\item Không có chữ số $1$: bốn chữ số $\ge 2$ có tích $36$: chỉ có $\{2,2,3,3\}$ (vì $36=2\cdot2\cdot3\cdot3$, không thể tách thêm).
\item Có ba chữ số $1$: chữ số còn lại phải là $36$ -- không được.
\end{itemize}

\begin{recallbox}\textbf{Đếm cách xếp $4$ chữ số:} cả bốn khác nhau: $4\cdot 3\cdot 2\cdot 1=24$ cách. Có đúng hai chữ số giống nhau (hai cái còn lại khác nhau và khác chúng): chọn $2$ vị trí cho cặp giống nhau trong $4$ vị trí ($6$ cách), rồi xếp hai chữ số còn lại vào $2$ chỗ trống ($2$ cách): $12$ cách. Hai cặp giống nhau (như $1,1,6,6$): chỉ cần chọn $2$ vị trí cho một cặp: $6$ cách.\end{recallbox}

\step{2}{Đếm}
\begin{center}
\begin{tabular}{|c|c|c|}
\hline
\rowcolor{Tint}Bộ chữ số & Loại & Số cách xếp\\
\hline
$\{1,1,4,9\}$ & một cặp giống & $12$\\
$\{1,1,6,6\}$ & hai cặp giống & $6$\\
$\{1,2,2,9\}$ & một cặp giống & $12$\\
$\{1,2,3,6\}$ & khác nhau & $24$\\
$\{1,3,3,4\}$ & một cặp giống & $12$\\
$\{2,2,3,3\}$ & hai cặp giống & $6$\\
\hline
\multicolumn{2}{|r|}{\textbf{Tổng}} & $\mathbf{72}$\\
\hline
\end{tabular}
\end{center}

\par\smallskip\textbf{Câu b).}
\step{3}{Cần ít chữ số nhất}
Số có ít chữ số hơn thì nhỏ hơn. Một hoặc hai chữ số có tích lớn nhất $9\cdot 9=81<360$, nên cần ít nhất $3$ chữ số.

\step{4}{Tìm bộ ba chữ số có tích $360$}
$360$ có thừa số nguyên tố $5$, mà chữ số duy nhất chứa thừa số $5$ là chính chữ số $5$. Vậy một chữ số là $5$, hai chữ số còn lại có tích $360:5=72=8\cdot 9$ (cách duy nhất với các chữ số $\le 9$). Bộ chữ số: $\{5,8,9\}$.

\step{5}{Xếp để được số nhỏ nhất}
Xếp các chữ số từ nhỏ đến lớn: $589$.

\par\smallskip\textbf{Câu c).}
\step{6}{Muốn số lớn nhất: càng nhiều chữ số càng tốt}
Các chữ số khác nhau, đều khác $0$, và một chữ số phải là $5$. Sáu chữ số khác nhau thì tích ít nhất là $1\cdot2\cdot3\cdot4\cdot5\cdot6=720>360$, nên có nhiều nhất $5$ chữ số.

\step{7}{Tìm các bộ $5$ chữ số khác nhau có tích $360$}
Một chữ số là $5$; bốn chữ số còn lại khác nhau, khác $5$, có tích $72=2^3\cdot 3^2$. Bốn chữ số khác nhau có tích $72$ mà nhỏ như vậy thì phải có chữ số $1$ (vì $2\cdot3\cdot4\cdot6=144>72$); ba chữ số còn lại khác nhau có tích $72$: $72=2\cdot4\cdot9=3\cdot4\cdot6$ (còn $2\cdot 6\cdot 6$, $3\cdot 3\cdot 8$ có chữ số lặp; $1\cdot 8\cdot 9$ lặp chữ số $1$). Được hai bộ: $\{1,2,4,5,9\}$ và $\{1,3,4,5,6\}$.

\step{8}{Xếp để được số lớn nhất}
Xếp từ lớn đến nhỏ: $95421$ và $65431$. Số lớn hơn là $95421$.

\begin{warnbox}Ở câu c), đừng chọn ngay số có chữ số đầu lớn như $985$ (tích $360$): số có $5$ chữ số luôn lớn hơn số có $3$ chữ số. Hãy ưu tiên \emph{số chữ số} trước, rồi mới đến chữ số đầu.\end{warnbox}

\verify{$5\cdot 8\cdot 9=360$; $9\cdot5\cdot4\cdot2\cdot1=360$. Máy tính duyệt mọi số có tới $6$ chữ số xác nhận: có $72$ số bốn chữ số tích $36$; số nhỏ nhất tích $360$ là $589$; số lớn nhất có chữ số khác nhau, tích $360$ là $95\,421$.}
\ans{a) $72$ số. \ b) $589$. \ c) $95\,421$.}
}

\begin{probbox}{46}{\stars{1}}
\textbf{Bảng cách đều} (mở rộng Bài 22). Trong một bảng ô vuông, các số trên mỗi hàng và trên mỗi cột đều tạo thành dãy cách đều.
\begin{enumerate}[label=\alph*),leftmargin=1.6em,itemsep=1pt]
\item Bảng $4\times 4$ có bốn số ở bốn góc là: góc trên trái $1$, góc trên phải $10$, góc dưới trái $19$, góc dưới phải $37$. Hãy điền đầy đủ bảng.
\item Tính tổng của tất cả $16$ số trong bảng ở câu a). So sánh với tổng bốn số ở góc, em thấy gì? Giải thích.
\item Chứng minh rằng trong một bảng $3\times 3$ như vậy, số ở ô chính giữa luôn bằng trung bình cộng của bốn số ở bốn góc.
\end{enumerate}
\begin{figC}
\foreach \rr in {0,1,2,3}{\foreach \cc in {0,1,2,3}{\draw[Navy,line width=0.9pt,fill=white] (\cc*1.1,-\rr*0.8) rectangle (\cc*1.1+1.1,-\rr*0.8-0.8);}}
\node at (0.55,-0.4) {$1$}; \node at (3.85,-0.4) {$10$};
\node at (0.55,-2.8) {$19$}; \node at (3.85,-2.8) {$37$};
\end{figC}
\end{probbox}
\probsol{
\textbf{Phân tích đề.} Bảng chỉ cho bốn số ở góc. Cột đầu và cột cuối (cũng như hàng đầu, hàng cuối) mỗi cái đã biết số đầu và số cuối, nên tìm được ngay. Sau đó dùng các hàng hoặc các cột ở giữa.

\begin{recallbox}Dãy cách đều $4$ số hạng: khoảng cách $d=(\text{số cuối}-\text{số đầu}):3$. Tổng $4$ số hạng $=2\times(\text{số đầu}+\text{số cuối})$ (ghép số $1$ với số $4$, số $2$ với số $3$: hai cặp có tổng bằng nhau). Dãy $3$ số hạng: số giữa $=$ trung bình cộng hai số ngoài.\end{recallbox}

\par\smallskip\textbf{Câu a).}
\step{1}{Hàng đầu và hàng cuối}
Hàng $1$: từ $1$ đến $10$, $d=(10-1):3=3$: \ $1,4,7,10$.\\
Hàng $4$: từ $19$ đến $37$, $d=(37-19):3=6$: \ $19,25,31,37$.

\step{2}{Bốn cột (biết số đầu và số cuối của mỗi cột)}
Cột $1$: $1\to 19$, $d=6$: $1,7,13,19$. \quad Cột $2$: $4\to 25$, $d=7$: $4,11,18,25$.\\
Cột $3$: $7\to 31$, $d=8$: $7,15,23,31$. \quad Cột $4$: $10\to 37$, $d=9$: $10,19,28,37$.

\step{3}{Kiểm tra các hàng ở giữa}
Hàng $2$: $7,11,15,19$ (cách đều $4$); hàng $3$: $13,18,23,28$ (cách đều $5$). Bảng thỏa mãn mọi điều kiện:
\begin{figC}
\foreach \rr in {0,1,2,3}{\foreach \cc in {0,1,2,3}{\draw[Navy,line width=0.9pt,fill=white] (\cc*1.1,-\rr*0.8) rectangle (\cc*1.1+1.1,-\rr*0.8-0.8);}}
\foreach \cc/\val in {0/1,1/4,2/7,3/10}{\node at (\cc*1.1+0.55,-0.4) {$\val$};}
\foreach \cc/\val in {0/7,1/11,2/15,3/19}{\node at (\cc*1.1+0.55,-1.2) {$\val$};}
\foreach \cc/\val in {0/13,1/18,2/23,3/28}{\node at (\cc*1.1+0.55,-2.0) {$\val$};}
\foreach \cc/\val in {0/19,1/25,2/31,3/37}{\node at (\cc*1.1+0.55,-2.8) {$\val$};}
\foreach \rr/\val in {0/22,1/52,2/82,3/112}{\node[anchor=west,font=\footnotesize,text=Slate] at (4.55,-\rr*0.8-0.4) {tổng hàng: $\val$};}
\end{figC}

\par\smallskip\textbf{Câu b).}
\step{4}{Tính tổng}
Tổng các hàng: $22+52+82+112=268$. Tổng bốn góc: $1+10+19+37=67$. Ta thấy $268=4\cdot 67$: \textbf{tổng cả bảng bằng $4$ lần tổng bốn góc.}

\step{5}{Giải thích (đúng với mọi bảng $4\times4$ như vậy)}
Mỗi hàng là dãy cách đều $4$ số, nên tổng mỗi hàng $=2\times$(số đầu hàng $+$ số cuối hàng). Cộng $4$ hàng:
\[\text{Tổng bảng}=2\times(\text{tổng cột }1+\text{tổng cột }4).\]
Cột $1$ và cột $4$ cũng là dãy cách đều $4$ số: tổng cột $1=2\times(1+19)$, tổng cột $4=2\times(10+37)$. Vậy
\[\text{Tổng bảng}=2\times\big(2\cdot(1+19)+2\cdot(10+37)\big)=4\times(1+10+19+37)=4\times 67=268.\]

\par\smallskip\textbf{Câu c).}
\step{6}{Đặt tên các ô của bảng $3\times 3$}
\begin{figC}
\foreach \rr in {0,1,2}{\foreach \cc in {0,1,2}{\draw[Navy,line width=0.9pt,fill=white] (\cc*1.1,-\rr*0.8) rectangle (\cc*1.1+1.1,-\rr*0.8-0.8);}}
\node at (0.55,-0.4) {$a$}; \node at (1.65,-0.4) {$\cdot$}; \node at (2.75,-0.4) {$b$};
\node at (0.55,-1.2) {$u$}; \node[text=AnsRed] at (1.65,-1.2) {$x$}; \node at (2.75,-1.2) {$v$};
\node at (0.55,-2.0) {$c$}; \node at (1.65,-2.0) {$\cdot$}; \node at (2.75,-2.0) {$d$};
\end{figC}
Bốn góc là $a,b,c,d$; ô giữa là $x$; $u$ và $v$ là ô giữa của cột trái và cột phải.

\step{7}{Dùng tính chất ``số giữa $=$ trung bình cộng hai số ngoài''}
Cột trái $a,u,c$ cách đều nên $u=\dfrac{a+c}{2}$. Cột phải: $v=\dfrac{b+d}{2}$. Hàng giữa $u,x,v$ cách đều nên
\[x=\frac{u+v}{2}=\frac{\frac{a+c}{2}+\frac{b+d}{2}}{2}=\frac{a+b+c+d}{4}.\]
Đó chính là trung bình cộng của bốn số ở góc.

\begin{proofbox}Vì sao trong dãy cách đều $p,q,r$ ta có $q=\dfrac{p+r}{2}$? Vì $q-p=r-q$ (cùng bằng khoảng cách $d$), nên $2q=p+r$.\end{proofbox}

\verify{Áp dụng câu c) cho phần $3\times 3$ ở góc trên trái của bảng câu a): bốn góc $1,7,13,23$, trung bình cộng $\dfrac{1+7+13+23}{4}=11$, đúng bằng ô giữa $11$.}
\ans{a) Xem bảng. \ b) Tổng $268=4\times 67$ (gấp $4$ lần tổng bốn góc). \ c) Xem chứng minh.}
}

\begin{probbox}{47}{\stars{2}}
\textbf{Số ``xoay vòng''} (mở rộng Bài 23). Trong các câu dưới đây, mỗi chữ cái biểu thị một chữ số, và $A\neq 0$.
\begin{enumerate}[label=\alph*),leftmargin=1.6em,itemsep=1pt]
\item Tìm số có sáu chữ số $\nd{2ABCDE}$ sao cho $\nd{2ABCDE}\times 3=\nd{ABCDE2}$.
\item Tìm tất cả các chữ số $d$ (từ $1$ đến $9$) sao cho có số có sáu chữ số $\nd{dABCDE}$ thỏa mãn $\nd{dABCDE}\times 3=\nd{ABCDEd}$.
\item Chứng minh rằng không có số có sáu chữ số $\nd{1ABCDE}$ nào thỏa mãn $\nd{1ABCDE}\times 2=\nd{ABCDE1}$.
\end{enumerate}
\end{probbox}
\probsol{
\textbf{Phân tích đề.} Cả ba câu đều là ``chuyển chữ số đầu xuống cuối''. Ta dùng lại Cách 1 của Bài 23: coi khối $\nd{ABCDE}$ là một số $N$ rồi lập đẳng thức. Lần này ta làm một lần cho mọi chữ số đầu $d$.

\begin{recallbox}Với $N=\nd{ABCDE}$: \ $\nd{dABCDE}=d\cdot 100\,000+N$ \ và \ $\nd{ABCDEd}=10\cdot N+d$. Vì $A\neq 0$ nên $N$ là số có đúng năm chữ số: $10\,000\le N\le 99\,999$.\end{recallbox}

\step{1}{Lập đẳng thức tổng quát cho câu a) và b)}
\[3\cdot(d\cdot 100\,000+N)=10N+d\ \Longrightarrow\ 300\,000\,d+3N=10N+d\ \Longrightarrow\ 299\,999\,d=7N.\]
Vì $299\,999=7\cdot 42\,857$ nên chia hai vế cho $7$:
\[N=42\,857\cdot d.\]

\step{2}{Câu a): $d=2$}
$N=42\,857\cdot 2=85\,714$. Số cần tìm là $285\,714$.

\step{3}{Câu b): điều kiện để $N$ có năm chữ số}
$N=42\,857\,d\le 99\,999$. Với $d=1$: $N=42\,857$ (năm chữ số, được). Với $d=2$: $N=85\,714$ (được). Với $d=3$: $N=128\,571$ có sáu chữ số -- quá lớn, không thể là khối $\nd{ABCDE}$. Với $d\ge 3$ thì $N$ còn lớn hơn nữa. Vậy chỉ có $d=1$ và $d=2$, ứng với hai số $142\,857$ và $285\,714$.

\begin{tipbox}Hai số này có cùng các chữ số $1,4,2,8,5,7$ xếp vòng tròn. Đó không phải ngẫu nhiên: $\dfrac17=0{,}(142857)$ và $\dfrac27=0{,}(285714)$. Nhân với $3$ ứng với việc ``xoay'' vòng tròn chữ số.\end{tipbox}

\step{4}{Câu c): lập đẳng thức}
Với $N=\nd{ABCDE}$:
\[2\cdot(100\,000+N)=10N+1\ \Longrightarrow\ 200\,000+2N=10N+1\ \Longrightarrow\ 199\,999=8N.\]

\step{5}{Câu c): tìm mâu thuẫn chẵn -- lẻ}
Vế phải $8N$ là số chẵn (vì $8$ chẵn). Vế trái $199\,999$ là số lẻ. Một số lẻ không thể bằng một số chẵn. Vậy không có số nào thỏa mãn.

\begin{proofbox}Có thể thấy ngay mà không cần lập đẳng thức: $\nd{1ABCDE}\times 2$ là số chẵn, nên chữ số tận cùng của nó phải chẵn; nhưng $\nd{ABCDE1}$ tận cùng bằng $1$ (lẻ). Mâu thuẫn.\end{proofbox}

\begin{warnbox}Ở câu b), nhiều bạn dừng lại khi tìm được $d=1$ (đáp án của Bài 23) và quên $d=2$; hoặc cho $d=3,4,\dots$ cũng được mà không kiểm tra $N$ có quá năm chữ số. Luôn kiểm tra điều kiện ``$N$ có đúng năm chữ số''.\end{warnbox}

\verify{$142\,857\times 3=428\,571$; \ $285\,714\times 3=857\,142$. Cả hai đều đúng là ``chữ số đầu chuyển xuống cuối''.}
\ans{a) $285\,714$. \ b) $d=1$ hoặc $d=2$ (các số $142\,857$ và $285\,714$). \ c) Xem chứng minh.}
}

\begin{probbox}{48}{\stars{2}}
\textbf{Chú ếch nhảy trên trục số} (mở rộng Bài 24). Một chú ếch đứng ở điểm $0$ trên trục số. Ở lần nhảy thứ $k$ ($k=1,2,3,\dots$), ếch nhảy đúng $k$ đơn vị, sang trái hoặc sang phải tùy ý.
\begin{enumerate}[label=\alph*),leftmargin=1.6em,itemsep=1pt]
\item Sau đúng $10$ lần nhảy, ếch có thể trở về điểm $0$ không? Vì sao?
\item Sau đúng $10$ lần nhảy, ếch có thể đứng ở những điểm nào? Có tất cả bao nhiêu điểm như vậy?
\item Ếch cần ít nhất bao nhiêu lần nhảy để đến được điểm $20$?
\end{enumerate}
\end{probbox}
\probsol{
\textbf{Phân tích đề.} Vị trí của ếch sau $n$ lần nhảy là $\pm1\pm2\pm3\cdots\pm n$ (dấu $+$ khi nhảy sang phải, dấu $-$ khi nhảy sang trái). Ý tưởng giống Bài 24: bắt đầu từ ``tất cả đều cộng'' rồi xem đổi dấu làm thay đổi kết quả thế nào.

\begin{recallbox}$1+2+\dots+n=\dfrac{n(n+1)}{2}$ (công thức Gauss). Ví dụ $1+2+\dots+10=\dfrac{10\cdot 11}{2}=55$.\end{recallbox}

\step{1}{Quy tắc ``đổi dấu''}
Nếu ếch nhảy hết sang phải, sau $10$ lần nó ở điểm $55$. Đổi một lần nhảy $k$ từ phải sang trái làm vị trí giảm $2k$ (mất $k$ và còn lùi thêm $k$). Vậy nếu tổng các bước nhảy sang trái là $s$ thì vị trí cuối là
\[55-2s.\]

\par\smallskip\textbf{Câu a).}
\step{2}{Tính chẵn lẻ}
$55-2s$ luôn là số \emph{lẻ} (số lẻ trừ số chẵn). Điểm $0$ là số chẵn, nên ếch \textbf{không thể} về điểm $0$ sau $10$ lần nhảy.

\par\smallskip\textbf{Câu b).}
\step{3}{$s$ có thể nhận những giá trị nào?}
$s$ là tổng của một nhóm (có thể rỗng) các số khác nhau trong $\{1,2,\dots,10\}$, nên $0\le s\le 55$. Ta chứng tỏ \emph{mọi} số nguyên $s$ từ $0$ đến $55$ đều đạt được.

\begin{proofbox}Lấy dần các số lớn nhất: $10, 9, 8, \dots$ Các tổng $0$, $10$, $10+9=19$, $19+8=27$, $27+7=34$, $34+6=40$, $40+5=45$, $45+4=49$, $49+3=52$, $52+2=54$, $54+1=55$ chia đoạn $[0;55]$ thành các khúc. Một số $s$ nằm giữa hai mốc liên tiếp, ví dụ $27<s\le 34$, thì viết $s=10+9+8+(s-27)$ với $s-27$ từ $1$ đến $7$ -- là một số chưa dùng. Tương tự cho mọi khúc khác (với $0<s\le 10$ chỉ cần lấy số $s$). Vậy mọi $s$ từ $0$ đến $55$ đều đạt được.\end{proofbox}

\step{4}{Các vị trí cuối cùng}
Với $s=0,1,2,\dots,55$, vị trí $55-2s$ lần lượt là $55,53,51,\dots,-53,-55$: tức là \textbf{mọi số lẻ từ $-55$ đến $55$}. Số các điểm: từ $-55$ đến $55$ có $28$ số lẻ dương ($1,3,\dots,55$) và $28$ số lẻ âm, tổng cộng $56$ điểm (cũng bằng số giá trị của $s$: $56$).

\begin{figC}
\draw[->,thick,Navy] (-5.9,0)--(5.9,0);
\foreach \x in {-5.5,-4.5,-3.5,-2.5,-1.5,-0.5,0.5,1.5,2.5,3.5,4.5,5.5}{\fill[KidOrange] (\x,0) circle (2.3pt);}
\draw[Navy] (0,0.1)--(0,-0.1); \node[below,font=\scriptsize] at (0,-0.1) {$0$};
\node[below,font=\scriptsize] at (-5.5,-0.1) {$-55$}; \node[below,font=\scriptsize] at (5.5,-0.1) {$55$};
\node[below,font=\scriptsize] at (0.5,-0.1) {$1$}; \node[below,font=\scriptsize] at (-0.5,-0.1) {$-1$};
\node[above,font=\scriptsize,text=Slate] at (0,0.15) {chỉ các số lẻ (hình minh họa, không đúng tỉ lệ)};
\end{figC}

\par\smallskip\textbf{Câu c).}
\step{5}{Hai điều kiện cần}
Sau $n$ lần nhảy, gọi $T=1+2+\dots+n$. Vị trí là $T-2s$, nên: (1) $T\ge 20$ (không thể đi xa hơn $T$); (2) $T-20$ phải chẵn, tức $T$ và $20$ cùng chẵn.

\step{6}{Thử lần lượt $n$}
\begin{center}
\begin{tabular}{|c|c|c|c|c|c|c|c|}
\hline
\rowcolor{Tint}$n$ & $1$ & $2$ & $3$ & $4$ & $5$ & $6$ & $7$\\
\hline
$T$ & $1$ & $3$ & $6$ & $10$ & $15$ & $21$ & $28$\\
\hline
Đạt $20$? & \multicolumn{5}{c|}{không ($T<20$)} & không ($T$ lẻ) & \textbf{có thể}\\
\hline
\end{tabular}
\end{center}

\step{7}{Chỉ ra cách nhảy với $n=7$}
Cần $28-2s=20$, tức $s=4$: chỉ cần nhảy sang trái ở lần thứ $4$:
\[1+2+3-4+5+6+7=20.\]
Vậy ít nhất là $7$ lần nhảy.

\begin{warnbox}Không đủ chỉ kiểm tra ``$T\ge 20$'': với $n=6$ thì $T=21\ge 20$ nhưng vị trí luôn lẻ, không thể là $20$. Đây chính là ý tưởng chẵn -- lẻ của Bài 24.\end{warnbox}

\verify{$1+2+3-4+5+6+7=20$. Với $n=6$: vị trí luôn là $21-2s$, là số lẻ, khác $20$. Máy tính duyệt mọi cách nhảy cũng xác nhận: sau $10$ lần có đúng $56$ vị trí, không có $0$; số lần ít nhất để tới $20$ là $7$.}
\ans{a) Không (vị trí luôn lẻ). \ b) Mọi số lẻ từ $-55$ đến $55$: $56$ điểm. \ c) $7$ lần.}
}

\begin{probbox}{49}{\stars{1}}
\textbf{Số chứa khối $2025$} (mở rộng Bài 28). Ta xét các số có sáu chữ số chứa bốn chữ số $2,0,2,5$ \emph{đứng liền nhau theo đúng thứ tự này} (ví dụ: $720\,251$, $202\,513$, $132\,025$). Có bao nhiêu số như vậy chia hết cho $45$?
\end{probbox}
\probsol{
\textbf{Phân tích đề.} Số có dạng khối $2025$ cộng thêm hai chữ số $a,b$ ở các vị trí khác nhau. ``Chia hết cho $45$'' cần được tách thành hai điều kiện dễ kiểm tra hơn.

\begin{recallbox}$45=5\cdot 9$, và $5$, $9$ không có ước chung nào khác $1$. Vì vậy một số chia hết cho $45$ khi và chỉ khi nó chia hết cho cả $5$ và $9$.
\textbf{Chia hết cho $5$:} chữ số tận cùng là $0$ hoặc $5$.
\textbf{Chia hết cho $9$:} tổng các chữ số chia hết cho $9$.\end{recallbox}

\step{1}{Điều kiện chia hết cho $9$}
Tổng chữ số của khối: $2+0+2+5=9$. Tổng chữ số của cả số là $9+a+b$, chia hết cho $9$ khi $a+b$ chia hết cho $9$, tức $a+b\in\{0;\ 9;\ 18\}$.

\step{2}{Ba vị trí của khối và đếm từng trường hợp}
\begin{itemize}[leftmargin=1.4em,itemsep=2pt]
\item \textbf{Dạng $\nd{2025ab}$} (chữ số tận cùng là $b$, nên $b\in\{0;5\}$; $a$ tùy ý):
  $b=0$: $a\in\{0;9\}$ $\to$ $202\,500$, $202\,590$. \ $b=5$: $a+5\in\{9;18\}$ $\to$ $a=4$ $\to$ $202\,545$. \ \textbf{Có $3$ số.}
\item \textbf{Dạng $\nd{a2025b}$} ($a\neq 0$, $b\in\{0;5\}$):
  $b=0$: $a=9$ $\to$ $920\,250$. \ $b=5$: $a=4$ $\to$ $420\,255$. \ \textbf{Có $2$ số.}
\item \textbf{Dạng $\nd{ab2025}$} (tận cùng là $5$ -- luôn chia hết cho $5$; $a\neq 0$):
  $a+b=9$: $a=1,2,\dots,9$ (với $b=9-a$) $\to$ $9$ số: $182\,025$, $272\,025$, \dots, $902\,025$. \ $a+b=18$: $a=b=9$ $\to$ $992\,025$. \ ($a+b=0$ không được vì $a\neq0$.) \ \textbf{Có $10$ số.}
\end{itemize}

\step{3}{Kiểm tra đếm trùng}
Một số bị đếm hai lần nếu khối $2025$ xuất hiện hai lần trong sáu chữ số, tức hai khối chồng lên nhau. Dịch khối đi $1$ vị trí: cần ``$025$'' trùng ``$202$'' -- sai; dịch $2$ vị trí: cần ``$25$'' trùng ``$20$'' -- sai. Vậy không có số nào bị đếm trùng.

\step{4}{Cộng lại}
$3+2+10=15$ số.

\begin{warnbox}(1) Đừng quên trường hợp $a+b=0$ ở dạng $\nd{2025ab}$ (số $202\,500$). (2) Ở dạng $\nd{a2025b}$ và $\nd{ab2025}$, chữ số đầu $a$ phải khác $0$. (3) Phải kiểm tra \emph{cả hai} điều kiện chia hết cho $5$ và cho $9$, không chỉ một.\end{warnbox}

\verify{Ví dụ $992\,025$: tận cùng $5$, tổng chữ số $9+9+2+0+2+5=27$ chia hết cho $9$, nên chia hết cho $45$ ($992\,025=45\cdot 22\,045$). Máy tính duyệt mọi số có sáu chữ số cũng được đúng $15$ số.}
\ans{$15$ số.}
}

\begin{probbox}{50}{\stars{2}}
\textbf{Trung bình cộng và trung vị} (mở rộng Bài 30). Một danh sách gồm năm số nguyên dương (có thể có số trùng nhau) có trung bình cộng bằng $12$ và trung vị bằng $10$.
\begin{enumerate}[label=\alph*),leftmargin=1.6em,itemsep=1pt]
\item Số lớn nhất trong danh sách có thể lớn nhất là bao nhiêu?
\item Số lớn nhất trong danh sách có thể nhỏ nhất là bao nhiêu?
\item Trả lời lại hai câu hỏi trên nếu năm số đôi một khác nhau.
\end{enumerate}
\end{probbox}
\probsol{
\textbf{Phân tích đề.} Xếp năm số từ nhỏ đến lớn: $a\le b\le c\le d\le e$. Trung vị là số giữa nên $c=10$. Trung bình cộng $12$ nên tổng là $5\cdot 12=60$. Ta cần làm cho $e$ lớn nhất hoặc nhỏ nhất. Đây là bài toán ``tối ưu'': phải vừa \emph{chứng minh} không thể vượt qua một giới hạn, vừa \emph{chỉ ra ví dụ} đạt đúng giới hạn đó.

\begin{recallbox}Trung bình cộng $=$ tổng $:$ số các số, nên tổng $=12\cdot 5=60$. Trung vị của năm số (đã xếp tăng dần) là số thứ ba.\end{recallbox}

\step{1}{Đẳng thức chính}
\[a+b+10+d+e=60\quad\Longrightarrow\quad a+b+d+e=50,\qquad a\le b\le 10\le d\le e.\]
\begin{barfig}
\foreach \x/\lab/\col in {0/a/tmBlue,1.3/b/tmBlue,2.6/10/tmGreen,3.9/d/tmOrange,5.2/e/tmOrange}{\blk{\x}{\x+1.2}{0}{0.6}{\col}{$\lab$}}
\node[anchor=west,font=\footnotesize] at (6.6,0) {tổng $=60$, \ trung vị $=10$};
\end{barfig}

\par\smallskip\textbf{Câu a): làm $e$ lớn nhất.}
\step{2}{Giới hạn trên}
Muốn $e$ lớn thì các số khác phải nhỏ nhất có thể: $a\ge 1$, $b\ge 1$, $d\ge 10$. Do đó $e=50-a-b-d\le 50-1-1-10=38$.
\step{3}{Ví dụ đạt được}
Danh sách $1,1,10,10,38$: tổng $60$, trung vị $10$. Vậy $e$ lớn nhất là $38$.

\par\smallskip\textbf{Câu b): làm $e$ nhỏ nhất.}
\step{4}{Giới hạn dưới}
Muốn $e$ nhỏ thì $a,b$ phải lớn nhất có thể, nhưng $a\le b\le 10$, nên $a+b\le 20$. Suy ra $d+e=50-(a+b)\ge 30$. Vì $d\le e$ nên $2e\ge d+e\ge 30$, tức $e\ge 15$.
\step{5}{Ví dụ đạt được}
Danh sách $10,10,10,15,15$: tổng $60$, trung vị $10$. Vậy $e$ nhỏ nhất là $15$.

\par\smallskip\textbf{Câu c): năm số đôi một khác nhau} ($a<b<10<d<e$).
\step{6}{$e$ lớn nhất}
Nhỏ nhất có thể: $a\ge 1$, $b\ge 2$, $d\ge 11$. Nên $e\le 50-1-2-11=36$. Ví dụ: $1,2,10,11,36$ (tổng $60$). Vậy $e$ lớn nhất là $36$.
\step{7}{$e$ nhỏ nhất}
Lớn nhất có thể: $b\le 9$, $a\le 8$, nên $a+b\le 17$ và $d+e\ge 50-17=33$. Vì $d<e$ nên $d\le e-1$, suy ra $2e-1\ge d+e\ge 33$, tức $e\ge 17$. Ví dụ: $8,9,10,16,17$ (tổng $60$). Vậy $e$ nhỏ nhất là $17$.

\begin{warnbox}Một bài ``lớn nhất/nhỏ nhất'' cần đủ \emph{hai} phần: (1) chứng minh không thể vượt qua giới hạn, và (2) một ví dụ cụ thể đạt đúng giới hạn. Chỉ có ví dụ thì chưa chắc đã là tốt nhất; chỉ có giới hạn thì chưa chắc đạt được.\end{warnbox}

\begin{parentbox}Có thể cho con dùng $60$ viên kẹo chia vào $5$ ``hộp'' xếp theo thứ tự, hộp giữa luôn có $10$ viên. Hỏi: muốn hộp cuối nhiều nhất thì các hộp khác nên có ít hay nhiều? Muốn hộp cuối ít nhất thì sao? Trò chơi này giúp con ``cảm'' được cách lập luận tối ưu.\end{parentbox}

\verify{$1+1+10+10+38=60$; $10+10+10+15+15=60$; $1+2+10+11+36=60$; $8+9+10+16+17=60$; tất cả đều có trung vị $10$. Máy tính duyệt mọi danh sách cũng xác nhận các giá trị $38$, $15$, $36$, $17$.}
\ans{a) $38$. \ b) $15$. \ c) Lớn nhất $36$, nhỏ nhất $17$.}
}

\end{document}
